%% ===========================================================================
%%  e-CALLISTO FITS Analyzer --- User Guide (version 3.1.0)
%%  Sahan S Liyanage, Astronomical and Space Science Unit,
%%  University of Colombo, Sri Lanka
%%
%%  Overleaf: Menu -> Compiler: pdfLaTeX (default). Biber and MakeIndex run
%%  automatically. Main document: main.tex.
%% ===========================================================================
\documentclass[
  paper=a4,
  fontsize=11pt,
  twoside=true,
  open=right,
  BCOR=8mm,
  DIV=11,
  headinclude=false,
  footinclude=false,
  parskip=half-,
  numbers=noenddot,
  toc=bibliography,
  toc=index,
  listof=totoc,
  chapterprefix=false,
  headings=big
]{scrbook}

\usepackage{ecallisto-handbook}

% Bibliography --------------------------------------------------------------
\usepackage[backend=biber, style=authoryear-comp, maxcitenames=2,
  maxbibnames=6, uniquelist=false, uniquename=false, dashed=false,
  sorting=nyt, giveninits=true, doi=true, url=true, isbn=false]{biblatex}
\addbibresource{references.bib}
\DeclareFieldFormat[article,inproceedings]{title}{#1}
\renewcommand*{\bibfont}{\small}

% Index ---------------------------------------------------------------------
\usepackage{imakeidx}
\makeindex[intoc, columns=2, title=Index]
\indexsetup{level=\chapter*, toclevel=chapter, firstpagestyle=plain}

% Faster drafts: compile only selected chapters by uncommenting and editing
% the line below (front matter and back matter are always typeset).
% \includeonly{chapters/16-solar-window,chapters/17-solar-processing}

\begin{document}

% ===========================================================================
\frontmatter
% ===========================================================================
%% Title page: clean typographic design ------------------------------------
\begin{titlepage}
\thispagestyle{empty}
\sffamily\raggedright
\vspace*{-0.6cm}
\noindent{\color{accentdark}\rule{\textwidth}{0.5pt}}\par
\vspace{3.5mm}
\noindent{\small\color{accent}\textls[200]{\MakeUppercase{Software User Guide}}\hfill
  \textls[160]{\MakeUppercase{Version \appversion}}}\par

\vspace*{5.2cm}

\noindent{\fontsize{50}{54}\selectfont\bfseries\color{accentdark}e-CALLISTO\par}
\vspace{2mm}
\noindent{\fontsize{50}{54}\selectfont\color{accent}FITS Analyzer\par}

\vspace{9mm}
\noindent{\color{solar}\rule{2.4cm}{2.4pt}}\par
\vspace{8mm}

{\rmfamily\itshape\fontsize{17}{24}\selectfont\color{black!72}\hyphenpenalty=10000
\noindent Visualization, processing and analysis\\
of solar radio dynamic spectra, solar images\\
and coronal mass ejections\par}

\vfill

\noindent{\Large\bfseries\color{accentdark}Sahan S Liyanage\par}
\vspace{2.5mm}
\noindent{\normalsize\color{black!65}Astronomical and Space Science Unit\par
\vspace{0.6mm}
\noindent University of Colombo, Sri Lanka\par}
\vspace{7mm}
\noindent{\color{accentdark}\rule{\textwidth}{0.5pt}}\par
\vspace*{-0.4cm}
\end{titlepage}

%% Copyright page (verso of the title page) -------------------------------
\thispagestyle{empty}
\null\vfill
{\small\setlength{\parskip}{0.55\baselineskip}
\noindent\textsf{\textbf{e-CALLISTO FITS Analyzer --- User Guide}}\\
For software version \appversion\\
First edition, October 2026

\noindent Copyright \textcopyright{} 2026 Sahan S Liyanage. All rights reserved. No part of
this publication may be reproduced, stored in a retrieval system or transmitted in any
form or by any means without the prior written permission of the author, except for
brief quotations in reviews and scholarly works.

\noindent The \appname{} software is open-source software released under the MIT License.
Its source code, installers and release notes are available from
\url{https://github.com/SaanDev/e-Callisto_FITS_Analyzer} and from the e-CALLISTO software
page, \url{https://www.e-callisto.org/Software/Callisto-Software.html}.

\noindent\textsf{\textbf{How to cite the software.}} If you use the \appname{} in your
research, please cite \textcite{liyanage2026}:\\[2pt]
\hspace*{1em}\begin{minipage}[t]{0.94\linewidth}
Liyanage, G.\,L.\,S.\,S., Adassuriya, J., Jayaratne, K.\,P.\,S.\,C., Monstein, C., and
Manoharan, P.\,K. (2026). e-CALLISTO FITS analyzer: a software framework for CALLISTO
solar radio data. \emph{RAS Techniques and Instruments}, 5, rzag056.
\href{https://doi.org/10.1093/rasti/rzag056}{doi:10.1093/rasti/rzag056}
\end{minipage}

\noindent\textsf{\textbf{Data acknowledgment.}} The radio spectra described in this guide
are recorded by the stations of the e-CALLISTO network. Solar images, X-ray and particle
fluxes, geomagnetic indices and event catalogs are retrieved from the archives listed in
Appendix~\ref{app:data-sources}; users should acknowledge those data providers as their
policies request.

\noindent\textsf{\textbf{Disclaimer.}} Every effort has been made to make this guide accurate
and complete. The software and this guide are provided ``as is'', without warranty of any
kind. Physical quantities derived with the software depend on models and assumptions that
are described in the relevant chapters; results should be validated independently before
they are used for scientific conclusions.

\noindent\textsf{\textbf{Trademarks.}} Windows is a trademark of Microsoft Corporation. macOS
and Apple silicon are trademarks of Apple Inc. Linux is the registered trademark of Linus
Torvalds. Debian is a trademark of Software in the Public Interest, Inc.; Ubuntu is a
trademark of Canonical Ltd. OriginPro is a trademark of OriginLab Corporation. Excel is a
trademark of Microsoft Corporation. Other names are used for identification only and may
be trademarks of their respective owners.

\noindent Typeset in Libertinus and Inconsolata with \LaTeX.\par}
\clearpage

%% Preface -------------------------------------------------------------------
\chapter{Preface}

The \appname{} began as a tool for looking at the solar radio dynamic spectra recorded by
the e-CALLISTO network and for measuring the drift of Type~II bursts. Over successive
releases it has grown into a framework that also retrieves archive data from several
radio instruments, places a radio event in its wider solar and geophysical context,
analyzes solar images from space-borne observatories, and reconstructs coronal mass
ejections in three dimensions. This guide describes all of that functionality as it
exists in version~\appversion.

\section*{About this guide}

This is a \emph{user guide}: a reference to what each window, menu, panel and control does,
what it expects as input, what it produces, and which limits apply. It is not a tutorial
and does not teach solar radio physics or image analysis. Where an operation consists of
several steps, the steps are given as a short numbered procedure; otherwise each feature
is described in the place where you meet it in the interface. Short \emph{Background}
boxes summarize the models behind computed quantities so that results can be interpreted
and reported correctly; the full set of equations is collected in
Appendix~\ref{app:formulas}.

The guide is written for observers, students and researchers who work with e-CALLISTO
data and related space-weather observations. No programming knowledge is required.

\section*{How this guide is organized}

\begin{description}
  \item[Part I, Getting Started] introduces the software, explains installation on
    Windows, macOS and Linux, tours the main window, and describes the help, update,
    citation and bug-report tools.
  \item[Part II, Working with Dynamic Spectra] covers opening and combining FITS files, the
    sidebar processing panels, display, RFI cleaning, annotation and measurement tools,
    the Type~II drift and shock analysis, band-splitting analysis, projects and export.
  \item[Part III, Radio Data Access] describes the e-CALLISTO FITS Downloader, the
    Multi-Station Comparison workspace, and the Learmonth and STEREO/SWAVES loaders.
  \item[Part IV, Solar-Event Context] describes the GOES X-ray, GOES proton, Dst, Kp and
    SOHO/LASCO CME catalog viewers, and the SunPy Multi-Mission Explorer.
  \item[Part V, Solar Image Analysis] covers the multi-mission imaging workspace: data
    retrieval, display and processing, measurement and CME kinematics, and potential-field
    source-surface modeling.
  \item[Part VI, Three-Dimensional CME Reconstruction] describes the GCS CME Fitting window
    for flux-rope and shock reconstruction from several coronagraph viewpoints.
  \item[The appendices] list keyboard shortcuts, file formats and outputs, data sources,
    models and formulas, and troubleshooting advice, followed by a glossary, the
    bibliography and an index.
\end{description}

\section*{Conventions}

Table~\ref{tab:conventions} lists the typographic conventions used throughout the guide.

\begin{table}[!htbp]
  \centering
  \caption{Typographic conventions used in this guide.}
  \label{tab:conventions}
  \small
  \begin{tabularx}{\linewidth}{@{}P{0.34\linewidth}L@{}}
    \toprule
    \thead{Example} & \thead{Meaning} \\
    \midrule
    \menu{File > Export As > Export Figure} & A command reached through a menu and its submenus. \\
    \gui{Subtract Background} & A button, field, check box, tab, panel or window as labeled in the interface. \\
    \kbd{Ctrl+S} & A keyboard shortcut. On macOS, \kbd{Ctrl} in a shortcut corresponds to the Command key \kbd{Cmd}. \\
    \file{project.efaproj} & A file name, file extension or folder. \\
    \callout{3} & A numbered callout in a screenshot, explained in the accompanying table. \\
    \textsf{\textbf{\textcolor{accent}{Note}}}, \textsf{\textbf{\textcolor{tipc}{Tip}}}, \textsf{\textbf{\textcolor{cautionc}{Caution}}} & Boxes with additional detail, a quicker way of doing something, or a warning about lost work or results that need care. \textsf{\textbf{\textcolor{scic}{Background}}} boxes summarize the model behind a computed quantity. \\
    \bottomrule
  \end{tabularx}
\end{table}

The screenshots were taken with the macOS version of the software in the light theme.
On Windows and Linux the window frames, fonts and menu placement follow the conventions
of the operating system, and in the dark theme the colors differ, but the controls and
their behavior are identical on all platforms.

\section*{Acknowledgments}

The \appname{} relies on data from the e-CALLISTO network and its archive, and on the many
public archives listed in Appendix~\ref{app:data-sources}. It is built on open-source
scientific software, including Qt for Python (PySide6), NumPy, SciPy, Matplotlib,
PyQtGraph, Astropy \parencite{astropy2022}, SunPy \parencite{sunpy2020}, aiapy
\parencite{barnes2020}, sunkit-magex \parencite{stansby2020} and reproject. Their authors and
maintainers are gratefully acknowledged.

\begin{flushright}
  \textit{Sahan S Liyanage}\\
  \textit{Colombo, October 2026}
\end{flushright}


\tableofcontents
\listoffigures
\listoftables

% ===========================================================================
\mainmatter
% ===========================================================================
\setpartpreamble{\vspace{2em}\begin{center}\begin{minipage}{0.72\textwidth}\centering\large
What the software does, how to install it, and a tour of the main window,
its menus and the built-in help.\end{minipage}\end{center}}
\part{Getting Started}
%% Chapter 1 -----------------------------------------------------------------
\chapter{Introduction}
\label{ch:introduction}

\section{What the \appname{} is}

The \appname{}\index{e-CALLISTO FITS Analyzer} is a cross-platform desktop application for
visualizing, processing and analyzing solar radio dynamic spectra recorded in FITS format
by the e-CALLISTO spectrometer network, and for placing those observations in their solar
and geophysical context. It runs natively on Windows, macOS and Linux (Debian and Ubuntu)
and is distributed as self-contained installers, so no separate Python installation is
required.

e-CALLISTO\index{e-CALLISTO network} is a world-wide network of CALLISTO radio
spectrometers that monitors the Sun around the clock \parencite{benz2009}. Each station
records dynamic spectra --- radio intensity as a function of frequency and time --- in
FITS files that normally span 15~minutes. A station may operate several receivers or
frequency bands, each identified by a two-digit \emph{focus code}\index{focus code}. File
names follow the pattern
\file{STATION_YYYYMMDD_HHMMSS_FC.fit.gz}, for example
\file{Australia-ASSA_20261004_000000_63.fit.gz} for focus code~63 of the Australia-ASSA
station starting at 00:00:00~UT on 4~October~2026.

The application merges fragmented observations across time and frequency, removes the
receiver background, mitigates radio-frequency interference (RFI)\index{RFI}, and isolates
solar radio bursts with an interactive mask. Dedicated tools derive the frequency drift
rate and shock parameters of Type~II bursts and estimate the coronal magnetic field\index{magnetic field} from
band-split emission. Around this core, the software retrieves archive data from several
radio instruments, displays X-ray, particle, geomagnetic and coronal mass ejection (CME)
data for the same interval, analyzes solar images from space-borne observatories, models
the coronal magnetic field, and reconstructs CMEs in three dimensions.

\section{Capabilities at a glance}

Table~\ref{tab:modules} summarizes the main modules, where they are opened, and where they
are described in this guide.

\begin{table}[!htbp]
  \centering
  \caption{Main modules of the \appname{}.}
  \label{tab:modules}
  \small
  \begin{tabularx}{\linewidth}{@{}P{0.27\linewidth}P{0.33\linewidth}Lr@{}}
    \toprule
    \thead{Module} & \thead{Opened from} & \thead{Purpose} & \thead{Chapter} \\
    \midrule
    Main window & Application start & Display and process dynamic spectra & \ref{ch:main-window}--\ref{ch:processing} \\
    Type~II analysis & Toolbar; \menu{Analysis} & Burst isolation, maxima, drift and shock parameters & \ref{ch:type-ii} \\
    Band-splitting analysis & \menu{Analysis > Type II Band-splitting} & Compression ratio, Alfvén speed\index{Alfvén speed}, magnetic field & \ref{ch:band-splitting} \\
    Projects and export & \menu{File} & Projects, figures, FITS, reports, batch export & \ref{ch:projects} \\
    FITS Downloader & \menu{Download}; toolbar & e-CALLISTO archive search, preview and import & \ref{ch:downloader} \\
    Multi-Station Comparison & \menu{View} & Stacked spectra from several stations & \ref{ch:comparison} \\
    Learmonth, STEREO/SWAVES & \menu{Solar Events > Radio Bursts} & Additional radio spectra & \ref{ch:other-radio} \\
    Solar-event viewers & \menu{Solar Events} & GOES X-rays and protons, Dst, Kp, CME catalog & \ref{ch:solar-events} \\
    SunPy Explorer & \menu{Solar Events > Archives} & Multi-mission archive search and plots & \ref{ch:sunpy} \\
    Solar Image Analysis & \menu{Analysis > Solar Image Analysis} & Imaging, measurement, CME kinematics, PFSS & \ref{ch:solar-window}--\ref{ch:pfss} \\
    GCS CME Fitting & \menu{Analysis > GCS CME Fitting…} & 3-D flux-rope and shock reconstruction & \ref{ch:gcs-window}--\ref{ch:gcs-fitting} \\
    \bottomrule
  \end{tabularx}
\end{table}

\section{Supported data}
\label{sec:supported-data}

The application reads and retrieves the data types listed in Table~\ref{tab:data-types}.
The archives behind each retrieval tool are listed in Appendix~\ref{app:data-sources}.

\begin{table}[!htbp]
  \centering
  \caption{Data types handled by the \appname{}.}
  \label{tab:data-types}
  \small
  \begin{tabularx}{\linewidth}{@{}P{0.36\linewidth}L@{}}
    \toprule
    \thead{Data} & \thead{Use in the application} \\
    \midrule
    CALLISTO FITS (\file{.fit}, \file{.fits}, \file{.fit.gz}, \file{.fits.gz}) & Loaded, combined, processed and analyzed in the main window. \\
    ARTEMIS-IV \texttt{ARTLOOK} FITS & Loaded in the main window with instrument-specific time and intensity scaling (Section~\ref{sec:artemis}). \\
    Learmonth Solar Radio Spectrograph archive & Converted to FIT files and imported (Section~\ref{sec:learmonth}). \\
    STEREO/SWAVES level-2 spectra & Displayed below the CALLISTO spectrum on a shared time axis (Section~\ref{sec:swaves}). \\
    Solar image FITS (SDO/AIA\index{SDO/AIA}, SDO/HMI\index{SDO/HMI}, SOHO/EIT\index{SOHO/EIT}, SOHO/LASCO\index{SOHO/LASCO}, STEREO/SECCHI, GOES/SUVI\index{GOES/SUVI}) & Displayed, processed and measured in Solar Image Analysis (Part~V). \\
    Helioviewer\index{Helioviewer} JPEG2000 images & Near-real-time previews and the coronagraph images used for GCS fitting (Part~VI). \\
    Synoptic magnetograms (GONG, HMI, ADAPT) & Boundary condition for PFSS modeling (Chapter~\ref{ch:pfss}). \\
    GOES XRS and SGPS, Kyoto Dst, GFZ Kp, SOHO/LASCO CME catalog & Context viewers and overlays (Chapter~\ref{ch:solar-events}). \\
    \bottomrule
  \end{tabularx}
\end{table}

\section{Workflow overview}

Figure~\ref{fig:workflow} shows how the main components relate to each other. Data enter
the main window from local files, the downloaders or a saved project; they are processed
in the sidebar; analysis windows take the processed spectrum as input; and every stage can
be saved in a project or exported. The context and imaging tools can be synchronized to
the time window of the loaded spectrum.

\begin{figure}[!htbp]
  \centering
  \begin{tikzpicture}[
      font=\sffamily\footnotesize,
      box/.style={draw=accent, fill=accent!6, rounded corners=2pt, align=center,
                  minimum height=9mm, text width=27mm, inner sep=3pt},
      ctx/.style={box, draw=solar, fill=solar!8},
      arr/.style={-latex, accent, line width=0.7pt},
      node distance=5mm and 7mm]
    \node[box] (src) {Local FITS files\\ Downloaders\\ Projects};
    \node[box, right=of src] (load) {Open and combine\\ (Ch.~\ref{ch:opening})};
    \node[box, right=of load] (proc) {Background, thresholds,\\ RFI cleaning (Ch.~\ref{ch:sidebar}, \ref{ch:processing})};
    \node[box, below=of proc] (iso) {Burst isolation,\\ maximum intensities};
    \node[box, left=of iso] (an) {Burst Analyzer:\\ drift and shock};
    \node[box, left=of an] (bs) {Band-splitting:\\ magnetic field};
    \node[box, below=of an] (out) {Projects, figures,\\ FITS, reports (Ch.~\ref{ch:projects})};
    \node[ctx, below=of bs] (ctx) {Solar-event viewers,\\ SunPy Explorer};
    \node[ctx, below=of iso] (img) {Solar Image Analysis,\\ GCS CME Fitting};
    \draw[arr] (src) -- (load);
    \draw[arr] (load) -- (proc);
    \draw[arr] (proc) -- (iso);
    \draw[arr] (iso) -- (an);
    \draw[arr] (an) -- (bs);
    \draw[arr] (an) -- (out);
    \draw[arr, solar, dashed] (proc.east) -| ++(4mm,-30mm) |- (img.east)
      node[pos=0.25, right, text=solar!80!black, align=left] {time\\window};
    \draw[arr, solar, dashed] (ctx) -- (out);
  \end{tikzpicture}
  \caption{Relationship between the main components of the \appname{}. Solid arrows show
    the radio-analysis chain; dashed arrows show context and imaging tools that use the
    time window of the loaded spectrum.}
  \label{fig:workflow}
\end{figure}

\section{What is new in version \appversion}

Version~\appversion{} adds or substantially extends the following features, all of which
are described in this guide:
\begin{itemize}
  \item automatic ridge tracking in the Maximum Intensities window
    (Section~\ref{sec:ridge-tracking});
  \item a selectable power-law time origin~$t_0$ and five coronal density models with a
    side-by-side comparison in the Burst Analyzer (Section~\ref{sec:analyzer});
  \item drag-and-drop opening and lists of recent files and projects
    (Section~\ref{sec:opening-files});
  \item the \gui{Timeline} panel for extending a spectrum with adjacent observations
    (Section~\ref{sec:timeline}) and a linear or logarithmic frequency axis
    (Section~\ref{sec:axis});
  \item separate background-subtraction controls that always start from the raw data
    (Section~\ref{sec:background});
  \item the \gui{Burst List}\index{FITS Downloader!Burst List} catalog tab and unified time~\(\times\)~frequency combination in
    the downloader (Chapter~\ref{ch:downloader});
  \item PFSS coronal-field modeling, calibration-level selection, differential-rotation
    compensation, circle fitting and higher-order kinematics in Solar Image Analysis
    (Part~V);
  \item multi-view GCS flux-rope and shock fitting with staged refinement, movies and
    fitting reports\index{GCS!fitting report} (Part~VI).
\end{itemize}

\section{Citing the software}
\label{sec:citing}
\index{citation}

If you use the \appname{} in your research, please cite \textcite{liyanage2026}. The
citation and a BibTeX entry can be copied from the application with the
\gui{Cite this Software}\index{Cite this Software} button (Section~\ref{sec:cite-dialog}). The BibTeX entry is:

\begin{tcolorbox}[enhanced, colback=black!3, colframe=rulegray, boxrule=0.5pt, arc=1pt,
  fontupper=\ttfamily\footnotesize, left=6pt, right=6pt]
@article\{liyanage2026ecallistofitsanalyzersoftware,\\
\hspace*{1.5em}title=\{e-CALLISTO FITS analyzer: a software framework for CALLISTO solar radio data\},\\
\hspace*{1.5em}author=\{Liyanage, G L S S and Adassuriya, J and Jayaratne, K P S C and Monstein, C and Manoharan, P K\},\\
\hspace*{1.5em}journal=\{RAS Techniques and Instruments\},\\
\hspace*{1.5em}volume=\{5\}, pages=\{rzag056\}, year=\{2026\},\\
\hspace*{1.5em}doi=\{10.1093/rasti/rzag056\}\\
\}
\end{tcolorbox}

When you publish results that depend on archive data retrieved through the application,
acknowledge the corresponding data providers as well (Appendix~\ref{app:data-sources}).

%% Chapter 2 -----------------------------------------------------------------
\chapter{Installation and Startup}
\label{ch:installation}
\index{installation}

Native installers are available for Windows, macOS and Linux. Each bundles the
application with everything it needs to run, including the Python runtime and all
scientific libraries, so no separate Python installation is required. Installers can be
downloaded from the official e-CALLISTO software page,
\url{https://www.e-callisto.org/Software/Callisto-Software.html}, or from the project's
GitHub releases page, \url{https://github.com/SaanDev/e-Callisto_FITS_Analyzer/releases}.

\section{System requirements}
\label{sec:requirements}

\begin{table}[!htbp]
  \centering
  \caption{System requirements for version \appversion.}
  \label{tab:requirements}
  \small
  \begin{tabularx}{\linewidth}{@{}P{0.24\linewidth}L@{}}
    \toprule
    \thead{Item} & \thead{Requirement} \\
    \midrule
    Windows & 64-bit Windows able to run x64 applications. Installation needs administrator approval. \\
    macOS & macOS~15 (Sequoia) or later on a Mac with Apple silicon (arm64). \\
    Linux & Debian, Ubuntu or a derivative on 64-bit x86 (\texttt{amd64}), with the Qt/OpenGL runtime libraries \texttt{libgl1}, \texttt{libegl1}, \texttt{libxkbcommon-x11-0} and \texttt{libxcb-cursor0} (installed automatically by \texttt{apt}). \\
    Disk space & About 3~GB for the application, plus space for the download caches. Solar image sequences can occupy several gigabytes. \\
    Display & 1280~\(\times\)~800 pixels or larger. A resolution of at least 1440~\(\times\)~900 is recommended for the Solar Image Analysis and GCS CME Fitting windows. \\
    Network & Required for the downloaders, archive searches, context viewers, update checks and bug reports. Local files, saved projects and cached downloads can be analyzed offline. \\
    \bottomrule
  \end{tabularx}
\end{table}

\section{Installing on Windows}
\index{installation!Windows}

\begin{steps}
  \item Download the installer, \file{e-CALLISTO_FITS_Analyzer_v3.1.0_Setup.exe}.
  \item Run the installer and approve the Windows administrator prompt.
  \item Follow the installation steps. The installer can optionally create a desktop
    shortcut.
  \item Start the application from the Start menu or the desktop shortcut.
\end{steps}

\begin{note}
  If a save location chosen by the application is not writable --- for example a folder
  inside \file{C:\\Program Files} --- the application asks you to choose another folder.
\end{note}

\section{Installing on macOS}
\index{installation!macOS}

\begin{steps}
  \item Download the disk image (\file{.dmg}) for Apple silicon (arm64).
  \item Open the disk image and drag the application to the \file{Applications} folder.
  \item Eject the disk image and start the application from \file{Applications} or
    Launchpad.
\end{steps}

\begin{note}
  If macOS reports on first launch that the application cannot be opened because its
  developer cannot be verified, open \menu{System Settings > Privacy \& Security} and
  click \gui{Open Anyway}\index{installation!macOS security prompt} for the \appname{}, or Control-click the application in Finder
  and choose \gui{Open}. This is needed only once.
\end{note}

\section{Installing on Linux}
\index{installation!Linux}

Download the Debian package and install it from a terminal in the download folder:
\begin{center}
  \texttt{sudo apt install ./e-callisto-fits-analyzer\_<version>\_amd64.deb}
\end{center}
The package installs a desktop entry, so the application appears in the desktop
environment's application menu. Using \texttt{apt} rather than \texttt{dpkg} also installs
any missing runtime libraries.

\section{Running from source}
\label{sec:from-source}
\index{installation!from source}

The source code is openly available at
\url{https://github.com/SaanDev/e-Callisto_FITS_Analyzer}. Running from source is intended
for developers and for platforms without an installer. With Python~3.12 or newer:
\begin{steps}
  \item Clone or download the repository and create a virtual environment in it.
  \item Install the pinned dependencies with \texttt{python scripts/install\_requirements.py}.
  \item Start the application with \texttt{python src/main.py}.
\end{steps}
On Windows, if Qt reports \texttt{DLL load failed} at start-up, the script
\file{packaging/windows/repair_windows_venv.ps1} rebuilds the virtual environment with the
supported Python version.

\section{Starting the application}
\label{sec:first-launch}

When the application starts it shows a splash screen and then opens the main window
(Chapter~\ref{ch:main-window}) with an empty plotting area. Shortly after start-up it
checks GitHub for a newer release and reports the result on the right of the status bar,
for example \gui{Updates: up to date} (Section~\ref{sec:updates}).

If the previous session ended unexpectedly while it had unsaved work, the application
offers to recover the latest autosave snapshot (Section~\ref{sec:autosave}).

The application remembers your preferences between sessions, including the theme and
rendering mode, which sidebar sections are expanded, the frequency-axis scale, the recent
files and projects, processing and display-range presets, and the JSOC\index{JSOC} e-mail address used
by Solar Image Analysis.

\section{Updating and uninstalling}

Use \menu{About > Check for Updates...}\index{updates!check} to look for a newer version; if one exists, the
installer can be downloaded from within the application (Section~\ref{sec:updates}).
Installing a new version over an existing one keeps your preferences.

To uninstall, use \menu{Settings > Apps} on Windows, move the application from
\file{Applications} to the Trash on macOS, or run
\texttt{sudo apt remove e-callisto-fits-analyzer} on Linux. Downloaded data remain in the
cache folders; empty them first with the \gui{Clear Cache}\index{cache!clearing} buttons of the downloaders if
you want to reclaim the space.

%% Chapter 3 -----------------------------------------------------------------
\chapter{The Main Window}
\label{ch:main-window}
\index{main window}

The main window is where dynamic spectra are displayed, processed and measured, and from
where every other tool is opened. This chapter describes its layout, toolbar, menus and
status bar. The sidebar panels are described in Chapter~\ref{ch:sidebar}.

\section{Layout}
\label{sec:layout}

The main window (Figure~\ref{fig:main-window}) consists of the regions listed in
Table~\ref{tab:main-regions}.

\screenshot{main-window/main_window}{Main window of the \appname{} before a file is loaded.}{fig:main-window}{Main window with the regions labeled.}

\begin{table}[!htbp]
  \centering
  \caption{Regions of the main window.}
  \label{tab:main-regions}
  \small
  \begin{tabularx}{\linewidth}{@{}P{0.24\linewidth}L@{}}
    \toprule
    \thead{Region} & \thead{Description} \\
    \midrule
    Menu bar & Access to every function, including those not on the toolbar (Section~\ref{sec:menus}). On macOS the menu bar is at the top of the screen. \\
    Toolbar & Icon buttons for the most frequent actions (Section~\ref{sec:toolbar}). \\
    \gui{Cite this Software}\index{Cite this Software} & Opens the citation dialog (Section~\ref{sec:cite-dialog}). \\
    Sidebar & Collapsible panels for the timeline, background subtraction, display thresholds, units, axis scale, graph properties, the analysis summary and the ruler readout (Chapter~\ref{ch:sidebar}). \\
    Main plotting area & The dynamic spectrum with its color bar. Interactive tools act here. \\
    Status bar & Messages on the left; the update status and the live cursor readout on the right (Section~\ref{sec:status-bar}). \\
    \bottomrule
  \end{tabularx}
\end{table}

A narrow arrow handle on the left edge of the plotting area hides or shows the sidebar to
give the plot more room. The divider between the sidebar and the plot can also be dragged.

\section{Toolbar}
\label{sec:toolbar}
\index{toolbar}

The toolbar (Figure~\ref{fig:toolbar}) contains 16 tools, described in
Table~\ref{tab:toolbar}. Hovering over a tool shows its name. Tools that need loaded or
processed data remain disabled until they can be used.

\screenshot{main-window/toolbar}{The main toolbar with its tools numbered.}{fig:toolbar}{Toolbar with numbered callouts 1--16.}

\begin{xltabular}{\linewidth}{@{}c P{0.28\linewidth} L P{0.14\linewidth}@{}}
  \caption{Toolbar tools.}\label{tab:toolbar}\\
  \toprule
  \thead{No.} & \thead{Tool} & \thead{Function} & \thead{Shortcut} \\
  \midrule
  \endfirsthead
  \toprule
  \thead{No.} & \thead{Tool} & \thead{Function} & \thead{Shortcut} \\
  \midrule
  \endhead
  \callout{1} & Open / Load & Load one or more CALLISTO FITS files from disk (Section~\ref{sec:opening-files}). & \kbd{Ctrl+O} \\
  \callout{2} & Download & Open the e-CALLISTO FITS Downloader (Chapter~\ref{ch:downloader}). & \\
  \callout{3} & Export & Export the displayed spectrum as a figure (Section~\ref{sec:export-figure}). & \kbd{Ctrl+E} \\
  \callout{4} & Export as FITS & Export the processed spectrum as a FITS file (Section~\ref{sec:export-fits}). & \kbd{Ctrl+F} \\
  \callout{5} & Save Project & Save the complete workspace as an \file{.efaproj}\index{efaproj@.efaproj file} project (Section~\ref{sec:projects}). & \kbd{Ctrl+S} \\
  \callout{6} & Undo & Undo the last processing change. & \kbd{Ctrl+Z} \\
  \callout{7} & Redo & Redo the last undone\index{undo} change. & \kbd{Ctrl+Shift+Z} \\
  \callout{8} & Estimate Drift Rate & Estimate the drift rate from points clicked along a burst (Section~\ref{sec:drift}). & \\
  \callout{9} & Isolate Burst & Draw a freehand mask around a burst (Section~\ref{sec:isolation}). & \\
  \callout{10} & Plot Maximum Intensities & Open the Maximum Intensities window (Section~\ref{sec:max-intensities}). & \\
  \callout{11} & Rectangular Zooming & Drag one rectangle to zoom into the plot. Available while navigation is locked (Section~\ref{sec:navigation}). & \\
  \callout{12} & Lock zooming and panning & Disable scroll-wheel zooming and panning, freezing the view and enabling rectangle zoom. & \\
  \callout{13} & Unlock zooming and panning & Re-enable scroll-wheel zooming and panning. & \\
  \callout{14} & Reset Selection & Clear the isolated-burst selection and drift markers. & \\
  \callout{15} & Reset to Raw & Discard all processing and show the raw spectrum. & \\
  \callout{16} & Reset All & Clear the workspace: data, processing, selections and analysis overlays. & \\
  \bottomrule
\end{xltabular}

\begin{caution}
  \gui{Reset All}\index{Reset All} clears the whole workspace. Save a project first if you want to keep
  the current state.
\end{caution}

\section{Plotting area and renderers}
\label{sec:renderers}
\index{renderer}\index{hardware acceleration}

The plotting area shows the dynamic spectrum with frequency on the vertical axis, time on
the horizontal axis and intensity as color. The plot title shows the source file name
and the processing state, for example \gui{-Background Subtracted} or
\gui{-Isolated Burst}.

Two renderers are available and are selected with \menu{View > Mode}\index{renderer!Classic and Modern}:
\begin{description}
  \item[Modern]\index{renderer!Modern} uses a hardware-accelerated canvas (PyQtGraph\index{PyQtGraph}) for fast zooming and
    panning of large or combined spectra. This is the default.
  \item[Classic] uses Matplotlib\index{Matplotlib}.
\end{description}
\menu{Processing > Hardware Acceleration > Enable} switches the accelerated live canvas on
or off. Both renderers support the same tools, annotations and overlays, and exported
figures look the same whichever renderer is in use (Section~\ref{sec:export-figure}).

\section{Menu bar}
\label{sec:menus}
\index{menus}

Table~\ref{tab:menus} summarizes the menus of the main window. The tables that follow list
every command, with its shortcut and the section where it is described.

\begin{table}[!htbp]
  \centering
  \caption{Menus of the main window.}
  \label{tab:menus}
  \small
  \begin{tabularx}{\linewidth}{@{}P{0.18\linewidth}L@{}}
    \toprule
    \thead{Menu} & \thead{Main functions} \\
    \midrule
    \menu{File} & Open and combine FITS files; recent files and projects; open, save and recover projects; project reports; figure, FITS, provenance and analysis-log export. \\
    \menu{Edit} & Undo and redo; extend the timeline with adjacent observations; reset to raw data; reset the workspace. \\
    \menu{Download} & Launch the e-CALLISTO FITS Downloader. \\
    \menu{Solar Events} & CME catalog, GOES X-ray and proton flux, geomagnetic indices, SunPy Explorer, radio-burst downloaders, time-window synchronization, GOES overlay and the SWAVES panel. \\
    \menu{View} & FITS header, display range and its presets, view configurations, Multi-Station Comparison, theme and rendering mode. \\
    \menu{Analysis} & Solar Image Analysis, GCS CME Fitting, maximum intensities, Type~II band-splitting, light curves and the ruler. \\
    \menu{Processing} & Hardware acceleration, RFI cleaning, annotations, presets, automatic outlier cleaning, restored analysis sessions and batch processing. \\
    \menu{About} & Update check, bug report and version information. \\
    \menu{Help} & The built-in user guide. \\
    \bottomrule
  \end{tabularx}
\end{table}

\begin{xltabular}{\linewidth}{@{}P{0.36\linewidth} L P{0.17\linewidth}@{}}
  \caption{File menu.}\label{tab:menu-file}\\
  \toprule
  \thead{Command} & \thead{Function} & \thead{Shortcut, see} \\
  \midrule\endfirsthead
  \toprule \thead{Command} & \thead{Function} & \thead{Shortcut, see} \\ \midrule\endhead
  \menu{Open} & Load one file, or combine a compatible multi-file selection. & \kbd{Ctrl+O}, §\ref{sec:opening-files} \\
  \menu{Recent Files} & Reopen one of the last ten FITS selections. & §\ref{sec:recent-files} \\
  \menu{Combine FITS Files...} & Validate, preview and import time, frequency or time~\(\times\)~frequency combinations. & §\ref{sec:combine-dialog} \\
  \menu{Open Project...} & Open a saved \file{.efaproj} project. & \kbd{Ctrl+Shift+O}, §\ref{sec:projects} \\
  \menu{Recent Projects} & Reopen one of the last ten projects. & §\ref{sec:projects} \\
  \menu{Save Project} & Save the workspace to the current project file. & \kbd{Ctrl+S} \\
  \menu{Save Project As...} & Save the workspace to a new project file. & \kbd{Ctrl+Shift+S} \\
  \menu{Recover Last Session...}\index{recovery!Recover Last Session} & Restore the most recent autosave snapshot. & §\ref{sec:autosave} \\
  \menu{Generate Project Report...}\index{project report} & Create a PDF summary of the dataset and analysis. & §\ref{sec:project-report} \\
  \menu{Export As > Export Figure}\index{export!figure} & Save the plot as an image or vector figure. & \kbd{Ctrl+E}, §\ref{sec:export-figure} \\
  \menu{Export As > Export to FIT}\index{export!FITS} & Save the processed spectrum as a FITS file. & \kbd{Ctrl+F}, §\ref{sec:export-fits} \\
  \menu{Export As > Export Provenance Report...}\index{provenance report} & Write Markdown and JSON provenance summaries. & §\ref{sec:provenance} \\
  \menu{Export As > Export Analysis Log...}\index{analysis log} & Write CSV and text summaries of the fit and shock results. & §\ref{sec:provenance} \\
  \bottomrule
\end{xltabular}

\begin{note}
  The disabled \menu{File > Save As} entry in the main window is reserved; projects are
  saved with \menu{Save Project} and \menu{Save Project As...}.
\end{note}

\begin{xltabular}{\linewidth}{@{}P{0.36\linewidth} L P{0.17\linewidth}@{}}
  \caption{Edit menu.}\label{tab:menu-edit}\\
  \toprule \thead{Command} & \thead{Function} & \thead{Shortcut, see} \\ \midrule\endfirsthead
  \toprule \thead{Command} & \thead{Function} & \thead{Shortcut, see} \\ \midrule\endhead
  \menu{Undo} & Undo the last processing or timeline change. & \kbd{Ctrl+Z} \\
  \menu{Redo} & Redo the last undone change. & \kbd{Ctrl+Shift+Z} \\
  \menu{Add Previous Observation} & Prepend the preceding 15-minute observation. & \kbd{Ctrl+Shift+Left}, §\ref{sec:timeline} \\
  \menu{Add Next Observation} & Append the following observation. & \kbd{Ctrl+Shift+Right}, §\ref{sec:timeline} \\
  \menu{Prefetch Adjacent Observations}\index{timeline!prefetch} & When ticked (default), download neighboring observations in the background so the next timeline step is instant. & §\ref{sec:timeline} \\
  \menu{Reset to Raw} & Remove all processing and return to the loaded data. & \\
  \menu{Reset All} & Clear processing, selections, analysis overlays and the loaded data. & \\
  \bottomrule
\end{xltabular}

\begin{xltabular}{\linewidth}{@{}P{0.36\linewidth} L P{0.17\linewidth}@{}}
  \caption{Download and Solar Events menus.}\label{tab:menu-solar-events}\\
  \toprule \thead{Command} & \thead{Function} & \thead{See} \\ \midrule\endfirsthead
  \toprule \thead{Command} & \thead{Function} & \thead{See} \\ \midrule\endhead
  \menu{Download > Launch FITS Downloader} & Open the e-CALLISTO FITS Downloader. & Ch.~\ref{ch:downloader} \\
  \menu{Solar Events > CMEs > SOHO/LASCO CME Catalog}\index{SOHO/LASCO} & Daily CME lists, parameters and movies. & §\ref{sec:lasco-catalog} \\
  \menu{Solar Events > Flares > GOES X-Ray Flux} & GOES XRS viewer and flare parameters. & §\ref{sec:goes-xrs} \\
  \menu{Solar Events > Energetic Particles > GOES SEP Proton Flux} & GOES SGPS proton flux viewer. & §\ref{sec:goes-sep} \\
  \menu{Solar Events > Geomagnetic > Kyoto Dst Index} & Dst index viewer. & §\ref{sec:dst} \\
  \menu{Solar Events > Geomagnetic > GFZ Kp Index} & Kp index viewer. & §\ref{sec:kp} \\
  \menu{Solar Events > Archives > SunPy Multi-Mission Explorer} & Multi-mission archive search and plotting. & Ch.~\ref{ch:sunpy} \\
  \menu{Solar Events > Radio Bursts > e-CALLISTO} & The e-CALLISTO FITS Downloader. & Ch.~\ref{ch:downloader} \\
  \menu{Solar Events > Radio Bursts > Learmonth} & Learmonth spectrograph browser. & §\ref{sec:learmonth} \\
  \menu{Solar Events > Radio Bursts > SWAVES} & STEREO/SWAVES loader. & §\ref{sec:swaves} \\
  \menu{Solar Events > Sync Current Time Window}\index{time synchronization} & Send the time range of the loaded spectrum to open context windows. & §\ref{sec:sync-time} \\
  \menu{Solar Events > GOES Overlay > Short(XRS-A)}\index{GOES!overlay}, \menu{Long(XRS-B)} & Draw GOES X-ray curves on the spectrum. & §\ref{sec:goes-overlay} \\
  \menu{Solar Events > SWAVES Panel}\index{STEREO/SWAVES!panel} & Show or hide the loaded SWAVES panel. & §\ref{sec:swaves} \\
  \bottomrule
\end{xltabular}

\begin{xltabular}{\linewidth}{@{}P{0.36\linewidth} L P{0.17\linewidth}@{}}
  \caption{View menu.}\label{tab:menu-view}\\
  \toprule \thead{Command} & \thead{Function} & \thead{See} \\ \midrule\endfirsthead
  \toprule \thead{Command} & \thead{Function} & \thead{See} \\ \midrule\endhead
  \menu{View FITS Header} & Show the primary header of the loaded file. & §\ref{sec:fits-header} \\
  \menu{Set Display Range...}\index{display range!dialog} & Enter exact time and frequency limits. & §\ref{sec:display-range} \\
  \menu{Save / Apply / Delete Display Range Preset...}\index{display range!presets} & Store and reuse named display ranges. & §\ref{sec:display-range} \\
  \menu{Export View Config...}, \menu{Import View Config...} & Share display settings as \file{.efaview.json}\index{efaview.json@.efaview.json file}. & §\ref{sec:view-config} \\
  \menu{Multi-Station Comparison...} & Open the comparison workspace. & Ch.~\ref{ch:comparison} \\
  \menu{Theme > System / Light / Dark} & Choose the color theme. & §\ref{sec:themes} \\
  \menu{Mode > Classic / Modern} & Choose the Matplotlib or hardware-accelerated renderer. & §\ref{sec:renderers} \\
  \bottomrule
\end{xltabular}

\begin{xltabular}{\linewidth}{@{}P{0.36\linewidth} L P{0.17\linewidth}@{}}
  \caption{Analysis menu.}\label{tab:menu-analysis}\\
  \toprule \thead{Command} & \thead{Function} & \thead{See} \\ \midrule\endfirsthead
  \toprule \thead{Command} & \thead{Function} & \thead{See} \\ \midrule\endhead
  \menu{Solar Image Analysis} & Open the multi-mission imaging workspace. & Ch.~\ref{ch:solar-window} \\
  \menu{GCS CME Fitting…} & Open the GCS window for the visible time range. & Ch.~\ref{ch:gcs-window} \\
  \menu{Maximum Intensities > Open Maximum Intensities} & Extract the peak frequency of each time channel. & §\ref{sec:max-intensities} \\
  \menu{Type II Band-splitting > Open Type II Band-splitting} & Open the band-splitting analyzer. & Ch.~\ref{ch:band-splitting} \\
  \menu{Plot Light Curves > Enter frequency} & Plot a light curve at a typed frequency. & §\ref{sec:light-curves} \\
  \menu{Plot Light Curves > Click on a frequency}\index{light curve!click mode} & Toggle click mode for light curves. & §\ref{sec:light-curves} \\
  \menu{Plot Light Curves > Settings...} & Light-curve appearance. & §\ref{sec:light-curves} \\
  \menu{Plot Light Curves > Clear light curve(s)} & Remove all light curves. & §\ref{sec:light-curves} \\
  \menu{Ruler Measurement} & Measure duration, frequency change and slope between two points. & §\ref{sec:ruler} \\
  \menu{Clear Ruler Measurement} & Remove the ruler. & §\ref{sec:ruler} \\
  \bottomrule
\end{xltabular}

\begin{xltabular}{\linewidth}{@{}P{0.36\linewidth} L P{0.17\linewidth}@{}}
  \caption{Processing, About and Help menus.}\label{tab:menu-processing}\\
  \toprule \thead{Command} & \thead{Function} & \thead{Shortcut, see} \\ \midrule\endfirsthead
  \toprule \thead{Command} & \thead{Function} & \thead{Shortcut, see} \\ \midrule\endhead
  \menu{Processing > Hardware Acceleration > Enable} & Toggle the accelerated live canvas. & §\ref{sec:renderers} \\
  \menu{Processing > RFI Cleaning > Open RFI Panel / Apply RFI / Reset RFI} & RFI cleaning pipeline. & §\ref{sec:rfi} \\
  \menu{Processing > Annotations > …} & Polygon, line and text annotations. & §\ref{sec:annotations} \\
  \menu{Processing > Presets > …} & Processing presets and the default preset. & §\ref{sec:presets} \\
  \menu{Processing > Maximum Intensity > Auto-Clean Isolated Burst Outliers}\index{outlier removal!automatic} & Remove outliers automatically when maxima are taken from an isolated burst. & §\ref{sec:outliers} \\
  \menu{Processing > Analysis Session > Open Restored Analysis}\index{restored analysis session} & Reopen the analysis windows of a restored session. & §\ref{sec:restored-analysis} \\
  \menu{Processing > Batch Processing > Open Batch Processor} & Export many files with consistent settings. & §\ref{sec:batch} \\
  \menu{About > Check for Updates...}\index{updates!check} & Look for a newer release. & §\ref{sec:updates} \\
  \menu{About > Report a Bug...} & Prepare a bug report with diagnostics. & §\ref{sec:bug-report} \\
  \menu{About > About} & Version and author information. & §\ref{sec:about} \\
  \menu{Help > User Guide} & Open the built-in user guide. & \kbd{F1}, §\ref{sec:in-app-guide} \\
  \bottomrule
\end{xltabular}

\section{Status bar}
\label{sec:status-bar}
\index{status bar}\index{cursor readout}

The left part of the status bar shows messages from the current operation, for example
the result of a drift-rate estimate or a reminder of how to finish a tool. The right part
shows:
\begin{itemize}
  \item the update status, for example \gui{Updates: up to date}; and
  \item the live cursor readout \gui{t = … | f = … MHz | I = …}, giving the time,
    frequency and intensity (in the selected units) under the mouse pointer.
\end{itemize}
The readout allows quick quantitative inspection of the spectrum without clicking.

\section{Themes}
\label{sec:themes}
\index{theme}

\menu{View > Theme}\index{theme} selects \gui{System} (follow the operating system), \gui{Light} or
\gui{Dark}. The theme applies to all windows of the application and to on-screen plots.
Exported figures and PDF reports are always drawn on a white page, whichever theme is in
use.

\section{Drag and drop}
\index{drag and drop}

FITS files and \file{.efaproj} projects can be dropped anywhere on the main window,
including the plot. A single FITS file opens on its own; several are combined exactly as
with \menu{File > Open}; a project file opens the project. Other file types are refused.
If the current work has unsaved changes, you are asked whether to save them first.

%% Chapter 4 -----------------------------------------------------------------
\chapter{Help, Updates and Support}
\label{ch:help}

This chapter describes the built-in user guide, the citation dialog, the update check, the
bug-report tool and the version information.

\section{The built-in user guide}
\label{sec:in-app-guide}
\index{user guide!built-in}

\menu{Help > User Guide} (\kbd{F1}) opens a condensed version of this guide in a separate
window (Figure~\ref{fig:in-app-guide}). The guide follows the application theme and offers:
\begin{itemize}
  \item a linked table of contents and cross-references between sections;
  \item \gui{Home} and \gui{Back} buttons for navigation;
  \item a \gui{Find in guide...} field with \gui{Find Next}; and
  \item \gui{Close}.
\end{itemize}
The GCS CME Fitting window has its own help, opened with \kbd{F1} in that window
(Section~\ref{sec:gcs-menus}).

\screenshot[0.7\linewidth]{dialogs/user_guide_dialog}{The built-in user guide window.}{fig:in-app-guide}{The \emph{e-CALLISTO FITS Analyzer User Guide} window opened with F1, showing the table of contents, the Home/Back buttons and the Find field.}

\section{Citing the software}
\label{sec:cite-dialog}
\index{citation!dialog}

The \gui{Cite this Software}\index{Cite this Software} button at the right end of the toolbar opens the citation
dialog (Figure~\ref{fig:cite-dialog}). It shows the recommended citation and a BibTeX
entry, with \gui{Copy Citation} and \gui{Copy BibTeX}\index{citation!BibTeX} buttons that place the text on the
clipboard. The citation is reproduced in Section~\ref{sec:citing}.

\screenshot[0.62\linewidth]{dialogs/citation_dialog}{The \gui{Cite this Software} dialog.}{fig:cite-dialog}{The Cite this Software dialog with the Citation and BibTeX boxes and the Copy Citation, Copy BibTeX and Close buttons.}

\section{Checking for updates}
\label{sec:updates}
\index{updates}

The application checks GitHub for a newer release shortly after it starts and shows the
result in the status bar (\gui{Updates: checking...}, \gui{Updates: up to date}, or a
notice that a new version is available). \menu{About > Check for Updates...}\index{updates!check} repeats the
check on demand:
\begin{itemize}
  \item if a newer version exists, the dialog shows the current and latest version numbers
    with a summary of the release notes, and offers to download the installer for your
    platform to a folder you choose;
  \item if you are up to date, the dialog confirms the installed version;
  \item if the check fails, for example without a network connection, an error message
    explains why.
\end{itemize}
Install the downloaded package as described in Chapter~\ref{ch:installation}.

\section{Reporting a bug}
\label{sec:bug-report}
\index{bug report}\index{diagnostics bundle}

\menu{About > Report a Bug...} opens the bug-report dialog
(Figure~\ref{fig:bug-report}). It contains:
\begin{description}[style=nextline]
  \item[Bug Title] a short summary, for example \emph{Crash when exporting PNG}.
  \item[What happened / Steps to reproduce / Expected behavior] a free-text description,
    pre-filled with a template.
  \item[Generate Diagnostics...] creates a ZIP bundle with a summary of the current session
    (loaded file, plot type, units, color map, combination state and project path),
    environment information (application version, operating system and Python version),
    the operation log, the processing and RFI settings, and a structured provenance
    summary of the loaded data. The path of the bundle is
    shown below the text.
  \item[Copy Issue Text] copies a formatted issue description, including the environment
    summary, to the clipboard.
  \item[Open GitHub Issue] opens a pre-filled issue draft on the project's GitHub page in
    your web browser, where you can attach the diagnostics bundle before submitting.
\end{description}

\screenshot[0.62\linewidth]{dialogs/bug_report_dialog}{The \gui{Report a Bug} dialog.}{fig:bug-report}{The Report a Bug dialog with a title and description filled in, after Generate Diagnostics has created a bundle (its path shown below the text box).}

\begin{note}
  Nothing is sent automatically. The issue is created only when you submit it on GitHub.
  Review the diagnostics bundle before attaching it if your file names or paths are
  confidential.
\end{note}

\section{About}
\label{sec:about}

\menu{About > About} shows the version number, the developer and the affiliation. The
Solar Image Analysis and Maximum Intensities windows have their own \menu{About} menus with
information about those tools.


\setpartpreamble{\vspace{2em}\begin{center}\begin{minipage}{0.72\textwidth}\centering\large
Loading, combining, processing and measuring CALLISTO dynamic spectra, the
Type~II analysis tools, and saving and exporting your work.\end{minipage}\end{center}}
\part{Working with Dynamic Spectra}
%% Chapter 5 -----------------------------------------------------------------
\chapter{Opening and Combining Data}
\label{ch:opening}

This chapter describes the file types the main window reads, the ways of opening them,
and the rules and options for combining several files into one dynamic spectrum.

\section{Supported files}
\index{FITS file!supported formats}

The main window reads CALLISTO dynamic spectra in FITS format, compressed or
uncompressed:
\begin{itemize}
  \item uncompressed files: \file{.fit}, \file{.fits};
  \item gzip-compressed files: \file{.fit.gz}, \file{.fits.gz}.
\end{itemize}
Upper-case suffixes such as \file{.FITS} are accepted. Combination relies on the standard
CALLISTO file name, \file{STATION_YYYYMMDD_HHMMSS_FC}, from which the station, start time
and focus code are read. Files converted from the Learmonth spectrograph
(Section~\ref{sec:learmonth}) follow the same conventions.

\subsection{ARTEMIS-IV files}
\label{sec:artemis}
\index{ARTEMIS-IV}

FITS files from the ARTEMIS-IV radio spectrograph (Thermopylae, Greece) in its
\texttt{ARTLOOK} format are recognized automatically. On loading:
\begin{itemize}
  \item the hours-based UT time axis is converted to the internal time axis;
  \item intensities are labeled and scaled as ADC counts, with the instrument's
    conversion of 58.51~counts per dB used for the dB scale; and
  \item the observation opens background subtracted, because the receiver's channel gains
    span more than an order of magnitude.
\end{itemize}

\section{Opening files}
\label{sec:opening-files}
\index{opening files}

Use any of the following:
\begin{itemize}
  \item \menu{File > Open}, the \gui{Open / Load}\index{toolbar!Open / Load} toolbar tool, or \kbd{Ctrl+O};
  \item drag the files onto the main window;
  \item \menu{File > Recent Files}\index{recent files} (Section~\ref{sec:recent-files});
  \item the \gui{Import} button of a downloader (Chapters~\ref{ch:downloader}
    and~\ref{ch:other-radio}).
\end{itemize}

When one file is selected, it is displayed at once. When several files are selected, the
application determines whether they can be combined in time, in frequency, or in both
(Section~\ref{sec:combine-rules}) and, if so, combines them into a single spectrum. If the
selection cannot be combined, a message explains why
(Figure~\ref{fig:combine-error}). If the current work has unsaved project changes, you are
asked whether to save them first.

After loading, the analysis controls and the \gui{Graph Properties} panel become active,
and the sidebar's \gui{Timeline} panel lists the observations that make up the dataset.

\subsection{Recent files}
\label{sec:recent-files}
\index{recent files}

\menu{File > Recent Files} lists the last ten FITS selections. A combined set is a single
entry, combined again when reopened; hovering over an entry shows its files. An entry
whose files have been moved or deleted is removed from the list with a message.
\gui{Clear Recent Files} empties the list.

\section{Combining files}
\label{sec:combine-rules}
\index{combining files}

CALLISTO observations are fragmented in time (15-minute files) and often in frequency
(several receivers or focus codes). The application can merge them along either axis or
both. The rules are summarized in Table~\ref{tab:combine-rules}.

\begin{table}[!htbp]
  \centering
  \caption{Requirements for combining FITS files.}
  \label{tab:combine-rules}
  \small
  \begin{tabularx}{\linewidth}{@{}P{0.2\linewidth}L@{}}
    \toprule
    \thead{Combination} & \thead{Requirements} \\
    \midrule
    Time & Same station and focus code; consecutive observations whose start times are 750--1050~s apart (normally 15~minutes). Sequences may continue across UTC midnight. \\
    Frequency & Same station, date and observation start time; distinct focus codes; identical time axes. Bands may be separated by a gap or may overlap (Section~\ref{sec:gap-overlap}). \\
    Time~\(\times\)~frequency & A complete grid of at least two consecutive timestamps and at least two focus codes, with the same set of focus codes at every timestamp. \\
    \bottomrule
  \end{tabularx}
\end{table}

A selection that contains extra files is searched for a combinable subset; for example,
selecting a whole day of files together with one stray file still finds the grid within
it. All files of one combined dataset must come from the same station.

For a time~\(\times\)~frequency grid, the bands are first combined at each timestamp with
one shared gap and overlap policy; the application checks that every result has the same
frequency grid, and then joins the results in time order. Incomplete grids, duplicate
timestamp and focus-code pairs, non-consecutive timestamps and incompatible axes are
rejected with a detailed message.

\screenshot[0.62\linewidth]{downloader/time_combine_err}{Error message for files that cannot be combined.}{fig:combine-error}{Warning that names the offending timestamps.}

Combined datasets keep their provenance: exported FITS files, projects and reports record
the original list of source files and the combination method (\texttt{time},
\texttt{frequency} or \texttt{time\_frequency}).

\subsection{Frequency gaps and overlaps}
\label{sec:gap-overlap}
\index{frequency gap}\index{frequency overlap}

When the selected bands leave a gap between them, or overlap, the \gui{Frequency Combine
Options} dialog asks how to handle it before the combined spectrum is imported
(Figures~\ref{fig:gap-options} and~\ref{fig:overlap-options}). The detected gap or overlap
range is shown at the top of the dialog. The chosen option is stored in the metadata of the
combined dataset.

\screenshot[0.62\linewidth]{downloader/freq_combine_gap_options}{Gap handling in the \gui{Frequency Combine Options} dialog.}{fig:gap-options}{Frequency Combine Options dialog for a gap.}

\begin{description}
  \item[Interpolated background fill] fills the gap with background values interpolated
    between the edges of the two bands.
  \item[Gray-hatched blank gap] leaves the gap without data and displays it as a
    gray-hatched region, so that it is clearly distinguished from observed data.
  \item[Average edge background fill] fills the gap with the average background level of
    the band edges on either side.
  \item[Zero fill] fills the gap with zeros.
\end{description}

\screenshot[0.62\linewidth]{downloader/freq_combine_overlap_options}{Overlap handling with a connection frequency.}{fig:overlap-options}{Frequency Combine Options dialog for an overlap.}

\begin{description}
  \item[Split at connection frequency] uses the lower band below the
    \gui{Connection frequency}\index{frequency overlap!connection frequency} and the higher band above it. The connection frequency is
    proposed automatically and can be edited.
  \item[Use lower-frequency file] keeps the data of the lower-frequency file in the
    overlapping region. The file name is shown in the option when there is a single overlap.
  \item[Use higher-frequency file] keeps the data of the higher-frequency file in the
    overlapping region.
  \item[Reject overlap] cancels the combination.
\end{description}

\section{The Combine FITS Files dialog}
\label{sec:combine-dialog}

\menu{File > Combine FITS Files...} opens a dialog for selecting, validating and previewing
a combination before it is imported (Figure~\ref{fig:combine-dialog}).

\begin{steps}
  \item Click \gui{Select FITS Files} and choose the files. The dialog inspects them and
    reports whether they form a valid time, frequency or time~\(\times\)~frequency
    combination.
  \item For a frequency combination, choose the gap and overlap handling and, if needed,
    the connection frequency in the \gui{Frequency Combine Options} group. In this dialog
    the overlap choices are labeled \gui{Keep low band in overlap} and \gui{Keep high
    band in overlap}.
  \item Click \gui{Combine and Preview}\index{Combine FITS Files dialog} to see the combined spectrum in the dialog.
  \item Click \gui{Import to Analyzer} to load the result into the main window.
\end{steps}

\screenshot[0.6\linewidth]{dialogs/combine_fits_dialog}{The \gui{Combine FITS Files} dialog.}{fig:combine-dialog}{The Combine FITS Files dialog after Select FITS Files and Combine and Preview, showing the gap/overlap options group and the combined preview image.}

\section{Viewing the FITS header}
\label{sec:fits-header}
\index{FITS header}

\menu{View > View FITS Header} shows the primary header of the loaded file
(Figure~\ref{fig:fits-header}). \gui{Save as .txt} writes the header to a text file.

\screenshot[0.62\linewidth]{dialogs/fits_header_viewer}{The FITS header viewer.}{fig:fits-header}{The FITS header viewer for a loaded CALLISTO file, showing the header cards and the Save as .txt and Close buttons.}

%% Chapter 6 -----------------------------------------------------------------
\chapter{The Sidebar Panels}
\label{ch:sidebar}
\index{sidebar}

The sidebar on the left of the main window groups the controls used most often while
working on a spectrum. It contains eight sections, from top to bottom:
\gui{Timeline}, \gui{Background Subtraction}, \gui{Noise Clipping Thresholds}, \gui{Units},
\gui{Axis}, \gui{Graph Properties}, \gui{Analysis Summary} and \gui{Ruler Measurement}.
Each section is a card with a clickable header that collapses or expands it; the state of
every card is remembered between sessions. The whole sidebar can be hidden with the arrow
handle on the left edge of the plot.

\section{Timeline}
\label{sec:timeline}
\index{timeline}

A loaded dataset knows which station, focus codes and observation times it came from. The
\gui{Timeline} panel (Figure~\ref{fig:timeline}) lists them and extends or trims the
dataset in place, without returning to the downloader.

\screenshot[0.45\linewidth]{main-window/timeline}{The \gui{Timeline} panel.}{fig:timeline}{Timeline panel with two observations.}

\begin{description}
  \item[◀ Prev / Next ▶] Fetch the observation before the first, or after the last, and
    time-combine it into the spectrum (\kbd{Ctrl+Shift+Left} and \kbd{Ctrl+Shift+Right},
    also in the \menu{Edit} menu).
  \item[Step selector (\(\times\)1, \(\times\)2, \(\times\)4, \(\times\)8)] The number of
    observations added per click.
  \item[Observation list] One row per observation, giving its start time and focus code(s).
  \item[Trim Start / Trim End] Remove the earliest or the latest observation. Trimming down
    to a single observation returns to a plain single-file load.
  \item[Status line] The station, the covered time range and the number of observations,
    for example \gui{BIR · 08:00–08:15 · 2 observations}. When a direction is unavailable
    the panel says why, for example the end of the archive day or a timestamp that supplies
    only part of the focus-code set.
\end{description}

Adjacent observations are looked for first among the loaded file's neighbors on disk, and
only then in the e-CALLISTO archive, so a set that has already been downloaded extends
without network access. Day listings are cached during the session. With
\menu{Edit > Prefetch Adjacent Observations}\index{timeline!prefetch} ticked (the default), neighboring files are
downloaded in the background so that the next step is instant; untick it to avoid
unrequested archive traffic.

Extending and trimming preserve your work:
\begin{itemize}
  \item annotations, the ruler measurement and drift picks are kept, and their time
    coordinates shift with the axis when an earlier observation is prepended;
  \item the active background subtraction and RFI cleaning are recalculated over the
    longer array rather than returning the view to the raw data;
  \item the display thresholds are kept;
  \item every extension or trim is a single undo\index{undo} step (\kbd{Ctrl+Z}), restoring the arrays,
    the source list and the annotations together.
\end{itemize}
After an extension or trim the plot is rescaled to show the full spectrum.

\section{Background Subtraction}
\label{sec:background}
\index{background subtraction}

The receiver background varies strongly from channel to channel, producing horizontal
stripes in a raw spectrum. The \gui{Background Subtraction} panel
(Figure~\ref{fig:bg-sub}) removes it.

\screenshot[0.45\linewidth]{main-window/bg_sub}{The \gui{Background Subtraction} panel.}{fig:bg-sub}{Background Subtraction panel.}

Choose a \gui{Method} and click \gui{Subtract Background}\index{background subtraction!Subtract Background}. The background is computed for
each frequency channel over the whole observation time:
\begin{description}
  \item[Mean] subtracts the mean of each frequency channel; the result stays in the raw
    intensity unit (Digits).
  \item[Median] subtracts the median of each frequency channel; the result stays in
    Digits. The median is less affected by bursts and RFI than the mean.
  \item[Median (dB)] subtracts the median of each channel and converts the result to
    decibels above that background, following the CALLISTO Plotutil\index{Plotutil} recipe used by the
    batch processor. This method is recommended for analysis.
\end{description}

The status line below the button names the method in use, for example
\gui{Applied: Median (dB)}\index{background subtraction!Median (dB)}. The subtraction is always computed from the raw data, so
choosing another method and clicking again replaces the previous result rather than
stacking on it. Re-applying the subtraction also clears any RFI cleaning or isolated burst
that was derived from the previous result. While a Median (dB) result is shown, the
intensity unit is locked to dB because the data are already in decibels.
\menu{Edit > Reset to Raw} removes the subtraction.

\section{Noise Clipping Thresholds}
\label{sec:thresholds}
\index{noise clipping thresholds}\index{display thresholds}

The \gui{Noise Clipping Thresholds} panel (Figure~\ref{fig:noise-clipping}) sets the
limits of the color scale.

\screenshot[0.45\linewidth]{main-window/noise_clipping}{The \gui{Noise Clipping Thresholds} panel.}{fig:noise-clipping}{Noise Clipping Thresholds panel.}

\begin{description}
  \item[Lower Threshold / Upper Threshold] High-resolution sliders for the lower and upper
    limits of the color scale (Vmin and Vmax). The display updates live while you drag;
    no Apply button is needed. The value card beside each label shows the limit. When the
    intensity is shown in Digits, a second line gives the corresponding dB value.
  \item[Logarithmic Threshold Scale] Makes the sliders logarithmic for finer control near
    zero.
\end{description}

The thresholds change the contrast only; the data are never clipped or modified. With
both sliders at zero the color scale spans the whole data range. Limits chosen on the raw
data are reset when you subtract the background, because they would saturate the new
scale. \menu{Processing > Presets > Raw FITS Percentile (5-98\%)}\index{presets!Raw FITS Percentile} sets the limits from the
distribution of the data on show (Section~\ref{sec:presets}).

\section{Units}
\label{sec:units}
\index{units}

The \gui{Units} panel (Figure~\ref{fig:units}) selects the units of the intensity and time
axes.

\screenshot[0.45\linewidth]{main-window/units}{The \gui{Units} panel.}{fig:units}{Units panel.}

\begin{description}
  \item[Intensity: Digits / dB] \gui{Digits} shows raw CALLISTO ADC values as stored in
    the file. \gui{dB} converts them with the CALLISTO receiver scale (see the glossary
    entry \emph{Digit}).
  \item[Time: Seconds / UT] \gui{Seconds} measures time from the start of the loaded data;
    \gui{UT} shows Universal Time. UT requires the observation start time
    (\texttt{TIME-OBS}\index{TIME-OBS keyword}) in the FITS header.
\end{description}
Switching units does not modify or lose data; it only changes the axis labels, the
color-bar label and the cursor readout.

\section{Axis}
\label{sec:axis}
\index{frequency axis!logarithmic}

The \gui{Axis} panel (Figure~\ref{fig:axis}) switches the frequency axis between
\gui{Linear}\index{height--time fit!order} (proportional to MHz) and \gui{Log} (proportional to \(\log_{10}\) MHz).

\screenshot[0.45\linewidth]{main-window/axis}{The \gui{Axis} panel.}{fig:axis}{Axis panel.}

A logarithmic axis spreads out the decametric end of the band, where Type~II and
Type~III bursts spend most of their drift, and places a fundamental--harmonic pair at a
constant separation. Ticks are anchored to the decades and labeled in MHz (for example 20,
30, 40, 50, 60, 70, 80 across a 20--80~MHz band), thinning as the span widens. Both
renderers support the logarithmic axis, and annotations, the ruler, light-curve picks and
drift points behave identically in both scales because their coordinates remain in MHz.
The choice is remembered between sessions. The STEREO/SWAVES panel always uses its own
logarithmic axis (Section~\ref{sec:swaves}).

\section{Graph Properties}
\label{sec:graph-properties}
\index{graph properties}\index{color map}

The \gui{Graph Properties} panel (Figure~\ref{fig:graph-properties}) controls the
appearance of the plot. It becomes active once a file is loaded. The settings apply to
the on-screen plot and to exported figures and reports.

\screenshot[0.45\linewidth]{main-window/graph_properties}{The \gui{Graph Properties} panel.}{fig:graph-properties}{Graph Properties panel.}

\begin{table}[!htbp]
  \centering
  \caption{Controls of the \gui{Graph Properties} panel.}
  \label{tab:graph-properties}
  \small
  \begin{tabularx}{\linewidth}{@{}P{0.16\linewidth}P{0.24\linewidth}L@{}}
    \toprule
    \thead{Group} & \thead{Control} & \thead{Function} \\
    \midrule
    Appearance & Colormap & \gui{Custom} (blue--red--yellow), \gui{viridis}, \gui{plasma}, \gui{inferno}, \gui{magma}, \gui{cividis}, \gui{turbo}, \gui{RdYlBu}, \gui{jet}, \gui{cubehelix} or \gui{bone\_r}. \\
               & Font family\index{graph properties!font} & Typeface of all plot text; \gui{Default} or any installed font. \\
    Text & Graph title & A custom title; leave empty for the default title. \\
         & Remove Titles & Hide the plot title. \\
    Font sizes & Tick labels (px) & 6--60~px (default 11). \\
               & Axis labels (px) & 6--60~px (default 12). \\
               & Title (px) & 6--80~px (default 14). \\
    Text style & Title, Axis labels, Tick labels & \gui{Bold} and \gui{Italic} for each text element. \\
    \bottomrule
  \end{tabularx}
\end{table}

\section{Analysis Summary}
\label{sec:analysis-summary}
\index{analysis summary}

The \gui{Analysis Summary} panel (Figure~\ref{fig:analysis-summary}) is a read-only summary
of the current analysis session: the fitted equation, \(R^2\), the fold number\index{fold number} and the
average shock speed and height from the Burst Analyzer (Section~\ref{sec:analyzer}). It
also reports when a saved session has been restored. Before any analysis it reads
\gui{No analysis session loaded.}

\screenshot[0.45\linewidth]{main-window/analysis_summary}{The \gui{Analysis Summary} panel.}{fig:analysis-summary}{Analysis Summary panel.}

\section{Ruler Measurement}
\label{sec:ruler-panel}
\index{ruler}

The \gui{Ruler Measurement} panel (Figure~\ref{fig:ruler-panel}) shows the result of the
ruler tool, started with \menu{Analysis > Ruler Measurement}
(Section~\ref{sec:ruler}): the two points \gui{P1} and \gui{P2} (time and frequency), the
\gui{Duration}, the \gui{Frequency drift} and the \gui{Slope} in MHz/s. \gui{Copy} copies
the readout to the clipboard and \gui{Clear} removes the measurement.

\screenshot[0.45\linewidth]{main-window/ruler_measurements}{The \gui{Ruler Measurement} panel.}{fig:ruler-panel}{Ruler Measurement panel.}

%% Chapter 7 -----------------------------------------------------------------
\chapter{Viewing, Processing and Annotation Tools}
\label{ch:processing}

This chapter describes how to navigate the spectrum, set an exact display range, clean
radio-frequency interference, use processing presets, annotate the plot, overlay light
curves, and make quick drift-rate and ruler measurements.

\section{Navigating the spectrum}
\label{sec:navigation}
\index{navigation}\index{zoom}

Table~\ref{tab:navigation} lists the mouse interactions available in the plotting area.

\begin{table}[!htbp]
  \centering
  \caption{Navigation and picking in the plotting area.}
  \label{tab:navigation}
  \small
  \begin{tabularx}{\linewidth}{@{}P{0.27\linewidth}L@{}}
    \toprule
    \thead{Interaction} & \thead{How} \\
    \midrule
    Zoom & Scroll the mouse wheel over the plot; the zoom is centered on the pointer. \\
    Pan & Hold the left mouse button and drag. \\
    Rectangle zoom & Click \gui{Lock zooming and panning}\index{navigation!lock and unlock} on the toolbar, then \gui{Rectangular Zooming}\index{zoom!rectangle}, and drag a rectangle. Repeat \gui{Rectangular Zooming} for each further zoom. \\
    Freeze or release the view & \gui{Lock zooming and panning} disables wheel zoom and panning; \gui{Unlock zooming and panning} re-enables them. \\
    Cursor readout & Move the pointer over the plot; time, frequency and intensity appear in the status bar. \\
    Estimate drift rate & Click points along the burst; right-click or double-click to finish (Section~\ref{sec:drift}). \\
    Isolate burst & Press, drag a loop around the burst and release (Section~\ref{sec:isolation}). \\
    Ruler & Click two points (Section~\ref{sec:ruler}). \\
    Light curve by click & With click mode on, click a frequency (Section~\ref{sec:light-curves}). \\
    Annotations & Line: click the start and end. Text: click the position. Polygon: click the vertices and right-click to close (Section~\ref{sec:annotations}). \\
    \bottomrule
  \end{tabularx}
\end{table}

\gui{Reset Selection}\index{Reset Selection} clears an isolated-burst selection and drift markers;
\gui{Reset to Raw}\index{Reset to Raw} discards all processing; \gui{Reset All}\index{Reset All} clears the workspace
(Table~\ref{tab:toolbar}).

\section{Setting an exact display range}
\label{sec:display-range}
\index{display range}

\menu{View > Set Display Range...}\index{display range!dialog} sets the visible window numerically, which is useful
for comparing several stations over identical axes (Figure~\ref{fig:display-range}).
\begin{itemize}
  \item Choose \gui{Use seconds from file start} or \gui{Use UT clock time}, then enter the
    start and stop times in the \gui{Seconds} or \gui{UT} group. UT input is unavailable when
    the file has no \texttt{TIME-OBS}\index{TIME-OBS keyword} value.
  \item Enter the lower and upper limits in the \gui{Frequency} group (MHz).
  \item The full time and frequency ranges of the data are shown for reference.
\end{itemize}

\screenshot[0.5\linewidth]{dialogs/display_range_dialog}{The \gui{Set Display Range} dialog.}{fig:display-range}{The Set Display Range dialog with the Seconds/UT choice, the UT time fields and the Frequency fields, and the full-range information lines.}

\subsection{Display-range presets}
\index{display range!presets}

A display range can be stored under a name with \menu{View > Save Display Range
Preset...}, applied later to another file with \menu{Apply Display Range Preset\index{display range!presets}...}, and
removed with \menu{Delete Display Range Preset...}. Presets are stored in your preferences
and are available in every session.

\subsection{View configurations}
\label{sec:view-config}
\index{view configuration}\index{efaview.json@.efaview.json file}

\menu{View > Export View Config...} saves the display range, units, thresholds, color map
and graph styling to a portable \file{.efaview.json}\index{efaview.json@.efaview.json file} file; \menu{View > Import View
Config...} applies such a file. View configurations can be shared with colleagues and are
also accepted by the batch processor (Section~\ref{sec:batch}) and the Multi-Station
Comparison (Chapter~\ref{ch:comparison}), so that many spectra can be rendered with
identical axes and appearance.

\section{RFI cleaning}
\label{sec:rfi}
\index{RFI cleaning}

Radio-frequency interference appears as bright horizontal streaks (persistent
narrow-band emitters) and short spikes. \menu{Processing > RFI Cleaning} offers
\gui{Open RFI Panel}, \gui{Apply RFI} and \gui{Reset RFI}. The \gui{RFI Cleaning} panel
(Figure~\ref{fig:rfi}) applies a deterministic pipeline to the dynamic spectrum:
\begin{steps}
  \item two-dimensional median smoothing with a window of
    \gui{Kernel (freq)}~\(\times\)~\gui{Kernel (time)}\index{RFI cleaning!kernel};
  \item detection of hot channels from a robust \(Z\)-score of each channel's level and
    variability;
  \item repair of the masked channels from their neighbors;
  \item per-channel clipping of high outliers above an upper percentile (the low side is
    preserved).
\end{steps}

\screenshot[0.5\linewidth]{dialogs/rfi_cleaning_dialog}{The \gui{RFI Cleaning} panel.}{fig:rfi}{The RFI Cleaning panel with Enable RFI cleaning, the kernel, Z-threshold and percentile fields, the Masked channels line after Preview, and the Preview/Apply/Reset/Close buttons.}

\begin{table}[!htbp]
  \centering
  \caption{RFI cleaning parameters.}
  \label{tab:rfi}
  \small
  \begin{tabularx}{\linewidth}{@{}P{0.21\linewidth}P{0.15\linewidth}L@{}}
    \toprule
    \thead{Parameter} & \thead{Range (default)} & \thead{Effect} \\
    \midrule
    Enable RFI cleaning & on/off (on) & Includes RFI cleaning in the processing chain. \\
    Kernel (time) & 1--31, odd (3) & Median window along time. Larger values suppress short spikes more strongly but can blur fast burst structure. \\
    Kernel (freq) & 1--31, odd (3) & Median window across frequency channels. Larger values reduce narrow-band striping but can widen spectral features. \\
    Channel Z threshold\index{RFI cleaning!channel Z threshold} & 0.5--20 (6.0) & Cut-off for hot-channel masking. Lower values mask more channels. \\
    Percentile clip\index{RFI cleaning!percentile clip} & 90--99.99 (99.5) & Upper cap applied to each channel. Lower values clip peaks more aggressively. \\
    \bottomrule
  \end{tabularx}
\end{table}

\begin{description}
  \item[Preview] shows the result without committing it. The \gui{Masked channels}\index{RFI cleaning!masked channels} line
    reports the number and indices of the channels flagged as hot.
  \item[Apply] commits the cleaning to the processed data.
  \item[Reset] restores the default parameters in the panel.
  \item[Close] closes the panel.
\end{description}

If channel streaks remain, lower the \gui{Channel Z threshold} gradually (for example from
6.0 to 5.0 and then 4.0). If burst detail looks over-smoothed, reduce the kernel sizes or
raise the threshold. \menu{Edit > Reset to Raw} reverts all applied processing, including
RFI cleaning.

\section{Processing presets}
\label{sec:presets}
\index{presets}

A processing preset stores the display thresholds and their scale, the background
subtraction method, the intensity and time units, the color map, the graph styling, the
RFI settings and the default annotation style. Presets are managed from
\menu{Processing > Presets}:
\begin{description}
  \item[Raw FITS Percentile (5-98\%)] sets the threshold limits to the 5th and 98th
    percentiles of the data currently shown --- a quick starting range for raw files.
  \item[Save Current as Preset...] stores the current settings under a name.
  \item[Apply Preset...] applies a stored preset to the loaded data.
  \item[Set Default Preset...] chooses a preset that is applied automatically to every
    file loaded afterwards.
  \item[Clear Default Preset] stops applying a default preset.
  \item[Delete Preset...] removes a stored preset.
\end{description}

\section{Annotations}
\label{sec:annotations}
\index{annotations}

Annotations mark features on the spectrum. They are drawn in both renderers, saved in
projects, and included in exported figures and reports. \menu{Processing > Annotations}
(Figure~\ref{fig:annotation-menu}) contains:
\begin{description}
  \item[Add Polygon] click the vertices; right-click to close the polygon.
  \item[Add Line] click the start point and then the end point.
  \item[Add Text] click where the label should go, then enter it in the text dialog.
  \item[Edit Text Label...] choose an existing label and change its text and style.
  \item[Move Text Label...] choose a label, then click its new position.
  \item[Toggle Visibility] show or hide all annotations.
  \item[Delete Last] remove the most recent annotation.
  \item[Clear All] remove every annotation.
\end{description}
The text dialog sets the \gui{Text}, \gui{Font}, \gui{Font Size}, \gui{Font Style}
(\gui{Bold}, \gui{Italic}) and \gui{Color} of a label.

\screenshot[0.24\linewidth]{dialogs/annotation_text_dialog}{The \menu{Processing > Annotations} menu.}{fig:annotation-menu}{The text annotation dialog with the Text, Font, Font Size, Font Style and Color rows.}

\section{Light curves}
\label{sec:light-curves}
\index{light curve}

A light curve is the intensity at one frequency as a function of time. Light curves are
drawn on top of the spectrum (Figure~\ref{fig:light-curves}) from
\menu{Analysis > Plot Light Curves}:
\begin{description}
  \item[Enter frequency] opens a dialog for typing a frequency in MHz; the available range
    of the data is shown.
  \item[Click on a frequency] toggles click mode. While it is on, clicking the spectrum
    plots the light curve at the clicked frequency.
  \item[Settings...] opens the \gui{Light Curve Settings}\index{light curve!settings} dialog (Table~\ref{tab:light-curves}).
  \item[Clear light curve(s)] removes all light curves without affecting the data.
\end{description}

\screenshot{dialogs/light_curve_overlay}{Light curves at two frequencies overlaid on a dynamic spectrum, with the \menu{Analysis > Plot Light Curves} submenu open.}{fig:light-curves}{A background-subtracted spectrum with two light curves at different frequencies, frequency labels shown, and the Light Curve Settings dialog open beside it.}

\begin{table}[!htbp]
  \centering
  \caption{Light-curve settings.}
  \label{tab:light-curves}
  \small
  \begin{tabularx}{\linewidth}{@{}P{0.2\linewidth}P{0.22\linewidth}L@{}}
    \toprule
    \thead{Setting} & \thead{Range} & \thead{Effect} \\
    \midrule
    Mode & Single / Multiple & \gui{Single light curve} replaces the previous curve; \gui{Multiple light curves} keeps adding curves. \\
    Color & color picker & Color of the curve. \\
    Thickness & 0.5--12 & Line width. \\
    Vertical Scale & 0.1--5 & Height of the curve relative to the plot. \\
    Opacity & 0.1--1 & Transparency of the curve. \\
    Line Style & solid, dashed, dotted & Line pattern. \\
    Labels & on/off & \gui{Show frequency labels} next to each curve. \\
    \bottomrule
  \end{tabularx}
\end{table}
\gui{Apply} applies the settings, \gui{Reset to Defaults} restores the defaults and
\gui{Close} closes the dialog. Light curves are saved in projects and can be included in
project reports.

\section{Estimating the drift rate}
\label{sec:drift}
\index{drift rate!estimate}

The \gui{Estimate Drift Rate}\index{drift rate!estimate tool} toolbar tool gives a quick estimate of a burst's frequency
drift from points that you click.
\begin{steps}
  \item Click \gui{Estimate Drift Rate}.
  \item Click two or more points along the burst.
  \item Right-click or double-click to finish.
\end{steps}
The points are joined by dashed segments, and the status bar shows the
\gui{Average Drift Rate} in MHz/s (the mean of the segment drift rates), the start and end
frequencies, and the duration. Segments of zero duration are ignored. \gui{Reset Selection}
clears the drift markers. For a fitted drift rate and shock parameters, use the Type~II
analysis tools (Chapter~\ref{ch:type-ii}).

\section{Ruler measurement}
\label{sec:ruler}

\menu{Analysis > Ruler Measurement} measures between two points clicked on the spectrum.
The readout appears in the sidebar's \gui{Ruler Measurement} panel
(Section~\ref{sec:ruler-panel}): the coordinates of both points, the duration, the
frequency change and the slope. \menu{Analysis > Clear Ruler Measurement} or the panel's
\gui{Clear} button removes it. The ruler is kept when the timeline is extended.

\section{Overlays from other instruments}

Two overlays add context directly to the main plot:
\begin{itemize}
  \item the GOES X-ray flux curves, enabled from \menu{Solar Events > GOES Overlay}\index{GOES!overlay}
    (Section~\ref{sec:goes-overlay}); and
  \item the STEREO/SWAVES spectrum below the CALLISTO spectrum
    (Section~\ref{sec:swaves}).
\end{itemize}

%% Chapter 8 -----------------------------------------------------------------
\chapter{Type II Burst Analysis}
\label{ch:type-ii}
\index{Type II burst}

Type~II solar radio bursts are slowly drifting emission lanes produced by shock waves
traveling outward through the corona. The plasma frequency falls with height, so the
drift of the emission to lower frequencies traces the shock's motion. The Type~II tools
fit the burst backbone and convert its frequency and drift into a shock height and speed
using a coronal electron-density model.

\section{The analysis chain}

The analysis uses four tools, normally applied in this order:
\begin{steps}
  \item \textbf{Burst isolation} restricts the data to the emission of interest
    (Section~\ref{sec:isolation}).
  \item \textbf{Maximum intensities} extracts the frequency of peak emission in each time
    channel, or tracks the burst ridge (Sections~\ref{sec:max-intensities}
    and~\ref{sec:ridge-tracking}).
  \item \textbf{Outlier removal} discards points that do not belong to the burst
    (Section~\ref{sec:outliers}).
  \item \textbf{The Burst Analyzer} fits the backbone and derives the drift rate, shock
    speed and shock height (Section~\ref{sec:analyzer}).
\end{steps}
Each tool works on the spectrum as currently displayed: background-subtracted,
RFI-cleaned or isolated. Figure~\ref{fig:type-ii-spectrum} shows a background-subtracted
spectrum containing a Type~II burst.

\screenshot{analysis/loaded_fits}{A background-subtracted dynamic spectrum containing a Type~II burst.}{fig:type-ii-spectrum}{Main window showing a background-subtracted spectrum with a Type II burst.}

\section{Burst isolation}
\label{sec:isolation}
\index{burst isolation}\index{lasso}

The \gui{Isolate Burst}\index{burst isolation!Isolate Burst tool} toolbar tool (tool~9) draws a freehand lasso mask around the
emission region of a burst. Background subtraction must be applied first.
\begin{steps}
  \item Click \gui{Isolate Burst}.
  \item Press the mouse button, drag a loop around the burst, and release.
\end{steps}
Only the data inside the loop are retained; everything outside is set to zero, which is the
background level of the background-subtracted spectrum. This removes unrelated emission and
residual noise from the analysis (Figure~\ref{fig:isolated}). Isolation can be undone\index{undo} with
\kbd{Ctrl+Z}. The plot
title changes to \gui{-Isolated Burst}. The mask follows the drawn path on the displayed
spectrum. \gui{Reset Selection}\index{Reset Selection} (tool~14) clears the selection.

\screenshot{analysis/isolated_burst}{An isolated Type~II burst lane.}{fig:isolated}{Main window after burst isolation.}

\section{The Maximum Intensities window}
\label{sec:max-intensities}
\index{maximum intensities}

The Maximum Intensities window (Figure~\ref{fig:max-intensities}) shows, for each time
channel of the displayed spectrum, the frequency at which the intensity is highest. These
points trace the burst backbone that the Analyzer fits. Open it with the
\gui{Plot Maximum Intensities}\index{maximum intensities!toolbar tool} toolbar tool (tool~10) or
\menu{Analysis > Maximum Intensities > Open Maximum Intensities}. When the maxima are taken
from an isolated burst, time channels with no emission inside the mask are excluded.

\screenshot{analysis/max_intensities}{The Maximum Intensities window. Without isolation, a persistent RFI channel near 89~MHz produces spurious points.}{fig:max-intensities}{Maximum Intensities window.}

\begin{table}[!htbp]
  \centering
  \caption{Controls of the Maximum Intensities window.}
  \label{tab:max-controls}
  \small
  \begin{tabularx}{\linewidth}{@{}P{0.3\linewidth}L@{}}
    \toprule
    \thead{Control} & \thead{Function} \\
    \midrule
    \gui{Select Outliers}\index{outlier removal!manual} & Draw a lasso around points to remove. \\
    \gui{Remove Outliers} & Remove the selected points. \\
    \gui{Track Ridge}\index{ridge tracking!Track Ridge} & Replace the per-column maxima by an automatically tracked ridge (Section~\ref{sec:ridge-tracking}). \\
    \gui{Pick Start}\index{ridge tracking!Pick Start} & Click a point on the burst from which ridge tracking starts. \\
    \gui{Fundamental} / \gui{Harmonic} & Whether the lane is fundamental or harmonic plasma emission (Section~\ref{sec:harmonic}). \\
    \gui{Analyze Burst}\index{Burst Analyzer!opening} & Open the Burst Analyzer with the current points. \\
    \menu{File > Save As} & Save the points as CSV: time channel, time (s) and frequency (MHz). \\
    \menu{File > Export As} & Save the plot as a figure (PNG, PDF, EPS, SVG, TIFF or JPG). \\
    \menu{Edit > Reset All} & Clear the selection, the ridge start point and the stored Analyzer state. \\
    \menu{Edit > Restore Per-Column Maxima}\index{ridge tracking!restore maxima} & Bring back the points replaced by ridge tracking. \\
    \menu{Analyze > Open Analyzer} & Open the Burst Analyzer. \\
    \menu{Analyze > Ridge Tracking Settings...}\index{ridge tracking!settings} & Parameters of the ridge tracker (Table~\ref{tab:ridge-settings}). \\
    \bottomrule
  \end{tabularx}
\end{table}

The status bar at the bottom of the window reports the result of each action, for example
the number of selected points.

\section{Ridge tracking}
\label{sec:ridge-tracking}
\index{ridge tracking}

Taking the brightest channel in every time column picks up RFI, a second burst or noise,
leaving scattered points to remove by hand. \gui{Track Ridge} follows the burst itself
instead:
\begin{itemize}
  \item Tracking starts from the brightest point of the burst, or from a point chosen with
    \gui{Pick Start} (marked with a star), and continues forwards and backwards in time,
    searching only a few channels around the ridge position in the previous column.
  \item A channel that is bright for most of the file is treated as RFI and is not chosen
    as a starting point. An error is reported if a picked start point is not above the
    background.
  \item Short dropouts in the burst are bridged. After a gap, tracking continues only once
    two consecutive columns show signal, so a single noise spike cannot extend the ridge.
    Tracking stops once the signal has stayed below the threshold for longer than the
    allowed gap.
  \item Each peak is refined to sub-channel frequency.
\end{itemize}
The tracked points replace the per-column maxima; outlier removal, the
fundamental/harmonic choice and \gui{Analyze Burst} work on them as before.
\menu{Edit > Restore Per-Column Maxima} brings the original points back.

\begin{table}[!htbp]
  \centering
  \caption{Ridge Tracking Settings.}
  \label{tab:ridge-settings}
  \small
  \begin{tabularx}{\linewidth}{@{}P{0.3\linewidth}LP{0.14\linewidth}@{}}
    \toprule
    \thead{Setting} & \thead{Description} & \thead{Default} \\
    \midrule
    Search window (channels) & Maximum movement of the ridge between two consecutive time columns (1--100). & 3 \\
    Threshold (\(\sigma\)) & Minimum height of a peak above the background, in robust standard deviations (0--50). & 3.0 \\
    Allowed gap (columns) & Number of consecutive columns below the threshold before tracking stops (0--1000). & 8 \\
    Only follow drift to lower frequency & Restricts the ridge to move toward lower frequencies with time, as in Type~II and Type~III bursts. & Off \\
    \bottomrule
  \end{tabularx}
\end{table}

\screenshot[0.5\linewidth]{dialogs/ridge_tracking_settings}{The \gui{Ridge Tracking Settings} dialog.}{fig:ridge-settings}{The Ridge Tracking Settings dialog with its four fields and OK/Cancel.}

\section{Outlier removal}
\label{sec:outliers}
\index{outlier removal}

Outliers are removed manually in the Maximum Intensities window:
\begin{steps}
  \item Click \gui{Select Outliers} and draw a lasso around the unwanted points
    (Figure~\ref{fig:outliers}).
  \item Click \gui{Remove Outliers}.
\end{steps}

\screenshot{analysis/outlier_removal}{Selecting outliers with the lasso in the Maximum Intensities window.}{fig:outliers}{Outlier selection in the Maximum Intensities window.}

When \menu{Processing > Maximum Intensity > Auto-Clean Isolated Burst Outliers}\index{outlier removal!automatic} is ticked
(the default), outliers are also removed automatically whenever the maxima are taken from
an isolated burst. The manual tools remain available in either case.

\section{Fundamental and harmonic emission}
\label{sec:harmonic}
\index{harmonic emission}

Type~II emission occurs at the local plasma frequency (fundamental) and at twice that
frequency (harmonic). Select \gui{Harmonic} when the fitted lane is the harmonic. The
observed frequencies and drift rates are then divided by two before the shock parameters
are calculated, so that heights and speeds always refer to the plasma frequency. Saved
summaries keep both the converted values and the observed reference values.

\section{The Burst Analyzer}
\label{sec:analyzer}
\index{Burst Analyzer}

The Burst Analyzer (Figure~\ref{fig:analyzer}) fits the burst backbone with a power law and
derives the frequency drift rate, shock speed and shock height, with the goodness of fit.
Open it with \gui{Analyze Burst} or \menu{Analyze > Open Analyzer} in the Maximum
Intensities window.

\screenshot{analysis/analyzer}{The Burst Analyzer window with the best fit, shock parameters and density-model comparison.}{fig:analyzer}{Burst Analyzer window.}

\begin{table}[!htbp]
  \centering
  \caption{Controls and displays of the Burst Analyzer.}
  \label{tab:analyzer}
  \small
  \begin{tabularx}{\linewidth}{@{}P{0.29\linewidth}L@{}}
    \toprule
    \thead{Control or display} & \thead{Function} \\
    \midrule
    \gui{Maximum Intensities} & Plot the extracted points. \\
    \gui{Best Fit} & Fit the points with a power law and plot the fit. \\
    \gui{t₀} & Time origin of the power law (Section~\ref{sec:t0}). \\
    \gui{Density model} & Coronal density model for heights and speeds (Section~\ref{sec:density-models}). \\
    \gui{Fold-number}\index{fold number} & Multiplier 1--4 applied to the model density; press \gui{Calculate} after changing it. \\
    \gui{Best Fit Equation} & The fitted power law. \\
    \gui{Fit Metrics} & \(R^2\) and RMSE of the fit. \\
    \gui{Shock Parameters} & Derived quantities with uncertainties (Table~\ref{tab:shock-params}). \\
    \gui{Density model comparison}\index{density model!comparison table} & Initial and average shock speed and height from all five models (Section~\ref{sec:model-comparison}). \\
    \gui{Save Graph}\index{Burst Analyzer!Save Graph} & Save the plot on show as a figure. \\
    \gui{Save Data}\index{Burst Analyzer!Excel export} & Save the results to an Excel workbook (Section~\ref{sec:analyzer-export}). \\
    \gui{Existing Excel File} & Append to an existing workbook instead of creating a new one. \\
    \gui{Extra Plots}\index{Burst Analyzer!extra plots} and \gui{Plot} & Additional plots (Section~\ref{sec:extra-plots}). \\
    \bottomrule
  \end{tabularx}
\end{table}

\subsection{Power-law fit and the time origin}
\label{sec:t0}
\index{power law}\index{time origin}

The backbone is fitted as
\begin{equation}
  f(t) = a\,(t - t_0)^{-b},
  \label{eq:power-law}
\end{equation}
where \(f\) is the frequency in MHz, \(t\) the time in seconds, and \(a\) and \(b\) the fitted
parameters. The drift rate at each point follows from the derivative,
\(\mathrm{d}f/\mathrm{d}t = -ab\,(t-t_0)^{-b-1}\). On the \gui{Best Fit} graph the points are
drawn as filled black squares and the fitted power law as a red curve, with a boxed legend.

Because a power law depends on where \(t=0\) lies, the fitted exponent, drift rate and shock
parameters depend on the chosen time origin. \(t_0\) is selected from the \gui{t₀} list next
to \gui{Best Fit}:
\begin{description}
  \item[File start] measures time from the start of the loaded data (the default).
  \item[Burst onset] places \(t_0\) one time sample before the first fitted point, so that
    every point lies on the curve.
  \item[Custom time] uses a time that you enter, in UT when the start time of the file is
    known and otherwise in seconds from the start of the file.
\end{description}
Changing \(t_0\) re-runs the fit automatically. The graph keeps the data's own time axis,
and the equation shows \((x - t_0)\) whenever \(t_0\) is not the file start.

\subsection{Shock parameters}

Table~\ref{tab:shock-params} defines the quantities in the \gui{Shock Parameters} block. The
heading names the density model and fold in use, for example
\gui{Shock Parameters (Newkirk 1-fold)}\index{density model!Newkirk}.

\begin{table}[!htbp]
  \centering
  \caption{Shock parameters reported by the Burst Analyzer.}
  \label{tab:shock-params}
  \small
  \begin{tabularx}{\linewidth}{@{}P{0.27\linewidth}L@{}}
    \toprule
    \thead{Quantity} & \thead{Definition} \\
    \midrule
    Average Frequency & Mean frequency of the maximum-intensity points (MHz), with its standard error. \\
    Average Drift Rate & Mean of the fitted drift rate along the backbone (MHz/s). \\
    Starting Frequency\index{shock parameters!starting frequency} & The 90th percentile of the fitted frequencies, taken as the start of the burst (MHz). \\
    Initial Shock Speed\index{shock parameters!initial speed} & Shock speed at the starting frequency (km/s). \\
    Initial Shock Height & Heliocentric height at the starting frequency (\(R_\odot\)). \\
    Average Shock Speed\index{shock parameters!average speed} & Mean shock speed along the backbone (km/s), with its standard error. \\
    Average Shock Height & Mean height along the backbone (\(R_\odot\)), with its standard error. \\
    \bottomrule
  \end{tabularx}
\end{table}

\begin{background}[Background: from drift rate to shock speed]
  The emission frequency is identified with the local plasma frequency,
  \(f_p \approx 8.98\times10^{-3}\sqrt{n_e}\)~MHz for an electron density \(n_e\) in
  cm\(^{-3}\). A density model \(n_e(r)\) therefore maps each frequency to a heliocentric
  height \(r\), and the drift rate to a radial speed,
  \begin{equation*}
    v = \left|\frac{\mathrm{d}f}{\mathrm{d}t}\right| \left|\frac{\mathrm{d}r}{\mathrm{d}f}\right|,
    \qquad \frac{\mathrm{d}r}{\mathrm{d}f} = \frac{2\,n_e}{f\,|\mathrm{d}n_e/\mathrm{d}r|}.
  \end{equation*}
  Uncertainties are propagated linearly from the drift-rate and frequency errors. The
  results are only as good as the density model, which describes an average corona and
  can differ from the actual medium by large factors; the fold number scales the model
  density to account for denser structures such as streamers.
\end{background}

\subsection{Density models}
\label{sec:density-models}
\index{density model}

Shock heights and speeds are calculated from one of five coronal electron-density models
(Table~\ref{tab:density-models}). Changing the model recalculates the shock parameters
immediately. The fold number (1--4) multiplies the model density; press \gui{Calculate}
after changing it. The formulas are given in Appendix~\ref{app:formulas}.

\begin{table}[!htbp]
  \centering
  \caption{Coronal electron-density models.}
  \label{tab:density-models}
  \small
  \begin{tabularx}{\linewidth}{@{}P{0.25\linewidth}L@{}}
    \toprule
    \thead{Model} & \thead{Reference} \\
    \midrule
    Newkirk (default) & \textcite{newkirk1961} \\
    Saito\index{density model!Saito} & \textcite{saito1977} \\
    Leblanc\index{density model!Leblanc} & \textcite{leblanc1998} \\
    Baumbach\index{density model!Baumbach--Allen}--Allen & \textcite{baumbach1937}; \textcite{allen1947} \\
    Mann & \textcite{mann1999}\index{density model!Mann} \\
    \bottomrule
  \end{tabularx}
\end{table}

Models other than Newkirk are evaluated between the photosphere (\(1\,R_\odot\)) and 1~AU.
A frequency that a model cannot reach in this range --- for example above about 116~MHz
for the 1-fold Saito model --- has no height, and is shown as a dash rather than an
extrapolated value.

\subsection{Density-model comparison}
\label{sec:model-comparison}

The \gui{Density model comparison} table below the shock parameters lists the initial shock
speed \(v_0\), initial height \(h_0\), average speed \(\bar v\) and average height \(\bar h\)
obtained with every model at the current fold, with the selected model in bold. It shows at
a glance how strongly the derived kinematics depend on the assumed corona.

\subsection{Additional plots}
\label{sec:extra-plots}

Choose an entry in \gui{Extra Plots} and click \gui{Plot} to show:
\gui{Shock Speed vs Shock Height}, \gui{Shock Speed vs Frequency} or
\gui{Shock Height vs Frequency}. These plots use the selected density model.

\subsection{Saving graphs and data}
\label{sec:analyzer-export}
\index{Excel export}

\gui{Save Graph} writes whichever plot is on show --- maximum intensities, best fit or an
additional plot --- as an OriginPro-style graph on a white page in PNG, PDF, EPS, SVG, TIFF
or JPG format.

\gui{Save Data} writes one row of results to an Excel workbook (\file{.xlsx}). With
\gui{Existing Excel File} ticked, you choose a workbook and the row is appended, which builds
an event table over many analyses. The columns are listed in Table~\ref{tab:excel-columns}.

\begin{table}[!htbp]
  \centering
  \caption{Columns of the Burst Analyzer Excel export.}
  \label{tab:excel-columns}
  \small
  \begin{tabularx}{\linewidth}{@{}rL@{}}
    \toprule
    \thead{No.} & \thead{Column} \\
    \midrule
    1--2 & Observation date; station \\
    3--5 & Best-fit equation; \(R^2\); RMSE \\
    6--9 & Average frequency and error; average drift rate and error \\
    10--11 & Starting frequency and error \\
    12--15 & Initial shock speed and error; initial shock height and error \\
    16--19 & Average shock speed and error; average shock height and error \\
    20 & Absolute value of the average drift rate \\
    21 & Density model \\
    22 & \(t_0\) in seconds from the start of the file \\
    \bottomrule
  \end{tabularx}
\end{table}

The density model and \(t_0\) are also saved in projects and named in project reports.
Projects saved before these options existed open with the Newkirk model and the file start
as \(t_0\), matching how they were calculated.

%% Chapter 9 -----------------------------------------------------------------
\chapter{Type II Band-Splitting Analysis}
\label{ch:band-splitting}
\index{band splitting}

The emission lanes of many Type~II bursts are split into two parallel bands, interpreted
as emission from the plasma ahead of (upper band) and behind (lower band) the shock front.
The frequency ratio of the two bands gives the shock's density compression ratio\index{compression ratio}, from
which the Alfvén Mach number\index{Alfvén Mach number}, the Alfvén speed\index{Alfvén speed} and the coronal magnetic field\index{magnetic field} can be
estimated \parencite{vrsnak2002}. The Type~II Band-splitting window performs this analysis
directly on the current noise-reduced dynamic spectrum.

\section{Opening the window}

Open the window with \menu{Analysis > Type II Band-splitting > Open Type II
Band-splitting}. The calculation needs the shock speed and the starting frequency from the
Burst Analyzer, so run the Analyzer (Section~\ref{sec:analyzer}) first. The
\gui{Analyzer Reference}\index{band splitting!Analyzer reference} panel reports \gui{Analyzer reference data loaded.} when these
values are available, and \gui{Analyzer input missing.} otherwise.

\section{Window layout}

The tools and panels of the window are numbered in Figure~\ref{fig:band-splitting} and
described in Table~\ref{tab:band-splitting-tools}.

\screenshot{analysis/typeII_band_splitting_window}{The Type~II Band-splitting window with its tools and panels numbered.}{fig:band-splitting}{Band-splitting window with numbered callouts.}

\begin{table}[!htbp]
  \centering
  \caption{Tools and panels of the Type~II Band-splitting window.}
  \label{tab:band-splitting-tools}
  \small
  \begin{tabularx}{\linewidth}{@{}cP{0.25\linewidth}L@{}}
    \toprule
    \thead{No.} & \thead{Tool or panel} & \thead{Function} \\
    \midrule
    \callout{1} & Add Points & Toggle point picking; click along the active band on the spectrum to add points. \\
    \callout{2} & Undo Last Point & Remove the last point added to the active band. \\
    \callout{3} & Clear Active Band & Remove all points of the active band. \\
    \callout{4} & Fit Active Band & Fit the points of the active band with a power law. \\
    \callout{5} & Fit Both Bands & Fit the upper and the lower band. \\
    \callout{6} & Save Plot & Save the fitted lanes over the spectrum, or the magnetic field versus height plot, as a figure. \\
    \callout{7} & BvR & Toggle the plot of magnetic field versus shock height (Section~\ref{sec:bvr}). \\
    \callout{8} & Settings & Graph settings: fonts, lane colors, line widths and marker sizes (Section~\ref{sec:bs-graph-settings}). \\
    \callout{9} & Controls panel & \gui{Active Band} (\gui{Upper Band} or \gui{Lower Band}), \gui{Shock Speed} (\gui{Initial Shock Speed}\index{shock parameters!initial speed} or \gui{Average Shock Speed}\index{shock parameters!average speed}) and \gui{Calculate}. \\
    \callout{10} & Results panels & \gui{Analyzer Reference}, \gui{Fitted Bands} and \gui{Averaged Band-Splitting Quantities} (Section~\ref{sec:bs-results}). \\
    \bottomrule
  \end{tabularx}
\end{table}

\section{Picking and fitting the bands}
\label{sec:bs-fitting}

\begin{steps}
  \item Select the \gui{Active Band} in the Controls panel \callout{9}.
  \item Switch on \gui{Add Points} \callout{1} and click points along that band on the
    spectrum. Points of the upper and lower bands are drawn in different colors.
  \item Select the other band and add its points.
  \item Fit each band with \gui{Fit Active Band} \callout{4}, or both at once with
    \gui{Fit Both Bands}\index{band splitting!fitting} \callout{5}.
\end{steps}
Each band is fitted with a power law \(f(t) = a\,t^{-b}\), where \(t\) is measured from the
start of the file. At least two points with positive time and frequency are needed per
band. The \gui{Fitted Bands} panel shows each fitted equation, \(R^2\), RMSE and the number
of selected points.

\section{Calculation and results}
\label{sec:bs-results}

Choose whether the \gui{Initial Shock Speed} or the \gui{Average Shock Speed} from the
Analyzer is used, and click \gui{Calculate}. The results panels contain:
\begin{description}
  \item[Analyzer Reference] the reference values imported from the Burst Analyzer: the
    fold number\index{fold number}, the starting frequency \(f_\mathrm{start}\), the average drift rate
    \(\langle\mathrm{d}f/\mathrm{d}t\rangle\), the initial and average shock speeds
    \(V_{s,0}\) and \(\langle V_s\rangle\), and the initial and average shock heights
    \(R_{p,0}\) and \(\langle R_p\rangle\).
  \item[Fitted Bands] the fitted equation, \(R^2\), RMSE and number of points of each band.
  \item[Averaged Band-Splitting Quantities] the averaging interval, the start time
    \(t_\mathrm{start}\) and band frequencies, the mean upper and lower frequencies, the
    bandwidth \(\langle\Delta f\rangle = \langle f_u - f_l\rangle\), the mean upper-band
    drift \(\langle\mathrm{d}f_u/\mathrm{d}t\rangle\), the compression ratio
    \(X = (\langle f_u\rangle/\langle f_l\rangle)^2\), the Alfvén Mach number \(M_A\), the
    Alfvén speed \(V_A\) and the magnetic field \(B\).
\end{description}
The averaging interval is the time span of the upper-band points; the two fitted curves are
sampled over it. If the lower-band fit has to be extrapolated over part of the interval, a
warning is shown. The calculation is refused if the compression ratio does not stay
between 0 and 4 over the interval.

\begin{background}[Background: band-splitting relations]
  For a perpendicular shock in a low-\(\beta\) plasma the compression ratio \(X\) and the
  Alfvén Mach number are related by \parencite{vrsnak2002}
  \begin{equation*}
    X = \left(\frac{f_u}{f_l}\right)^2, \qquad
    M_A = \sqrt{\frac{X\,(X+5)}{2\,(4-X)}}.
  \end{equation*}
  The Alfvén speed follows from the shock speed \(V_s\) taken from the Analyzer,
  \(V_A = V_s/M_A\), and the magnetic field from
  \(B = V_A\sqrt{\mu_0\,m_p\,n_e}\), with the electron density obtained from the plasma
  frequency, \(n_e = (f/8.98\times10^{-3}\,\mathrm{MHz})^2\)~cm\(^{-3}\). In the averaged
  result the density is taken at the Analyzer's starting frequency.
\end{background}

\section{Magnetic field versus height}
\label{sec:bvr}
\index{magnetic field!versus height}

After a successful calculation, the \gui{BvR}\index{band splitting!B vs R plot} tool \callout{7} plots the magnetic field
against shock height (Figure~\ref{fig:bvr}). Along the averaging interval the field is
computed at each sample time from the local compression ratio, with the density taken from
the upper-band frequency. The heights are calculated with the Newkirk\index{density model!Newkirk} model and the
Analyzer's fold number, and the profile is fitted with a power law \(B = a\,R^{-b}\). Click
\gui{BvR} again to return to the spectrum.

\screenshot{analysis/BvR}{Magnetic field versus shock height in the band-splitting window.}{fig:bvr}{Band-splitting window showing the B vs R plot.}

\begin{note}
  The shock speed is taken from the Burst Analyzer and therefore follows the density model
  selected there (Section~\ref{sec:density-models}). The heights of the magnetic field
  versus height profile always use the Newkirk model, with the fold number set in the
  Analyzer.
\end{note}

\section{Graph settings}
\label{sec:bs-graph-settings}

The \gui{Settings} tool \callout{8} opens the \gui{Type II Graph Settings}\index{band splitting!graph settings} dialog
(Figure~\ref{fig:bs-settings}):
\begin{description}
  \item[Text Formatting] font family, and size, bold and italic for the title, the axis
    labels and the tick labels.
  \item[Upper Band, Lower Band] color and width of the fit line, and color and size of
    the picked points.
  \item[B vs R] color and width of the fit line, and color and size of the markers.
\end{description}
\gui{Apply} applies the settings, \gui{Reset to Defaults} restores the defaults and
\gui{Close} closes the dialog.

\screenshot[0.45\linewidth]{dialogs/type_ii_graph_settings}{The \gui{Type II Graph Settings} dialog.}{fig:bs-settings}{The Type II Graph Settings dialog showing the Text Formatting, Upper Band, Lower Band and B vs R groups.}

\section{Saving plots}

\gui{Save Plot} \callout{6} exports either the fitted lanes over the dynamic spectrum or the
magnetic field versus height fit, whichever is on show, as an OriginPro-style graph in PNG,
PDF, EPS, SVG, TIFF or JPG format. The lane colors and fonts from the graph settings are
kept.

\begin{caution}
  Derived magnetic-field values may not be accurate for all events or under all
  assumptions. They depend on the shock geometry, the density model and the
  interpretation of the split bands. Confirm results against known or independently
  validated events before using them for scientific conclusions.
\end{caution}

%% Chapter 10 ----------------------------------------------------------------
\chapter{Projects, Reports and Export}
\label{ch:projects}

This chapter describes how the state of an analysis is saved and recovered, and how
figures, processed data, reports and batches of files are exported.

\section{Projects}
\label{sec:projects}
\index{project}\index{efaproj@.efaproj file}

A project file (\file{.efaproj}\index{efaproj@.efaproj file}) stores the complete state of the analyzer so that work can
be resumed exactly where it was left. A project contains:
\begin{itemize}
  \item the loaded or combined data, with the list of source files and the combination
    method;
  \item the processing state: background subtraction, RFI cleaning, isolated burst;
  \item the view: display range, thresholds, units, axis scale, color map and graph
    properties;
  \item annotations, light curves and their settings, and a loaded STEREO/SWAVES panel
    (restored without downloading it again);
  \item the active processing preset, the operation log and the last time-window
    synchronization;
  \item the analysis session: maximum-intensity points, the fundamental/harmonic choice,
    the Analyzer's fit, density model, fold and \(t_0\), and band-splitting points and fits.
\end{itemize}

Projects are managed from the \menu{File} menu:
\begin{description}
  \item[Save Project] (\kbd{Ctrl+S}, or the toolbar tool) saves to the current project
    file, asking for a name the first time.
  \item[Save Project As...] (\kbd{Ctrl+Shift+S}) saves to a new file.
  \item[Open Project...] (\kbd{Ctrl+Shift+O}) opens a project. A project file can also be
    dropped onto the main window.
  \item[Recent Projects]\index{recent projects} lists the last ten projects opened or saved;
    \gui{Clear Recent Projects} empties the list, and an entry whose file is missing is
    removed with a message.
\end{description}
Before a project is opened, or new data are loaded, you are asked whether to save unsaved
changes.

\subsection{Restored analysis sessions}
\label{sec:restored-analysis}

When a project with an analysis session is opened, the \gui{Analysis Summary} panel reports
that the session has been restored. \menu{Processing > Analysis Session > Open Restored
Analysis} reopens the windows that were open when the project was saved --- the Maximum
Intensities window, the Burst Analyzer and the Type~II Band-splitting window --- with their
points, fits and settings.

\section{Autosave and recovery}
\label{sec:autosave}
\index{autosave}\index{recovery}

While data are loaded and there are unsaved changes, the application writes a recovery
snapshot every three minutes. The ten most recent snapshots are kept. If the application
did not close normally, it offers to recover the latest snapshot at the next start.
\menu{File > Recover Last Session...}\index{recovery!Recover Last Session} restores the most recent snapshot at any time.

\section{Exporting figures}
\label{sec:export-figure}
\index{export!figure}\index{OriginPro style}

\menu{File > Export As > Export Figure}\index{export!figure} (\kbd{Ctrl+E}, or the \gui{Export} toolbar tool)
saves the main plot as a PNG, PDF, EPS, SVG, TIFF or JPG file. Whatever renderer and theme
are on screen, the exported figure is drawn as a publication-style (``OriginPro-style'')
graph in light mode:
\begin{itemize}
  \item Arial type, a closed black frame with inward major and minor ticks on all four
    sides, a framed color scale, and time ticks on whole UT minutes;
  \item the current zoom, color map, display range, linear or logarithmic frequency axis,
    and the titles and font sizes from \gui{Graph Properties};
  \item every visible overlay: light curves, annotations, the GOES X-ray overlay, the ruler
    measurement and the STEREO/SWAVES panel.
\end{itemize}
The plot on screen keeps its own look. On Windows, if the proposed folder is not writable
(for example inside \file{C:\\Program Files}), you are asked to choose another folder.

The Maximum Intensities window, the Burst Analyzer, the band-splitting window, the
Multi-Station Comparison and the Solar Image Analysis tracking panel export their graphs in
the same style.

\section{Exporting FITS files}
\label{sec:export-fits}
\index{export!FITS}

\menu{File > Export As > Export to FIT}\index{export!FITS} (\kbd{Ctrl+F}, or the \gui{Export as FITS} toolbar
tool) writes the spectrum on show as a new FITS file, for downstream analysis or machine
learning. The processed data are exported if processing has been applied, and the raw data
otherwise; combined datasets are exported as one file.

Before saving, choose the data type (\texttt{BITPIX}\index{BITPIX keyword}): \gui{Auto} (the recommended type is
shown), \gui{8 (unsigned byte)}, \gui{16 (signed int16)} or \gui{32 (signed int32)}.
Integer types are offered because some FITS viewers cannot read 64-bit floating-point data.
The exported file keeps the original header, updated with:
\begin{itemize}
  \item \texttt{HISTORY} records naming the application, the plot type, the intensity
    units and the source file(s);
  \item \texttt{COMBINED}, \texttt{COMBMETH} (combination method) and \texttt{NFILES} for
    combined data;
  \item \texttt{BUNIT}, \texttt{DATAMIN} and \texttt{DATAMAX};
\end{itemize}
and an \texttt{AXIS} binary-table extension with the \texttt{FREQUENCY} and \texttt{TIME}
columns, as in standard CALLISTO files (see Appendix~\ref{app:file-formats}).

\section{Provenance reports and analysis logs}
\label{sec:provenance}
\index{provenance report}\index{analysis log}

For reproducibility, two structured reports can be exported from \menu{File > Export As}:
\begin{description}
  \item[Export Provenance Report...] writes a Markdown (\file{.md}) and a JSON
    (\file{.json}) summary with the sections \emph{App}, \emph{Data Source},
    \emph{Processing}, \emph{RFI}, \emph{Annotations}, \emph{Time Sync} and
    \emph{Operation Log}.
  \item[Export Analysis Log...] writes \file{<name>_analysis_log.csv} and
    \file{<name>_analysis_log.txt} with the Analyzer's fit parameters and the derived shock
    quantities.
\end{description}
These reports are suitable for laboratory notebooks, collaboration hand-overs and the
preparation of papers.

\section{Project report PDF}
\label{sec:project-report}
\index{project report}

\menu{File > Generate Project Report...}\index{project report} creates a consolidated PDF of the current project
or loaded dataset. A progress dialog is shown while the report is built. The report
contains the sections listed below; sections without content are omitted.
\begin{itemize}
  \item \emph{Data Source} and \emph{Selected FITS Header};
  \item \emph{Processing}, and \emph{RFI, Annotations, and Light Curves};
  \item \emph{Analysis Summary} and \emph{Type II Results}, including the fit equation (with
    \(t_0\) when it is not the file start), the density model, and band-splitting output;
  \item \emph{Figures}: the raw and background-subtracted spectra, light curves with the
    spectrum, the maximum-intensity fit, the Type~II band-splitting plots, the
    STEREO/SWAVES split view\index{STEREO/SWAVES!split view}, and any available GOES X-ray, GOES proton, Dst and Kp context
    plots;
  \item \emph{Full FITS Header Appendix}.
\end{itemize}
Every graph in the report is drawn in the same style as \gui{Export Figure}, on a white page
even in the dark theme.

\section{Batch processing}
\label{sec:batch}
\index{batch processing}

\menu{Processing > Batch Processing > Open Batch Processor} opens the \gui{Batch FIT
Processing} dialog (Figure~\ref{fig:batch}), which renders every FIT file in a folder as a
PNG image with consistent settings.

\screenshot[0.62\linewidth]{dialogs/batch_processing_dialog}{The \gui{Batch FIT Processing} dialog.}{fig:batch}{The Batch FIT Processing dialog with input and output folders selected, the Output Type, Processing Options and Aligned View groups, and the status line.}

\begin{description}
  \item[Input Folder, Output Folder] the folder of FIT files to process and the folder for
    the PNG images (one per input file, at 300~dpi; existing names are not overwritten).
  \item[Output Type] \gui{Raw} or \gui{Background Subtracted}.
  \item[Processing Options] the \gui{Background Method} (\gui{Mean}, \gui{Median} or
    \gui{Plotutil Median (dB)}\index{Plotutil}\index{background subtraction!Median (dB)}) and the \gui{Colormap}. Plotutil Median (dB) converts digits
    to dB with the CALLISTO scale before subtracting the median, and defaults to a display
    range of \(-1\) to 8~dB.
  \item[Aligned View] \gui{Use current display range} and \gui{Use current visual settings}
    apply the main window's display range and appearance; \gui{Load View Config...} applies
    a saved \file{.efaview.json}\index{efaview.json@.efaview.json file} file instead. These options give many station spectra
    identical time and frequency axes.
  \item[Start, Cancel, Close] start the run, cancel it after the current file, or close the
    dialog. A summary is shown when the run finishes, including any files that failed.
\end{description}


\setpartpreamble{\vspace{2em}\begin{center}\begin{minipage}{0.72\textwidth}\centering\large
Retrieving radio data from the e-CALLISTO archive, comparing stations, and
working with Learmonth and STEREO/SWAVES spectra.\end{minipage}\end{center}}
\part{Radio Data Access}
%% Chapter 11 ----------------------------------------------------------------
\chapter{The e-CALLISTO FITS Downloader}
\label{ch:downloader}
\index{FITS Downloader}\index{e-CALLISTO archive}

The e-CALLISTO FITS Downloader retrieves FITS files directly from the e-CALLISTO data
archive. Open it with the \gui{Download} toolbar tool, \menu{Download > Launch FITS
Downloader}, or \menu{Solar Events > Radio Bursts > e-CALLISTO}.

The window has four tabs --- \gui{Single Station}\index{FITS Downloader!Single Station}, \gui{Multi-Station Event}\index{FITS Downloader!Multi-Station Event},
\gui{Spectral Overview}\index{FITS Downloader!Spectral Overview} and \gui{Burst List}\index{FITS Downloader!Burst List} --- described in the sections below. Station
lists are read from the archive whenever a date changes, so they offer only the stations
that have data on the chosen UTC days; a line under the date reports how many, for example
\gui{42 stations have data on 2026-10-04 (UTC)}.

\section{Common features}

\subsection{The download cache}
\index{cache!downloader}

Previewed and downloaded files are kept in a persistent cache on disk, so re-previewing,
importing or extending a dataset never fetches the same file twice. The size of the cache is
shown at the bottom left of the window (\gui{Cache: …}). \gui{Open Cache Folder} opens the
cache in the file manager and \gui{Clear Cache}\index{cache!clearing} empties it.

\subsection{Preview, download, import and compare}

The tabs share four actions on the files you have ticked:
\begin{description}
  \item[Preview] opens the \gui{FITS Preview}\index{FITS Downloader!preview} window with a background-subtracted
    spectrogram of the selection. A selection that can be combined in time, in frequency or
    in both is previewed as it would be imported; otherwise each file appears in its own
    tab.
  \item[Download] asks for a folder (\gui{Select Download Folder}) and saves the files
    there.
  \item[Import] loads the selection into the main window, combining it automatically when
    possible (Section~\ref{sec:combine-rules}). Gap and overlap options are requested when
    needed.
  \item[Compare] opens the Multi-Station Comparison with the selection
    (Chapter~\ref{ch:comparison}).
\end{description}
During downloads a progress panel shows the overall progress, transfer statistics and an
estimated time remaining; \gui{Show file details}\index{FITS Downloader!download progress} expands a table with the progress, speed
and status of every file.

\section{Single Station}
\label{sec:single-station}

The \gui{Single Station} tab (Figure~\ref{fig:single-station}) lists all files of one
station for a complete UTC day.

\screenshot{downloader/single_station}{The \gui{Single Station} tab.}{fig:single-station}{Single Station tab.}

\begin{steps}
  \item In \gui{Observation Parameters}, choose the \gui{Date} and then the \gui{Station}.
  \item Click \gui{Show Available FITS}.
  \item In \gui{Available FITS Files (Whole Day, grouped by focus code)}, expand a focus
    code and tick files, or tick a whole focus-code group. Each group heading shows the
    number of files and the time span covered, for example
    \gui{Focus 63 — 34 files (00:00–08:14)}, so the receivers that covered an event are
    visible at a glance.
  \item Use the buttons in \gui{Actions}: \gui{Select All}, \gui{Deselect All},
    \gui{Download Selected}, \gui{Preview Selected}, \gui{Compare} or \gui{Import}.
\end{steps}

When several files are ticked, the downloader determines whether they can be combined in
time, in frequency or in both, using the rules of Section~\ref{sec:combine-rules}. Files that
cannot be combined are reported with an explanation (Figure~\ref{fig:combine-error}).

\section{Multi-Station Event}
\label{sec:multi-station-event}

The \gui{Multi-Station Event} tab (Figure~\ref{fig:multi-station-event}) finds the files
recorded by several stations during the same event in a single search.

\screenshot{downloader/multi_station_download}{The \gui{Multi-Station Event} tab.}{fig:multi-station-event}{Multi-Station Event tab.}

\begin{steps}
  \item In \gui{Event Window}, set \gui{Start (UTC)} and \gui{Stop (UTC)}. The station list
    is updated to the stations with data in that window (windows of up to seven days).
  \item In \gui{Stations}, tick the stations to search. \gui{Filter stations} narrows the
    list by name; \gui{Select All} and \gui{Clear} tick or untick all stations.
  \item Click \gui{Search Matching FITS}. The \gui{Matching FITS Files} table lists the
    station, UTC time, file name and receiver (focus code) of every file overlapping the
    window.
  \item Use \gui{Select All}, \gui{Clear Selection}, \gui{Download Selected}, \gui{Import}
    or \gui{Compare}.
\end{steps}

A compatible selection is imported directly into the main window, combined in time, in
frequency or in both. Files from several stations cannot be combined into one spectrum; use
\gui{Compare} to view them side by side. The comparison workspace opens only when
\gui{Compare} is clicked.

\section{Spectral Overview}
\label{sec:spectral-overview}
\index{spectral overview}

The \gui{Spectral Overview} tab (Figure~\ref{fig:spectral-overview}) assembles a station's
observations for a complete UTC day into one overview figure.

\screenshot{downloader/spectral_overview}{The \gui{Spectral Overview} tab with a full-day overview for one focus code.}{fig:spectral-overview}{Spectral Overview tab.}

\begin{description}
  \item[Date (UTC), Station] the day and station to summarize.
  \item[Focus code] \gui{All available focus codes}, or a single focus code to generate or
    regenerate on its own.
  \item[Generate Overview] downloads (or reuses from the cache) every file of the day and
    builds the overview. The status line reports the focus codes generated and the number of
    files loaded. \gui{Cancel} stops a running generation.
  \item[Overview Preview] one tab per focus code. Each overview shows the day as six
    four-hour panels (00--04~UT to 20--24~UT), background-subtracted with a day-wide median
    (dB) baseline and a common color scale; intervals without data are marked
    \emph{No data}.
  \item[Export...] saves the overview of the selected tab as a PNG, PDF, EPS, SVG or TIFF
    figure at 300~dpi on a white background.
\end{description}

\section{Burst List}
\label{sec:burst-list}
\index{burst list}\index{Monstein catalog}

The \gui{Burst List} tab (Figure~\ref{fig:burst-list}) gives access to the e-CALLISTO solar
radio burst catalog compiled by C.~Monstein \parencite{monsteinlist}, so that the FITS files
for a listed burst can be found together with its type and reporting stations. The
\gui{Monstein source catalog} link opens the original lists in a web browser.

\screenshot{downloader/burst_list}{The \gui{Burst List} tab with a selected Type~II event and its FITS files.}{fig:burst-list}{Burst List tab.}

\procedure{Catalog filters}
\begin{description}
  \item[From (UTC), Through (UTC)] the date range to load; click \gui{Load Events}.
  \item[Burst type] \gui{All types} or one of the burst types found in the loaded catalog (for example II or III).
  \item[Reporting station] events reported by a given station.
  \item[Min. stations] events reported by at least this many stations (\gui{Any} for no
    limit).
  \item[Search] free-text search in the type, stations or remarks.
  \item[Reset Filters] clears the filters while keeping the selected dates.
\end{description}
A progress bar with a \gui{Cancel} button reports catalog loading, FITS searches, preview
preparation and downloads. The final state remains visible after completion.

\procedure{1.~Choose a burst}
The event table lists the start and end time (UTC), type, number of stations, and the
observatories and remarks of each matching event, with a count such as
\gui{1 of 2 events match the filters}. Hovering over a row shows the original catalog
record and its source URL. Missing periods, source notices and loading errors are reported
below the table.

\procedure{2.~Preview or import FITS}
\begin{description}
  \item[Station] \gui{All reported stations}, one of them, or the name of any archive
    station typed in the box.
  \item[Extra time] padding added on both sides of the burst (0--120~min, default 2~min).
  \item[Find FITS] searches the archive. Overlapping 15-minute files and events that cross
    UTC midnight are included. The table lists the station, start time, focus code and file
    name of each file found.
  \item[Select All, Clear Selection, Preview, Download, Import, Compare] act on the ticked
    files as described above; the number of selected files is shown.
\end{description}
The divider between the two tables can be dragged to resize them.

\begin{note}
  Catalog coverage varies by year, station and type information may be unavailable in
  older lists, and a listed detection does not guarantee that matching FITS data are
  archived. Catalog times are UTC with one-minute precision.
\end{note}

%% Chapter 12 ----------------------------------------------------------------
\chapter{Multi-Station Comparison}
\label{ch:comparison}
\index{Multi-Station Comparison}

The Multi-Station Comparison workspace displays dynamic spectra from several stations in
vertically stacked panels with a shared time axis, so that the same event can be compared
across stations, receivers and frequency ranges (Figure~\ref{fig:comparison}). It is opened
from \menu{View > Multi-Station Comparison...} or with the \gui{Compare} button of the
downloader (Chapter~\ref{ch:downloader}).

\screenshot{downloader/multi_station_comparison}{The Multi-Station Comparison workspace with four stations observing the same event.}{fig:comparison}{Multi-Station Comparison window.}

The workspace follows the rendering mode of the application: in Modern mode the panels are
hardware-accelerated when available, and in Classic mode they are drawn with Matplotlib\index{Matplotlib}.

\section{Stations and files}

The \gui{Stations / Files} group lists the loaded files as \emph{station -- date -- file
name}. \gui{Add Files...} adds FITS files from disk, \gui{Remove} removes the selected entry
and \gui{Clear} removes all entries. Files that can be combined are merged automatically
before display: each station's files are combined in time, in frequency, or as a complete
time~\(\times\)~focus-code grid (Section~\ref{sec:combine-rules}), so a mixed-station
selection yields one panel per station. Hovering over an entry shows any warning raised
while reading it.

\section{Display controls}

\begin{table}[!htbp]
  \centering
  \caption{Display controls of the Multi-Station Comparison.}
  \label{tab:comparison-controls}
  \small
  \begin{tabularx}{\linewidth}{@{}P{0.15\linewidth}P{0.2\linewidth}L@{}}
    \toprule
    \thead{Group} & \thead{Control} & \thead{Function} \\
    \midrule
    View & Time alignment\index{Multi-Station Comparison!time alignment} & \gui{UT clock} aligns the panels on Universal Time; \gui{Seconds from file start} aligns each panel on its own start. UT alignment requires \texttt{TIME-OBS}\index{TIME-OBS keyword} in every file; otherwise the workspace switches to seconds and says so. \\
    Appearance & Units & \gui{Digits} or \gui{dB}. \\
               & Colormap & Color map of all panels. \\
               & Color scale\index{Multi-Station Comparison!color scale} & \gui{Shared scale} (one scale for all panels), \gui{Per-station auto scale}, or \gui{Manual scale} with \gui{Low} and \gui{High} limits. \\
    Noise Reduction & Target & \gui{All panels} or one panel. \\
                    & Method & \gui{None}, \gui{Mean background}, \gui{Median background}, \gui{Robust background} or \gui{Noise clipping}. \\
                    & Colormap & Color map of the noise-reduced panels. \\
                    & Thresholds & With \gui{Noise clipping}, \gui{Lower threshold} and \gui{Upper threshold} sliders and a \gui{Logarithmic Threshold Scale}\index{display thresholds!logarithmic scale} option, as in the main window. \\
    \bottomrule
  \end{tabularx}
\end{table}

\section{Display range and view configurations}

The \gui{Actions} group contains:
\begin{description}
  \item[Set Range...] enter a shared time and frequency range for all panels.
  \item[Reset Range] return to the full range of the data.
  \item[Load View Config...] apply a \file{.efaview.json}\index{efaview.json@.efaview.json file} view configuration
    (Section~\ref{sec:view-config}). The visual settings are applied; when the files differ
    in date, set UT bounds with \gui{Set Range...}.
  \item[Export...] export the comparison (Section~\ref{sec:comparison-export}).
\end{description}

\section{Ruler measurement}

In the \gui{Measurement} group, \gui{Start Ruler}\index{Multi-Station Comparison!ruler} activates a ruler: click two points in one
panel to measure the duration, frequency drift and slope, shown under \gui{Readout}.
\gui{Clear} removes the measurement.

\section{Exporting a comparison}
\label{sec:comparison-export}

\gui{Export...} opens the \gui{Export Comparison} dialog (Figure~\ref{fig:comparison-export}):
\begin{description}
  \item[Visible View] exports the comparison exactly as it appears in the window.
  \item[Grid / Table Layout] exports a compact publication-style matrix of all panels, with
    a shared time axis, each station's native frequency range and a white background. Its
    options are the number of \gui{Columns} (\gui{Auto} or a fixed number), an
    \gui{Overall title}, and the \gui{Raster resolution} (DPI) for bitmap formats.
\end{description}
The figure can be saved as PNG, PDF, EPS, SVG or TIFF.

\screenshot[0.5\linewidth]{dialogs/comparison_export_dialog}{The \gui{Export Comparison} dialog.}{fig:comparison-export}{The Export Comparison dialog with Grid / Table Layout selected, showing the Columns, Overall title and Raster resolution fields and the explanatory note.}

%% Chapter 13 ----------------------------------------------------------------
\chapter{Learmonth and STEREO/SWAVES Data}
\label{ch:other-radio}

Besides the e-CALLISTO archive, the application reads data from two further radio
instruments: the ground-based Learmonth Solar Radio Spectrograph in Western Australia, and
the space-based STEREO/WAVES (SWAVES) radio receivers. Both are opened from
\menu{Solar Events > Radio Bursts}.

\section{Learmonth Solar Radio Spectrograph}
\label{sec:learmonth}
\index{Learmonth}

The Learmonth Solar Radio Spectrograph, operated by the Australian Bureau of Meteorology,
records daily archive files (\file{.srs}\index{srs@.srs file}) in two receiver bands. \menu{Solar Events > Radio
Bursts > Learmonth} opens the \gui{Learmonth Solar Radio Browser}
(Figure~\ref{fig:learmonth}), which converts selected parts of a daily file into FIT files
for the main window.

\screenshot[0.7\linewidth]{radio-sources/learmonth_browser}{The Learmonth Solar Radio Browser.}{fig:learmonth}{The Learmonth Solar Radio Browser after Show Available Data, with the Daily Raw File information, the list of available 15-minute chunks (one highlighted for the chosen time) and the Convert buttons.}

\begin{description}
  \item[Observation Parameters] choose the \gui{Date} and a \gui{Time}, then click
    \gui{Show Available Data}. The daily file is read from the local cache when available
    and downloaded otherwise. The chunk containing the chosen time is highlighted.
  \item[Daily Raw File] information about the daily archive file. \gui{Download} saves
    the raw \file{.srs} file to a folder you choose.
  \item[Available 15-Minute Chunks] the time ranges into which the day is divided.
    \gui{Select All} and \gui{Deselect All} change the selection.
  \item[Convert to FIT] converts the selected chunks to CALLISTO-style FIT files in a folder
    you choose.
  \item[Convert and Import] converts the selected chunks and imports them into the main
    window, combined in time when consecutive.
\end{description}
Converted files carry a standard FITS header with the frequency and time axes and the
observatory position, so they can be processed and analyzed exactly like e-CALLISTO data.

\section{STEREO/SWAVES}
\label{sec:swaves}
\index{STEREO/SWAVES}\index{SWAVES}

SWAVES on the twin STEREO spacecraft \parencite{bougeret2008} records radio spectra from
2.6~kHz to 16~MHz, directly below the CALLISTO band. A burst that drifts out of the
ground-based range can therefore be followed into the interplanetary medium on the same
figure. \menu{Solar Events > Radio Bursts > SWAVES} opens the \gui{STEREO/SWAVES Radio
Spectrograph} dialog (Figure~\ref{fig:swaves-dialog}).

\screenshot[0.55\linewidth]{radio-sources/swaves_dialog}{The STEREO/SWAVES dialog.}{fig:swaves-dialog}{The STEREO/SWAVES Radio Spectrograph dialog with the Observation Window (UTC) fields, the Use CALLISTO Window, Cancel and Download and Plot buttons, and the summary line.}

\begin{description}
  \item[Start date, Start time, Duration] the UTC window to load (0.1--24~h; default
    4~h). The data are NASA/SPDF level-2 one-minute averages served as one file per UTC day;
    a window that crosses midnight fetches and joins both days. The archive starts on
    27~October~2006, and downloaded day files are cached and reused.
  \item[Spacecraft] \gui{STEREO-A (Ahead)}\index{STEREO} or \gui{STEREO-B (Behind)}. STEREO-B is
    unavailable for dates after 1~October~2014, when contact with the spacecraft was lost.
  \item[Sync padding] extra time added on each side when the window is taken from the
    CALLISTO spectrum (0--720~min; default 30~min). A Type~II or Type~III burst takes tens
    of minutes to hours to drift from the corona into the SWAVES band, so an exact match
    would cut off the part you want to see.
  \item[Use CALLISTO Window] copies the time range of the loaded CALLISTO spectrum, widened
    by the sync padding. \menu{Solar Events > Sync Current Time Window}\index{time synchronization} updates an open
    SWAVES dialog in the same way.
  \item[Download and Plot] loads the data; \gui{Cancel} stops a running load.
\end{description}

\subsection{The split view}
\label{sec:swaves-split}

Once SWAVES data are loaded, the plotting area of the main window splits into two panels:
CALLISTO above and SWAVES below, on a shared time axis, so that panning or zooming either
panel moves both (Figure~\ref{fig:swaves-split}). The CALLISTO interval is outlined on the
SWAVES panel. The SWAVES frequency axis is always logarithmic and its intensity is in
decibels above the instrument background, so the panel keeps its own color bar and color
scaling while following the color map chosen in the sidebar. Loading SWAVES with no
CALLISTO file open plots it across the whole area; it moves into the split view\index{STEREO/SWAVES!split view} as soon as a
FITS file is loaded.

\screenshot{radio-sources/swaves_split_view}{CALLISTO and STEREO/SWAVES spectra in the split view.}{fig:swaves-split}{Main window in split view: a CALLISTO spectrum with a Type II/III burst in the upper panel and the SWAVES spectrum below on the shared time axis, with the CALLISTO interval outlined.}

\menu{Solar Events > SWAVES Panel}\index{STEREO/SWAVES!panel} hides and restores the panel without discarding the
loaded data. The split view is included in exported figures, saved in and restored from
projects without downloading the data again, and added to project reports.

\begin{note}
  The drawing, lasso, drift and measurement tools act on the CALLISTO panel only, and the
  data-reduction controls of the sidebar do not apply to the SWAVES panel.
\end{note}


\setpartpreamble{\vspace{2em}\begin{center}\begin{minipage}{0.72\textwidth}\centering\large
X-ray, particle, geomagnetic and CME catalog context for a radio event, and
the SunPy Multi-Mission Explorer.\end{minipage}\end{center}}
\part{Solar-Event Context}
%% Chapter 14 ----------------------------------------------------------------
\chapter{Solar-Event Viewers}
\label{ch:solar-events}
\index{solar-event viewers}

The viewers in the \menu{Solar Events} menu place a radio event in its wider context: the
X-ray flare, energetic protons, the geomagnetic response and any associated coronal mass
ejection. Each viewer is an independent window that fetches its data from the archive
listed in Appendix~\ref{app:data-sources}. The viewers can also be started from source as
standalone programs (Section~\ref{sec:from-source}).

\section{Common controls}

The GOES and geomagnetic viewers share a common layout:
\begin{itemize}
  \item a time-range selector with \gui{Start} and \gui{End} rows (date, and hour, minute or
    three-hour slot) and, where applicable, the spacecraft;
  \item a \gui{Preset} list with \gui{Apply} for common ranges;
  \item a plot button that fetches and plots the data, with a progress indicator;
  \item \gui{Save Plot} (PNG) and \gui{Save Data}\index{Burst Analyzer!Excel export} (CSV);
  \item a reset or clear button; and
  \item an information panel summarizing the plotted data or a selected interval.
\end{itemize}

\section{Synchronizing the time window}
\label{sec:sync-time}
\index{time synchronization}

\menu{Solar Events > Sync Current Time Window}\index{time synchronization} sends the time range of the spectrum loaded in
the main window to every open context window: the GOES X-ray and proton viewers, the Dst and
Kp viewers and the STEREO/SWAVES dialog plot the new range; the SOHO/LASCO\index{SOHO/LASCO} CME catalog
searches the day of the middle of the range; and the SunPy Explorer takes the range as its
query window. The status bar reports which windows were updated, or that none were open.

\section{GOES X-ray flux}
\label{sec:goes-xrs}
\index{GOES!X-ray flux}

\menu{Solar Events > Flares > GOES X-Ray Flux} opens the \gui{GOES XRS Plotter}
(Figure~\ref{fig:goes-xrs}), which plots the soft X-ray flux measured by the GOES X-Ray
Sensor in its short (XRS-A, 0.05--0.4~nm) and long (XRS-B, 0.1--0.8~nm) channels.

\screenshot{space-weather/goes_xrs}{The GOES XRS Plotter with a selected flare interval.}{fig:goes-xrs}{The GOES XRS Plotter showing XRS-A and XRS-B on a logarithmic flux axis with flare-class levels, a dragged selection over a flare, and the Flare Information panel filled in.}

\begin{description}
  \item[Start, End, Use] the UTC range and the spacecraft (GOES-8 to GOES-19; default
    GOES-16). \gui{Preset: Whole day (00:00–23:59)} selects the whole of the start day.
  \item[Plot XRS Data] downloads the daily science files from NOAA and plots both channels
    on a logarithmic flux axis.
  \item[Selection] drag across the plot to select a flare interval. The
    \gui{Flare Information}\index{GOES!flare information} panel then shows the \gui{Selection window}, the
    \gui{Peak X-ray Flux (W/m²)} and \gui{Peak Time (UTC)} of the long channel, the
    \gui{Rise Time} (selection start to peak), the \gui{Decay Time} (peak to selection end),
    the \gui{Flare Duration} and the \gui{GOES Class}\index{GOES!flare class} derived from the peak flux.
  \item[Save Plot, Save Data, Reset] save the figure as PNG, save the fluxes as CSV, or clear
    the plot and selection.
\end{description}

\subsection{GOES overlay on the spectrum}
\label{sec:goes-overlay}
\index{GOES!overlay}

\menu{Solar Events > GOES Overlay > Short(XRS-A)}\index{GOES!overlay} and \menu{Long(XRS-B)} draw the GOES
X-ray curves directly on the dynamic spectrum in the main window
(Figure~\ref{fig:goes-overlay}), with a flare-class guide (A, B, C, M, X) on the right. The
overlay does not modify the spectrum data. The application chooses a GOES satellite
appropriate for the date automatically, falling back across legacy and modern satellites
when a file is unavailable. The overlay is included in exported figures.

\screenshot{space-weather/goes_overlay}{GOES X-ray curves overlaid on a dynamic spectrum.}{fig:goes-overlay}{Main window with a dynamic spectrum and both GOES Overlay channels enabled, showing the right-hand flare-class guide (A/B/C/M/X).}

\section{GOES proton flux}
\label{sec:goes-sep}
\index{GOES!proton flux}\index{solar energetic particles}

\menu{Solar Events > Energetic Particles > GOES SEP Proton Flux} opens the
\gui{GOES SEP Proton Flux Plotter} (Figure~\ref{fig:goes-sep}), which plots proton fluxes
from the GOES Solar and Galactic Proton Sensor (SGPS) archive at NOAA over multi-day UTC
ranges.

\screenshot{space-weather/goes_sep}{The GOES SEP Proton Flux Plotter.}{fig:goes-sep}{The GOES SEP Proton Flux Plotter for a multi-day range during an SEP event, with the ~10 MeV and ~100 MeV channels plotted, a selected interval, and the SEP Analysis panel filled in.}

\begin{description}
  \item[Start, End, Use] the UTC range and the spacecraft: \gui{Auto (GOES-19 -> GOES-16)}
    tries the newest spacecraft first, or a specific GOES satellite can be chosen.
  \item[Preset] \gui{Last 24 hours}, \gui{Last 72 hours}, \gui{Last 7 days} or
    \gui{Current UTC day}.
  \item[Plot SEP Proton Flux] fetches and stitches the daily files and plots the channels
    closest to about 10~MeV and 100~MeV, with hover readouts.
  \item[Selection] drag across the plot to measure an event.
  \item[SEP Analysis] the spacecraft, channels, number of samples and source files, the
    peaks of both channels in the plot, and for the selection its window, peak (low
    channel), peak time, rise time, decay time and duration.
  \item[Save Plot, Save Data, Clear Selection] save the figure as PNG, save the stitched
    flux table as CSV, or clear the selection.
\end{description}

\section{Kyoto Dst index}
\label{sec:dst}
\index{Dst index}

\menu{Solar Events > Geomagnetic > Kyoto Dst Index} opens the \gui{Dst Index} viewer
(Figure~\ref{fig:dst}) for the hourly Dst index from the World Data Center for
Geomagnetism, Kyoto \parencite{nose2015}. Dst measures the strength of the ring current;
large negative values indicate geomagnetic storms.

\screenshot{space-weather/dst_index}{The Kyoto Dst Index viewer.}{fig:dst}{The Dst Index viewer for a storm interval, showing the Dst curve with the Quiet, −50, −100 and −200 nT guide lines and the Dst Summary panel.}

\begin{description}
  \item[Start, End] date and hour. \gui{Preset}: \gui{Last 24 hours}, \gui{Last 7 days},
    \gui{Current month} or \gui{Previous month}.
  \item[Plot Dst Data] fetches and plots the index with storm-level guides at 0 (quiet),
    \(-50\), \(-100\) and \(-200\)~nT.
  \item[Dst Summary] the number of samples, the data sources (final, provisional or
    real-time), and the minimum, maximum (each with its time) and mean Dst.
  \item[Save Plot, Save Data, Reset] PNG figure, CSV table, or clear.
\end{description}

\section{GFZ Kp index}
\label{sec:kp}
\index{Kp index}

\menu{Solar Events > Geomagnetic > GFZ Kp Index} opens the \gui{Kp Index} viewer
(Figure~\ref{fig:kp}) for the three-hourly planetary Kp index from GFZ Potsdam
\parencite{matzka2021}.

\screenshot{space-weather/kp_index}{The GFZ Kp Index viewer.}{fig:kp}{The Kp Index viewer for a storm interval, showing colored Kp bars with the Kp 5, 7 and 8 guide lines and the Kp Summary panel.}

\begin{description}
  \item[Start, End] date and three-hour slot (\gui{3h Slot}). \gui{Preset}:
    \gui{Last 24 hours}, \gui{Last 7 days}, \gui{Current month} or \gui{Previous month}.
  \item[Plot Kp Data] plots the index as bars colored by activity (quiet below 5, storm from
    5, severe from 7), with guide lines at Kp~5, 7 and 8.
  \item[Kp Summary] the number of intervals, the maximum, minimum and mean Kp, and the
    number of storm intervals (\(\mathrm{Kp}\geq5\)).
  \item[Save Plot, Save Data, Reset] PNG figure, CSV table, or clear.
\end{description}

\section{SOHO/LASCO CME catalog}
\label{sec:lasco-catalog}
\index{CME catalog}\index{SOHO/LASCO}

\menu{Solar Events > CMEs > SOHO/LASCO CME Catalog} opens the
\gui{SOHO/LASCO CME Catalog Tool} (Figure~\ref{fig:lasco-catalog}), which reads the CDAW\index{CDAW}
catalog of CMEs observed by SOHO/LASCO \parencite{yashiro2004}.

\screenshot{space-weather/lasco_cme_catalog}{The SOHO/LASCO CME Catalog Tool.}{fig:lasco-catalog}{The SOHO/LASCO CME Catalog Tool after Search for a day with several CMEs, with one CME selected, its details shown at the right and the movie status panel with the Play Movie, Open in Browser and Copy URL buttons.}

\begin{description}
  \item[Year, Month, Day, Search] list the CMEs of the chosen day. A progress bar shows the
    retrieval.
  \item[CME table] one row per CME with the columns \gui{Date and Time}, \gui{Central PA},
    \gui{Angular Width}, \gui{Linear Speed}, \gui{Accel}, \gui{Mass},
    \gui{Kinetic Energy}, \gui{MPA} (measurement position angle) and \gui{Remarks}.
    Clicking a row shows its details beside the table.
  \item[Play Movie] plays the LASCO movie associated with the selected CME (double-clicking a
    row does the same) in a viewer with \gui{Back}, \gui{Forward}, \gui{Reload},
    \gui{Open in Browser}, \gui{Open Raw Movie} and \gui{Copy URL}.
  \item[Open in Browser, Copy URL] open the CDAW page of the selected CME, or copy its
    address.
\end{description}

%% Chapter 15 ----------------------------------------------------------------
\chapter{SunPy Multi-Mission Explorer}
\label{ch:sunpy}
\index{SunPy Multi-Mission Explorer}

The SunPy Multi-Mission Explorer searches public solar archives through SunPy
\parencite{sunpy2020}, downloads the selected records into an application-managed cache,
and plots images and time series for quick inspection. Open it with
\menu{Solar Events > Archives > SunPy Multi-Mission Explorer}. For in-depth image analysis,
use Solar Image Analysis (Part~V).

\screenshot[0.8\linewidth]{sunpy/sunpy_explorer}{The SunPy Multi-Mission Solar Explorer.}{fig:sunpy-explorer}{The SunPy Multi-Mission Solar Explorer after a search, with the Archive Query group filled in, records listed in Search Results (some ticked), and the Analysis and Export group.}

\section{Supported instruments}

Table~\ref{tab:sunpy-instruments} lists the instruments that can be searched. Map
(image) products are shown in the plot window as image sequences; GOES/XRS is a time
series.

\begin{table}[!htbp]
  \centering
  \caption{Instruments available in the SunPy Multi-Mission Explorer.}
  \label{tab:sunpy-instruments}
  \small
  \begin{tabularx}{\linewidth}{@{}P{0.27\linewidth}P{0.17\linewidth}L@{}}
    \toprule
    \thead{Instrument} & \thead{Type} & \thead{Selectable options} \\
    \midrule
    SDO/AIA\index{SDO/AIA} & Map & Wavelength 94, 131, 171, 193, 211, 304, 335, 1600, 1700~Å \\
    SDO/HMI\index{SDO/HMI} & Map & Product: magnetogram, continuum, dopplergram \\
    SOHO/EIT\index{SOHO/EIT} & Map & Wavelength 171, 195, 284, 304~Å \\
    SOHO/LASCO\index{SOHO/LASCO} C2, C3 & Map & Detector \\
    STEREO-A\index{STEREO}/B SECCHI & Map & Detector EUVI (171, 195, 284, 304~Å), COR1, COR2, HI1, HI2 \\
    GOES/SUVI\index{GOES/SUVI} & Map & Wavelength 94, 131, 171, 195, 284, 304~Å; level 1b or 2 \\
    PROBA2/SWAP\index{PROBA2/SWAP} & Map & --- \\
    GOES/XRS & Time series & GOES satellite \\
    \bottomrule
  \end{tabularx}
\end{table}

STEREO-B entries return data only for 2007--2014, before contact with the spacecraft was
lost.

\section{Searching the archives}

The \gui{Archive Query} group contains:
\begin{description}
  \item[Spacecraft, Instrument] the observatory and instrument. The \gui{Detector},
    \gui{Wavelength (A)}, \gui{GOES Satellite}, product and \gui{Level} fields appear when
    they apply to the selected instrument.
  \item[Sample (s)] the sampling cadence in seconds; 0 requests the full archive cadence,
    which is slower and returns more records.
  \item[Start (UTC), End (UTC), Max Records] the time range and the maximum number of
    records returned.
  \item[Search Archives] runs the query. \gui{Retry Last Query} repeats it, for example after
    a network error.
  \item[Use Analyzer Time Window] sets the time range from the spectrum in the main window and
    searches.
  \item[Open Cache Folder, Clear Cache] manage the download cache; its size is shown.
\end{description}

The \gui{Search Results} table lists, for each record, the start and end time, source,
instrument, provider, file identifier and size. Tick the records to download with the
\gui{Use} column or \gui{Select All} and \gui{Clear Selection}, then click
\gui{Download \& Load Selected}\index{Solar Image Analysis!Load Selected}. \gui{Retry Failed} downloads again only the rows that
failed, and \gui{Stop} cancels a running search or download. Cached files can be reopened
without a network connection.

\section{The plot window}

\gui{Open Plot Window} in the \gui{Analysis and Export} group opens the
\gui{SunPy Plot and Animation}\index{SunPy Multi-Mission Explorer!plot window} window with the loaded data (Figure~\ref{fig:sunpy-plot}).

\screenshot[0.72\linewidth]{sunpy/sunpy_plot_window}{The SunPy Plot and Animation window in map mode.}{fig:sunpy-plot}{The SunPy Plot and Animation window showing a coronagraph or EUV image sequence with the frame slider, the Run Diff, AIA Limb (EUVI), NRGF, playback, FPS and ROI controls.}

In \emph{map mode} the controls are:
\begin{description}
  \item[Frame slider] steps through the sequence; the frame number is shown.
  \item[Run Diff] shows running differences\index{running difference} between consecutive frames.
    \menu{Difference > Base Difference}\index{base difference} shows each frame minus the first.
  \item[AIA Limb (EUVI)] on STEREO/EUVI images, draws the solar limb as seen from Earth,
    showing which part of the disk is visible to SDO/AIA.
  \item[NRGF] applies the normalizing-radial-graded filter \parencite{morgan2006} to
    coronagraph images (COR1, COR2, LASCO), flattening the radial fall-off of brightness to
    reveal faint CME fronts.
  \item[Play, Pause, Rewind, FPS] animate the sequence at 1--30 frames per second.
  \item[Select ROI, Reset ROI] draw a rectangular region of interest on the map, for which
    the minimum, maximum, mean, median, standard deviation and 95th and 99th percentiles are
    computed.
\end{description}
In \emph{time-series mode} (GOES/XRS) the window plots the X-ray channels and derives basic
flare summary metrics.

\section{Exporting}

The \gui{Analysis and Export} group of the explorer provides \gui{Save Files As...}, which
copies the downloaded data files to a folder of your choice; \gui{Export Plot}, which saves
the current plot as PNG, PDF, SVG, TIFF or JPG; and \gui{Export Analysis CSV}, which saves
the ROI statistics or flare summary.


\setpartpreamble{\vspace{2em}\begin{center}\begin{minipage}{0.72\textwidth}\centering\large
The multi-mission imaging workspace: obtaining and displaying solar images,
measuring eruptions, and modeling the coronal magnetic field.\end{minipage}\end{center}}
\part{Solar Image Analysis}
%% Chapter 16 ----------------------------------------------------------------
\chapter{The Solar Image Analysis Window}
\label{ch:solar-window}
\index{Solar Image Analysis}

Solar Image Analysis is a multi-mission imaging workspace for extreme-ultraviolet (EUV)
disk images, magnetograms, coronagraph images and heliospheric images. It retrieves image
sequences from the solar archives or reads local FITS files, displays and processes them,
measures eruptions and their kinematics, and models the coronal magnetic field\index{magnetic field}. Open it with
\menu{Analysis > Solar Image Analysis} in the main window.

This chapter describes the window, its menus and sessions, and how image data are obtained.
Display and processing are described in Chapter~\ref{ch:solar-processing}, measurements in
Chapter~\ref{ch:solar-measurements} and magnetic-field modeling in Chapter~\ref{ch:pfss}.

\section{Supported observables}
\label{sec:observables}
\index{observable}

Table~\ref{tab:observables} lists the observables that can be retrieved. Local FITS files
from these and other solar instruments can be loaded as well.

\begin{table}[!htbp]
  \centering
  \caption{Observables available in Solar Image Analysis.}
  \label{tab:observables}
  \small
  \begin{tabularx}{\linewidth}{@{}P{0.25\linewidth}P{0.18\linewidth}L@{}}
    \toprule
    \thead{Instrument} & \thead{Class} & \thead{Observables} \\
    \midrule
    SDO/AIA\index{SDO/AIA} \parencite{lemen2012} & EUV/UV disk & 94, 131, 171, 193, 211, 304, 335, 1600 and 1700~Å \\
    SDO/HMI\index{SDO/HMI} \parencite{scherrer2012} & Magnetograph & Magnetogram, Continuum (Intensitygram), Dopplergram; vector field (Section~\ref{sec:hmi-vector}) \\
    SOHO/EIT\index{SOHO/EIT} & EUV disk & 171, 195, 284 and 304~Å \\
    SOHO/LASCO\index{SOHO/LASCO} \parencite{brueckner1995} & Coronagraph & C2, C3 \\
    STEREO-A\index{STEREO}/B SECCHI \parencite{howard2008} & EUV disk, coronagraph, heliospheric imager & EUVI 171, 195, 284 and 304~Å; COR1, COR2; HI1, HI2 \\
    GOES/SUVI\index{GOES/SUVI} & EUV disk & 94, 131, 171, 195, 284 and 304~Å \\
    \bottomrule
  \end{tabularx}
\end{table}

The sidebar shows only the tool groups that apply to the instrument class of the loaded data
(or, before loading, of the selected observable); see Table~\ref{tab:solar-cards}.

\section{Window layout}
\label{sec:solar-layout}

The window is divided into the regions numbered in Figure~\ref{fig:solar-window} and
described in Table~\ref{tab:solar-regions}.

\screenshot{solar/solar_window_overview}{The Solar Image Analysis window with its regions numbered.}{fig:solar-window}{Solar Image Analysis window with an AIA 193 Å sequence loaded and Measurements enabled. Add numbered callouts: 1 menu bar, 2 sidebar cards, 3 sidebar handle, 4 coordinate readout and Export MP4, 5 Measure toolbar, 6 image canvas, 7 CME tracking panel, 8 playback bar, 9 Details drawer.}

\begin{table}[!htbp]
  \centering
  \caption{Regions of the Solar Image Analysis window.}
  \label{tab:solar-regions}
  \small
  \begin{tabularx}{\linewidth}{@{}cP{0.22\linewidth}L@{}}
    \toprule
    \thead{No.} & \thead{Region} & \thead{Description} \\
    \midrule
    \callout{1} & Menu bar & \menu{Session}, \menu{Data}, \menu{Analysis}, \menu{Movie}, \menu{Export} and \menu{About} (Section~\ref{sec:solar-menus}). \\
    \callout{2} & Sidebar & Collapsible cards that read from top to bottom as the workflow: get data, choose records, analyze, then display, export and instrument-specific tools. \\
    \callout{3} & Sidebar handle & Hides or shows the sidebar. \\
    \callout{4} & Header bar & Live cursor readout --- distance from Sun center in \(R_\odot\), position angle, heliographic coordinates on the disk, and pixel --- and the one-click \gui{Export MP4}\index{movie export!MP4} button. \\
    \callout{5} & Measure toolbar & The \gui{Measurements}\index{measurement tools!enabling} switch, the \gui{Ruler}\index{measurement tools!ruler}, \gui{Profile}\index{measurement tools!profile}, \gui{Region Stats}\index{measurement tools!region statistics}, \gui{Track CME}\index{Track CME}, \gui{Circle Fit}\index{circle fit} and \gui{Clear} tools, and \gui{Pan / Zoom}\index{Solar Image Analysis!pan and zoom} (Chapter~\ref{ch:solar-measurements}). \\
    \callout{6} & Image canvas & The current frame, titled with instrument, wavelength and time. \\
    \callout{7} & CME tracking panel & Measurement table, live height--time graph and kinematic fit; active when \gui{Measurements} is ticked (Section~\ref{sec:tracking-panel}). \\
    \callout{8} & Playback bar & Rewind, previous, play, pause and next buttons, the frame scrubber, the frame counter (\gui{Frame n / N}) and the \gui{Speed} in frames per second. \\
    \callout{9} & Details drawer\index{Solar Image Analysis!Details drawer} & Results and warnings of every operation. Click \gui{Details} to collapse it; drag the divider above it to resize. \\
    \bottomrule
  \end{tabularx}
\end{table}

The playback speed set in the playback bar is used for playback and for exported movies.
While a sequence plays, or while a contrast slider is dragged, the canvas switches to a fast
preview that subsamples the display; full resolution returns as soon as playback stops.
Measurements, statistics and exports always use the full-resolution data.

\subsection{Menus}
\label{sec:solar-menus}

\begin{xltabular}{\linewidth}{@{}P{0.38\linewidth} L@{}}
  \caption{Menus of the Solar Image Analysis window.}\label{tab:solar-menus}\\
  \toprule \thead{Command} & \thead{Function} \\ \midrule\endfirsthead
  \toprule \thead{Command} & \thead{Function} \\ \midrule\endhead
  \menu{Session > Open Session…} & Open an \file{.ecsolar}\index{ecsolar@.ecsolar file} session (Section~\ref{sec:solar-sessions}). \\
  \menu{Session > Save Session} / \menu{Save Session As…} & Save the session (\kbd{Ctrl+S} / \kbd{Ctrl+Shift+S}). \\
  \menu{Data > Fetch Archive Records} & Search the archive (\gui{Fetch Records}\index{Solar Image Analysis!Fetch Records}). \\
  \menu{Data > Find Latest Available}\index{Solar Image Analysis!Find Latest} & Find the newest available data (\gui{Find Latest}). \\
  \menu{Data > SOHO Live Preview (Helioviewer)}\index{Helioviewer} & Near-real-time SOHO quicklook images (Section~\ref{sec:live-preview}). \\
  \menu{Data > Load Selected (to cache)}\index{Solar Image Analysis!Load Selected} & Download and load the ticked records. \\
  \menu{Data > Save Selected to Disk…} & Download the ticked records to a folder. \\
  \menu{Data > Upload FITS Files}\index{Solar Image Analysis!local files} & Load local FITS files. \\
  \menu{Data > Use Analyzer Window} & Copy the time range of the main window's spectrum. \\
  \menu{Data > Stop Download/Search} & Cancel the running operation. \\
  \menu{Data > Reset All (clear cache \& defaults)} & Clear the cache and restore default settings. \\
  \menu{Analysis > Plot Frames}, \menu{Running Difference}\index{running difference}, \menu{Composite} & Display modes (Section~\ref{sec:solar-analysis-card}). \\
  \menu{Analysis > Identify Active Regions}, \menu{Fetch NOAA/HEK Labels}\index{active regions!NOAA labels} & Active-region tools (Section~\ref{sec:active-regions}). \\
  \menu{Analysis > Compare Viewpoint…} & Compare with a second observer (Section~\ref{sec:compare-viewpoint}). \\
  \menu{Analysis > GCS CME Fitting…} & Open GCS fitting at the time of the displayed frame (Chapter~\ref{ch:gcs-window}). \\
  \menu{Analysis > Reset Loaded Frames} & Restore the frames as loaded. \\
  \menu{Movie > Rewind}, \menu{Back}, \menu{Play}, \menu{Pause}, \menu{Forward} & Playback. \\
  \menu{Movie > Build Movie}\index{movie export} & Export a movie (Section~\ref{sec:solar-movies}). \\
  \menu{Export > Export Plot}, \menu{Export Cropped FITS}\index{crop!export FITS}, \menu{Export Regions CSV}, \menu{Export MP4} & Exports (Section~\ref{sec:solar-movies}). \\
  \menu{About > About Solar Image Analysis…} & Version information. \\
  \bottomrule
\end{xltabular}

\subsection{Sidebar cards}

Table~\ref{tab:solar-cards} lists the sidebar cards in order, and the instrument classes for
which each is shown.

\begin{table}[!htbp]
  \centering
  \caption{Sidebar cards of Solar Image Analysis.}
  \label{tab:solar-cards}
  \small
  \begin{tabularx}{\linewidth}{@{}P{0.3\linewidth}P{0.3\linewidth}L@{}}
    \toprule
    \thead{Card} & \thead{Shown for} & \thead{See} \\
    \midrule
    Data Source & all & Section~\ref{sec:data-source} \\
    Archive Results & all & Section~\ref{sec:archive-results} \\
    Analysis & all (disk tools for disk imagers and magnetographs only) & Section~\ref{sec:solar-analysis-card} \\
    Processing Level & AIA and SUVI & Section~\ref{sec:processing-level} \\
    Display \& Crop & all & Section~\ref{sec:display-crop} \\
    Solar Derotation & all & Section~\ref{sec:derotation} \\
    Movie Export & all & Section~\ref{sec:solar-movies} \\
    Coronagraph Tools & coronagraphs & Section~\ref{sec:nrgf} \\
    Overlay Layers & coronagraphs & Section~\ref{sec:overlay-layers} \\
    Heliospheric Imager (J-map) & HI1, HI2 & Section~\ref{sec:jmap} \\
    Magnetic Vector Field (HMI) & HMI & Section~\ref{sec:hmi-vector} \\
    Potential Field (PFSS) & disk imagers, magnetographs, synoptic maps & Chapter~\ref{ch:pfss} \\
    Active Regions & disk imagers, magnetographs & Section~\ref{sec:active-regions} \\
    \bottomrule
  \end{tabularx}
\end{table}

\section{Obtaining image data}

\subsection{The Data Source card}
\label{sec:data-source}

The \gui{Data Source} card (Figure~\ref{fig:solar-data-source}) defines what to retrieve and
how.

\screenshot[0.5\linewidth]{solar/solar_data_source}{The \gui{Data Source} card.}{fig:solar-data-source}{The Data Source card (Observable, Start/End, Cadence, Max records, Use Analyzer Time Window, Download options, buttons) above the Archive Results card with a list of records, some ticked.}

\begin{description}
  \item[Observable] the instrument and wavelength or product (Table~\ref{tab:observables}).
  \item[Start (UTC), End (UTC)] the time range.
  \item[Cadence (s)] the sampling interval between frames (0--3600~s; default 120~s).
  \item[Max records] the maximum number of records to list (1--5000; default 120).
  \item[Use Analyzer Time Window] copies the start and end times of the spectrum in the main
    window, so that the images match the radio burst being studied.
\end{description}

\procedure{Download options (SDO)}
\begin{description}
  \item[Source] \gui{Auto (JSOC → VSO)}\index{VSO}\index{JSOC} tries the Joint Science Operations Center (JSOC)
    first and falls back to the Virtual Solar Observatory (VSO); \gui{JSOC (fast)} returns
    direct, compressed AIA files and is usually much faster; \gui{VSO (classic)} uses the
    VSO only.
  \item[JSOC notify e-mail] JSOC requires a one-time, free registration of an e-mail address
    at \url{https://jsoc.stanford.edu/ajax/register_email.html}. The address is remembered.
  \item[Frame size] \gui{Full disk (4096²)}, \gui{Binned ½ (2048²)}, \gui{Binned ¼ (1024²)}
    or \gui{Cutout (region)}\index{Solar Image Analysis!cutout}. Binned and cutout frames are produced by JSOC, so only the
    reduced data are transferred; they require the JSOC source and e-mail. For a cutout,
    enter the \gui{Centre X″}, \gui{Centre Y″}, \gui{Width″} and \gui{Height″} in arcseconds
    from disk center.
  \item[High resolution AIA (best crop)] requests full-resolution AIA products. Files are
    larger, but cropped views keep more detail.
\end{description}
An estimate of the download size is shown below the options when records are selected.

\procedure{Actions}
\begin{description}
  \item[Fetch Records] searches the archive; the matching records appear in
    \gui{Archive Results}. \gui{Stop} cancels a running search or download; files already
    completed remain in the cache.
  \item[Find Latest] scans the archive backwards to find the newest data actually available
    for the observable. This is useful for SOHO/LASCO, whose calibrated archive can lag real
    time by many months, and may take a minute or two.
  \item[Upload FITS…] loads solar FITS files from disk without an archive search.
  \item[Live Preview (Helioviewer)] shown for SOHO/LASCO and SOHO/EIT
    (Section~\ref{sec:live-preview}).
\end{description}

\subsection{The Archive Results card}
\label{sec:archive-results}

The \gui{Archive Results} table lists the records found, with the columns \gui{Use},
\gui{Start UTC}, \gui{Source}, \gui{Size} and \gui{File ID}. Tick records in the \gui{Use}
column, or use \gui{Select All} and \gui{Deselect All}, then:
\begin{itemize}
  \item \gui{Load Selected} downloads the ticked records into the working cache and loads
    them for analysis; or
  \item \gui{Save to Disk…} downloads them to a folder you choose, for keeping.
\end{itemize}
The \gui{Analysis} card then summarizes what is loaded (instrument, number of frames and
configuration).

\begin{note}[SOHO data]
  The archive serves two files for every SOHO/EIT observation, the calibrated Level~1 file
  and the raw level-zero file. The window chooses one per observation: Level~1 where it
  exists and the raw file otherwise. Level~1 processing runs about a year behind the raw
  archive, so recent searches return raw frames and older ones calibrated frames; a window
  that straddles the boundary loads as two configuration groups. The calibrated SOHO/LASCO
  archive also lags real time. For a date beyond the newest archived data, the search falls
  back to the nearest available data and says so. Use \gui{Live Preview} to see the most
  recent SOHO images.
\end{note}

\subsection{Live Preview (Helioviewer)}
\label{sec:live-preview}
\index{Helioviewer}

\gui{Live Preview (Helioviewer)} opens the \gui{SOHO/… Near-Real-Time Preview (Helioviewer)}
window (Figure~\ref{fig:live-preview}) with the newest quicklook image of SOHO/LASCO C2 or C3
or SOHO/EIT, usually updated within about an hour. These browse images are for situational
awareness, not for analysis.

\screenshot[0.5\linewidth]{solar/helioviewer_preview}{The SOHO near-real-time preview.}{fig:live-preview}{The SOHO/LASCO Near-Real-Time Preview (Helioviewer) window showing the latest C2 image, the Detector selector, the Latest Frame button, and the Movie (time range · step) group with playback controls.}

\begin{description}
  \item[Passband / Detector, Latest Frame] choose the EIT passband or LASCO detector, and
    reload the newest image.
  \item[Movie (time range · step)] set \gui{Start (UTC)}, \gui{End (UTC)}, \gui{Step}
    (minutes) and \gui{Max frames}, then \gui{Build Movie} to fetch one frame per step and
    play them with \gui{◀}, \gui{Play}, \gui{Pause}, \gui{▶} and an FPS setting.
    \gui{Cancel} stops the fetch. The frame cap protects the Helioviewer service.
  \item[Open in Browser] opens the view on the Helioviewer web site.
\end{description}

\subsection{Processing level}
\label{sec:processing-level}
\index{calibration level}

Downloads always fetch the archive's base product. For instruments that publish more than
one processing level, the \gui{Processing Level} card selects the level that is displayed and
analyzed:
\begin{description}
  \item[SDO/AIA] \gui{Level 1 (as downloaded)}, the archive product; or
    \gui{Level 1.5 (registered here)}\index{calibration level!AIA level 1.5}, computed locally with aiapy \parencite{barnes2020},
    which applies the master pointing, rotates the image to solar north, rescales it to
    0.6~arcsec per pixel and centers the Sun.
  \item[GOES/SUVI] \gui{Level 1b (as downloaded)}, single calibrated exposures; or
    \gui{Level 2 (archive composite)}, NOAA's high-dynamic-range composites, which are
    downloaded separately when selected.
\end{description}
\gui{Apply Level}\index{calibration level!Apply Level} re-derives the loaded sequence at the selected level; any crop or
derotation is re-applied afterwards. The status line reports the level in use. The card is
hidden for instruments that offer only one level.

\section{Sessions}
\label{sec:solar-sessions}
\index{session!solar}\index{ecsolar@.ecsolar file}

A Solar Image Analysis session (\file{.ecsolar}) is a single file that reopens an analysis
exactly where it was left. It embeds the original FITS frames, so it survives a cleared
cache or a move to another computer (at the cost of a large file), together with the data
source, the display, crop and playback state, the CME measurements (height--time picks,
circle fits and fit order) and the PFSS settings. The PFSS solution itself is not stored; one
click on \gui{Compute PFSS}\index{PFSS!compute} reproduces it. Sessions are opened and saved from the
\menu{Session} menu (\kbd{Ctrl+S}, \kbd{Ctrl+Shift+S}).

%% Chapter 17 ----------------------------------------------------------------
\chapter{Image Display and Processing}
\label{ch:solar-processing}

This chapter describes the tools of Solar Image Analysis that display, enhance, combine and
export image sequences once they are loaded (Chapter~\ref{ch:solar-window}).

\section{The Analysis card}
\label{sec:solar-analysis-card}

The \gui{Analysis} card (Figure~\ref{fig:solar-cards}a) shows a summary of the loaded data and
the main display modes and analysis tools.

\begin{figure}[!htbp]
  \centering
  \begin{minipage}[b]{0.47\linewidth}\centering
    \includegraphics[width=\linewidth]{figures/solar/solar_analysis_cards.png}\par\smallskip
    {\small (a) \gui{Analysis} card}
  \end{minipage}\hfill
  \begin{minipage}[b]{0.47\linewidth}\centering
    \includegraphics[width=\linewidth]{figures/solar/solar_display_cards.png}\par\smallskip
    {\small (b) \gui{Display \& Crop} card}
  \end{minipage}
  \caption{The \gui{Analysis} and \gui{Display \& Crop} cards of Solar Image Analysis.}
  \label{fig:solar-cards}
\end{figure}

\begin{description}
  \item[Plot Frames] displays the loaded sequence as plain frames.
  \item[Running Diff] displays each frame minus the previous one, which makes moving and
    changing features stand out.
  \item[Base Diff] displays each frame minus the first, showing the total change since the
    start --- useful for eruptions, coronal dimmings and EUV waves.
  \item[Composite] builds an RGB tri-color image from the first three AIA frames or, when a
    magnetogram overlay has been loaded, draws HMI polarity contours on the current AIA frame
    (red positive, blue negative).
  \item[Mag Overlay…] loads an HMI magnetogram FITS file for the composite; the box beside it
    sets the contour level (10--2000~G; default 100~G).
  \item[Light Curve] plots the mean intensity of a region against time
    (Section~\ref{sec:region-lightcurve}).
  \item[Active Regions], \gui{NOAA/HEK Labels}\index{active regions!NOAA labels} detect and label active regions
    (Section~\ref{sec:active-regions}).
  \item[Compare Viewpoint…] compares the frame with a second observer
    (Section~\ref{sec:compare-viewpoint}).
  \item[Reset Frames] restores the frames as downloaded, discarding crops, processing
    levels, derotation, composites and detected regions.
\end{description}
\gui{Composite}\index{composite image}, \gui{Mag Overlay…}\index{magnetogram overlay}, \gui{Active Regions} and \gui{NOAA/HEK Labels} are
available only for disk imagers and magnetographs. The difference modes also determine the
content of exported movies (Section~\ref{sec:solar-movies}).

\section{Display and cropping}
\label{sec:display-crop}
\index{display!solar images}\index{crop}

The \gui{Display \& Crop} card (Figure~\ref{fig:solar-cards}b) controls the appearance of the image:
\begin{description}
  \item[Renderer] \gui{PyQtGraph}\index{PyQtGraph} (default; fast, and required for clicked PFSS seeds and the
    interactive crop rectangle) or \gui{Matplotlib}\index{Matplotlib}.
  \item[Colormap] instrument color maps (for example \texttt{sdoaia193}, \texttt{hmimag},
    \texttt{soholasco2}, \texttt{stereocor2}, \texttt{euvi195}, \texttt{goes-rsuvi171}) and
    general ones (\texttt{inferno}, \texttt{magma}, \texttt{plasma}, \texttt{viridis},
    \texttt{cividis}, \texttt{gray}, \texttt{RdBu\_r}, \texttt{hot}). A suitable map is
    chosen automatically for each observable.
  \item[Scale] \gui{linear} or \gui{log}.
  \item[Clip low, Clip high] the lower and upper percentiles mapped to the ends of the color
    scale (defaults 1\% and 99.9\%); dragging gives a live preview.
  \item[Solar Limb] draws the photospheric limb.
  \item[Coordinate Grid, Grid frame] overlays a graticule of meridians and parallels in the
    selected frame --- Heliocentric Inertial (HCI, the default), Stonyhurst or Carrington. The
    same frame is used for the heliographic coordinates of the cursor readout.
  \item[Colorbar, Region Overlays] show the color bar and the detected active regions.
\end{description}

\procedure{Cropping}
Under \gui{Crop (arcsec)}, enter \gui{X min / max} and \gui{Y min / max} in arcseconds from
disk center, or tick \gui{Rectangle crop}\index{crop!rectangle} to drag a rectangle on the image that fills in
the fields. \gui{Apply Crop} crops all loaded frames to the rectangle. Cropping is done
locally after the files are loaded. \gui{Export Cropped FITS}\index{crop!export FITS} saves the cropped frames as
FITS files with updated coordinate metadata.

The crop rectangle also defines the region used by \gui{Region Stats}\index{measurement tools!region statistics}, the light curve, the
derotation and the \gui{Current crop} PFSS seeding.

\section{Solar derotation}
\label{sec:derotation}
\index{derotation}\index{differential rotation}

Over a sequence of several hours, a fixed rectangle slowly drifts onto neighboring
surface because of solar differential rotation. The \gui{Solar Derotation} card keeps a
selected region on the same piece of the Sun:
\begin{description}
  \item[Derotate selected area] enables derotation of the crop region.
  \item[Mode] \gui{Track region (no resampling)}\index{derotation!track mode} moves the cut-out window with the rotation
    and keeps the original pixel values, so it is the mode for photometry and light curves.
    \gui{Reproject to reference time}\index{derotation!reproject mode} resamples every frame onto the reference time's grid,
    which corrects foreshortening and gives all frames one coordinate system, making
    differencing meaningful, at the cost of interpolating the data.
  \item[Reference] \gui{First frame}, \gui{Current frame} or \gui{Specific frame...} (with its
    number): the frame from which the region is read and to whose time every other frame is
    rotated.
  \item[Padding] the extra margin cut around the region before reprojecting (0--600~arcsec;
    default 60~arcsec), so that the interpolation has data up to the edge. Reproject mode only.
  \item[Apply Derotation, Reset to Full Frames] apply the derotation, or discard the cut-out
    and restore the full frames at the current processing level.
\end{description}
The status line reports whether derotation is active.

\section{Region light curve}
\label{sec:region-lightcurve}
\index{light curve!EUV}

\gui{Light Curve} plots the mean intensity of the crop region (or of the whole frame if
\gui{Rectangle crop} is off), converted to DN/s using the exposure time, against time
(Figure~\ref{fig:region-lightcurve}). The time of peak intensity is marked and, when a
spectrum is loaded in the main window, its time window is shaded so that the EUV brightening
can be timed against the radio burst. At least two frames are needed.

\screenshot[0.72\linewidth]{solar/region_light_curve}{The \gui{Region Light Curve} window.}{fig:region-lightcurve}{The Region Light Curve window for an AIA flare region, showing the DN/s curve, the EUV peak marker and the shaded radio burst window.}

\section{Active regions}
\label{sec:active-regions}
\index{active regions}

The \gui{Active Regions} card detects bright regions in the current frame and lists them:
\begin{description}
  \item[Threshold] the intensity percentile above which pixels count as bright (50--100;
    default 98).
  \item[Min px] the minimum area of a region in pixels (default 12).
  \item[Region table] one row per region with \gui{ID}, \gui{Label}, \gui{NOAA} number,
    \gui{Centroid}, \gui{Area}, \gui{Peak} intensity and \gui{Source}.
\end{description}
Click \gui{Active Regions} in the \gui{Analysis} card to run the detection, and
\gui{NOAA/HEK Labels} to fetch the NOAA active-region numbers from the Heliophysics Event
Knowledge base (HEK) and attach them to the detected regions; the status line reports whether
metadata are loaded. Detected regions are drawn on the image when \gui{Region Overlays}\index{active regions!overlays} is
ticked. \gui{Export Regions CSV} saves the table with centroids, bounding boxes and
intensities. Detection works on local files; the labels need network access.

\section{Comparing viewpoints}
\label{sec:compare-viewpoint}
\index{Compare Viewpoint}

\gui{Compare Viewpoint…} opens a window that compares the current frame (viewpoint~A) with a
second observable (viewpoint~B) taken near the same time --- for example a STEREO image
against the Earth view.
\begin{steps}
  \item Choose the \gui{Second viewpoint} observable and the search window (\gui{±}, 1--720~min;
    default 30~min).
  \item Click \gui{Fetch Viewpoint B}. Viewpoint~B is downloaded, reprojected onto the
    coordinate frame of A, and shown beside it.
  \item Tick \gui{Blink}\index{Compare Viewpoint!blink} to alternate between the two images in place.
\end{steps}
The \gui{Observer separation}\index{Compare Viewpoint!observer separation} line gives the angle between the two observers.

\section{Coronagraph tools}
\label{sec:nrgf}
\index{NRGF}

For coronagraph data the \gui{Coronagraph Tools} card offers \gui{NRGF radial filter}\index{NRGF}, the
normalizing-radial-graded filter \parencite{morgan2006}. It flattens the steep radial fall-off
of coronal brightness so that faint CME fronts become visible at all heights. It applies to
the plain frame view; difference images already remove the background.

\section{Overlay layers}
\label{sec:overlay-layers}
\index{overlay layers}\index{coronagraph composite}

A CME is never fully visible in one instrument: the eruption starts on the disk (AIA, EIT,
EUVI), crosses the low corona (LASCO~C2, COR1) and expands into the outer corona (LASCO~C3,
COR2), and each coronagraph is blind inside its occulter. For coronagraph views the
\gui{Overlay Layers} card blends several instruments into one continuous image from the disk
to the outer corona.

\begin{description}
  \item[Add, Remove] choose an observable and add it to the layer list, or remove the
    selected layer. Layers are painted widest field of view first. Untick a layer to leave it
    out without losing its settings.
  \item[Per-layer settings] for the selected layer: \gui{Colormap}, \gui{Scale} (log or
    linear), \gui{Opacity}, \gui{Midtones}\index{overlay layers!midtones} (gamma; 50 suits most solar images), and the
    \gui{Inner} and \gui{Outer} edges of its field of view in \(R_\odot\). Pixels inside the
    inner edge stay transparent so that the layer below shows through.
  \item[Time match] how far from each loaded frame's time a layer frame may be taken
    (1--720~min; default 30~min). Widen it for slow cadences or patchy archives.
  \item[Build Composite] downloads each layer's matching frames, reprojects them onto the
    loaded series and blends them, in the background with a progress bar. The result replaces
    the view and works with the measurement tools, plot export and movie export.
  \item[Clear] drops the composite and restores the originally loaded frames.
\end{description}

\begin{note}
  Each layer is color-mapped before blending, because EUV disk images and the white-light
  corona differ by orders of magnitude. Coronagraph layers are reprojected under a
  spherical-screen assumption, flagged as approximate in the build notes: cross-observer
  coronagraph overlays are morphological context, not photometry. SOHO/LASCO\index{SOHO/LASCO} files carry no
  observer position, so an Earth-based observer is assumed (accurate to about 1\%).
\end{note}

\section{Heliospheric imager J-maps}
\label{sec:jmap}
\index{J-map}\index{heliospheric imager}

For STEREO HI1 and HI2 data the \gui{Heliospheric Imager (J-map)} card builds a
time--elongation map \parencite{sheeley1999}:
\begin{description}
  \item[Background] \gui{Median background} removes the static F-corona and star field with
    the per-pixel temporal median; \gui{Previous frame} uses a running difference\index{running difference}.
  \item[PA] the position angle of the slit, counter-clockwise from the image \(+x\) axis
    (90° points up; default 90°). Set it along the CME direction.
  \item[Build J-map] stacks the slit profiles of all frames into the \gui{HI J-map} window
    (elongation in degrees against frame index). An outward-moving CME appears as a slanted
    bright track.
\end{description}

\section{HMI vector magnetic field}
\label{sec:hmi-vector}
\index{HMI!vector field}

For SDO/HMI\index{SDO/HMI} data the \gui{Magnetic Vector Field (HMI)} card overlays the measured
photospheric vector field:
\begin{description}
  \item[Load Vector FITS…] loads the \texttt{hmi.B\_720s} field, inclination and azimuth
    segments of each time step from disk (with the optional disambiguation segment, applied
    automatically).
  \item[Download Vector (JSOC)] downloads \texttt{hmi.B\_720s} for the query time window
    through JSOC\index{JSOC}; this needs the registered JSOC e-mail. Each 720-s time step is four segments
    of about 50~MB. If no HMI image is loaded, the derived vertical-field (\(B_z\))
    magnetogram is plotted with the vectors.
  \item[Show vector field] overlays the field on the HMI frame. The transverse component is
    drawn as \gui{Arrows} (length proportional to strength) and/or \gui{Streamlines}\index{HMI!streamlines}, and
    \gui{|B| layer} tints the image by total field strength. Arrows are red where the vertical
    field points toward the observer (\(+B_z\)) and blue where it points away.
  \item[Spacing, Threshold] the arrow grid spacing (16--512 detector pixels; default 64) and
    the minimum transverse field drawn (50--2000~G; default 200~G). The transverse noise floor
    of \texttt{hmi.B\_720s} is about 100~G.
\end{description}

\section{Movies and still images}
\label{sec:solar-movies}
\index{movie export}

\begin{description}
  \item[Movie Export card] choose the \gui{Content} (\gui{Frames}, \gui{Running Difference}
    or \gui{Base Difference}\index{base difference}) and the \gui{Format} (\gui{MP4} or \gui{GIF}), and click
    \gui{Build Movie…}\index{movie export}. Movies use the playback speed set under the viewer and the current
    color map, scale, clipping and crop.
  \item[Export MP4] in the header bar makes a one-click MP4 of the loaded sequence with the
    current display settings.
  \item[Export Plot] saves the current frame, with its overlays, as an image.
  \item[Export Cropped FITS] saves the cropped frames as FITS files.
  \item[Export Regions CSV] saves the active-region table.
\end{description}
The same exports are available from the \menu{Export} and \menu{Movie} menus.

%% Chapter 18 ----------------------------------------------------------------
\chapter{Measurements and CME Kinematics}
\label{ch:solar-measurements}
\index{measurement tools}\index{CME kinematics}

The Measure toolbar above the image provides click-driven tools for distances, intensity
profiles, region statistics and CME tracking. The CME tracking panel to the right of the
image collects the measurements of a sequence and fits their kinematics.

\section{Enabling the tools}

Tick \gui{Measurements}\index{measurement tools!enabling} in the Measure toolbar to enable the tools; the CME tracking panel
becomes available at the same time. One tool is active at a time. A right-click or \kbd{Esc}
cancels a pick in progress without leaving the tool. Results are shown in the status line and
appended to the Details drawer\index{Solar Image Analysis!Details drawer}. \gui{Pan / Zoom}\index{Solar Image Analysis!pan and zoom}, at the right of the toolbar, enables
dragging and scroll-wheel zooming of the image; the zoom is kept while the movie plays, and
unticking it returns to the full frame.

\begin{table}[!htbp]
  \centering
  \caption{Tools of the Measure toolbar.}
  \label{tab:measure-tools}
  \small
  \begin{tabularx}{\linewidth}{@{}P{0.17\linewidth}L@{}}
    \toprule
    \thead{Tool} & \thead{What it does} \\
    \midrule
    \gui{Ruler}\index{measurement tools!ruler} & Two clicks give the plane-of-sky distance in arcseconds, \(R_\odot\) and Mm, and the position angle (measured from north toward east). \\
    \gui{Profile}\index{measurement tools!profile} & Two clicks plot the intensity along the cut, for example across a loop, a filament or a CME front, in the \gui{Intensity Profile} window. \\
    \gui{Region Stats}\index{measurement tools!region statistics} & Summarizes the crop rectangle: pixel count, mean, median, minimum, maximum, standard deviation and intensity-weighted centroid (arcsec). Enable \gui{Rectangle crop}\index{crop!rectangle} first to choose the region. \\
    \gui{Track CME}\index{Track CME} & One click per frame on the leading edge of a CME. The height is its distance from disk center. Best for limb events. \\
    \gui{Circle Fit}\index{circle fit} & Three or more clicks along a circular front per frame. Best for eruptions near disk center. \\
    \gui{Clear} & Resets every measurement: picks, circles, table and overlays. \\
    \bottomrule
  \end{tabularx}
\end{table}

The 3-D reconstruction of CMEs from several viewpoints is performed in the separate GCS CME
Fitting window (Part~VI), opened from \menu{Analysis > GCS CME Fitting…}.

\section{Tracking a CME}
\label{sec:track-cme}
\index{Track CME}

\gui{Track CME} measures a height--time profile with one click per frame
(Figure~\ref{fig:track-cme}).
\begin{steps}
  \item Go to the first frame in which the CME front is visible.
  \item Select \gui{Track CME} and click the leading edge. The pick is added to the tracking
    table and, with \gui{Auto-advance frame after each pick}\index{Track CME!auto-advance} ticked (the default), the next
    frame is shown.
  \item Repeat on each frame, then choose the fit order and click \gui{Fit Height–Time}\index{height--time fit}.
\end{steps}
The height is the plane-of-sky distance of the clicked point from Sun center in \(R_\odot\);
the position angle (PA) is recorded with it. The tool works on any imager --- EUVI, AIA,
SUVI, COR, LASCO and HI. Clicking again on a frame replaces its pick.

\screenshot{solar/track_cme}{CME tracking with the height--time graph in the tracking panel.}{fig:track-cme}{A LASCO C2 running-difference frame with leading-edge picks, and the CME tracking panel showing the table, the height–time plot with a quadratic fit, and the speed/acceleration text.}

\section{Fitting a circle to a CME front}
\label{sec:circle-fit}
\index{circle fit}

For eruptions near disk center the front expands as a growing dome rather than moving
outwards from the limb, and a single leading-edge click has no well-defined origin. \gui{Circle
Fit} models the CME as a uniformly expanding sphere resting on the solar surface
(Figure~\ref{fig:circle-fit}).
\begin{steps}
  \item Select \gui{Circle Fit} and click three or more points along the visible front. The
    least-squares circle is drawn from the third click and tightens with each further point.
  \item Click \gui{Commit Circle}\index{circle fit!commit} (\kbd{Ctrl+Return}). The frame's result is added to the table
    and the next frame is shown.
  \item Repeat on each frame, then click \gui{Fit Height–Time}.
\end{steps}
Each row gives the height \(h = 1\,R_\odot + 2r\), where \(r\) is the fitted radius (the
sphere's top is one diameter above the surface), the radius in \(R_\odot\), the position angle
of the center and the number of points. The fit residual, the center and the leading-edge
distance from disk center appear in the row's tooltip. Re-committing a frame replaces its
entry.
\begin{itemize}
  \item Click a wide arc. A short arc is fitted, but it constrains the radius very weakly, and
    the panel warns when the arc spans less than about 30°.
  \item \gui{Lock centre}\index{circle fit!lock center} freezes the center at its current fitted value so that later frames
    fit the radius alone. Use it once the dome has stopped drifting.
\end{itemize}
The fit uses Taubin\index{circle fit!Taubin method}'s method \parencite{taubin1991}, which is not biased toward small radii on
partial arcs, and works in helioprojective arcseconds so that the front remains a circle when
the plate scale differs between the axes.

\screenshot{solar/circle_fit}{Circle fitting of an on-disk CME front.}{fig:circle-fit}{An AIA or EUVI difference image of a disk-center eruption with clicked points and the fitted circle, and the tracking panel showing circle-fit rows, the Lock centre and Commit Circle controls.}

\section{The CME tracking panel}
\label{sec:tracking-panel}
\index{tracking panel}

The tracking panel shows the measurements of the active tool:
\begin{description}
  \item[Table] for \gui{Track CME}: \gui{Time (UT)}, \gui{t (s)} since the first pick,
    \gui{Height (R☉)} and \gui{PA (°)}; for \gui{Circle Fit}: additionally \gui{Radius (R☉)}
    and the number of points \gui{N}.
  \item[Height--time graph] the measured heights against time, with the live fit.
  \item[Fit] the polynomial order: \gui{Linear}\index{height--time fit!order}, \gui{Quadratic} or \gui{Cubic}. Changing it
    refits the points already in the graph.
  \item[Fit Height–Time] fits the points and reports the speed and acceleration below the
    graph.
  \item[Clear Picks / Clear Circles] removes the measurements of the active tool.
  \item[Export] \gui{Table as CSV…}, or \gui{Graph (PNG, PDF, EPS, SVG, TIFF, JPG)…} drawn in
    light mode and the OriginPro style.
\end{description}

\begin{background}[Background: height--time fits]
  \begin{description}
    \item[Linear] \(h = h_0 + v t\): one plane-of-sky speed for the whole track. The
      acceleration is taken from a companion quadratic fit, as in the CDAW\index{CDAW} CME catalog.
    \item[Quadratic] \(h = h_0 + v_0 t + \tfrac12 a t^2\): constant acceleration, with the
      speed reported at the first and last frame.
    \item[Cubic] adds a constant jerk, with the speed and acceleration reported at both ends.
  \end{description}
  Every value carries a \(1\sigma\) uncertainty propagated from the fit covariance, so the
  correlations between the polynomial terms are included. A fit of degree \(n\) needs \(n+1\)
  points, and \(n+2\) before an uncertainty can be estimated; with fewer points the value is
  shown without an error rather than with a misleading zero. Heights from \gui{Track CME} are
  plane-of-sky distances, so the derived speeds generally underestimate the true radial
  speed; \gui{Circle Fit} heights rely on the assumption of a sphere expanding from the
  surface.
\end{background}

The measurements --- picks, circles, the lock state and the fit order --- are saved in
\file{.ecsolar}\index{ecsolar@.ecsolar file} sessions (Section~\ref{sec:solar-sessions}).

%% Chapter 19 ----------------------------------------------------------------
\chapter{Coronal Magnetic Field Modeling (PFSS)}
\label{ch:pfss}
\index{PFSS}\index{potential field source surface}

An EUV image shows where the coronal plasma is, not where the magnetic field\index{magnetic field} is. The
\gui{Potential Field (PFSS)} card of Solar Image Analysis extrapolates the coronal magnetic
field from a photospheric magnetogram with the potential-field source-surface model
\parencite{schatten1969,altschuler1969} and draws the resulting field lines over the disk
(Figure~\ref{fig:pfss}), so that the modeled loop topology and coronal holes can be compared
with the image. The card is shown for disk imagers, magnetographs and synoptic maps.

\screenshot{solar/pfss_overlay}{PFSS field lines over an AIA image.}{fig:pfss}{An AIA 193 Å image with PFSS open field lines (red/blue), gray closed loops and open-field boundaries drawn, and the Potential Field (PFSS) card visible with its settings and status line.}

\begin{background}[Background: the PFSS model]
  The model assumes a current-free (potential) field between the photosphere and a spherical
  \emph{source surface} a few solar radii out, where the field is forced to be radial to
  mimic the solar wind dragging it open. Field lines that reach the source surface are
  \emph{open} and feed the solar wind; the others are \emph{closed} loops. The open-field
  footpoint regions are the model's coronal holes, and the polarity-inversion line on the
  source surface marks where the heliospheric current sheet leaves the model.
\end{background}

\begin{caution}
  AIA images carry no magnetic-field information, so the boundary condition comes from a
  \emph{synoptic} magnetogram --- a full-Sun map assembled from central-meridian observations
  over a whole solar rotation. The model therefore describes the global field around your
  observation, never the instantaneous field, and the far side is always at least two weeks
  old. The card shows how far the map's date is from the displayed frame; that offset is the
  number to judge the result by.
\end{caution}

\section{Settings}

\begin{table}[!htbp]
  \centering
  \caption{Settings of the \gui{Potential Field (PFSS)} card.}
  \label{tab:pfss}
  \small
  \begin{tabularx}{\linewidth}{@{}P{0.19\linewidth}L@{}}
    \toprule
    \thead{Setting} & \thead{Description} \\
    \midrule
    Magnetogram & The boundary condition (Section~\ref{sec:pfss-magnetograms}). For \gui{GONG ADAPT}\index{PFSS!ADAPT}, \gui{Realization} (0--11) chooses one of the twelve flux-transport realizations; for \gui{Local FITS file...}, \gui{Browse…} selects the file. The line below names the loaded map and its time offset from the frame. \\
    Source surface\index{PFSS!source surface} & Height at which the field is forced radial (1.5--10~\(R_\odot\); default 2.5~\(R_\odot\)). Lowering it opens more field. \\
    Radial cells\index{PFSS!radial cells} & Radial grid resolution between the photosphere and the source surface (10--150; default 35). More is smoother and slower. \\
    Seeds & Where field lines start (Section~\ref{sec:pfss-seeds}). \\
    Density & Seed density (4--64; default 12, about 150 field lines). It controls both speed and legibility; the hint beside it estimates the number of lines. \\
    \bottomrule
  \end{tabularx}
\end{table}

\subsection{Magnetograms}
\label{sec:pfss-magnetograms}
\index{magnetogram!synoptic}

\begin{description}
  \item[GONG synoptic] the ground-based GONG synoptic\index{PFSS!GONG} map \parencite{harvey1996}, updated daily
    and available for any date from 2006 on. No registration is needed; this is the usual
    choice.
  \item[HMI synoptic (JSOC)] SDO's own synoptic map, instrumentally consistent with AIA.
    Searching is free, but downloading stages a JSOC\index{JSOC} export and needs the registered e-mail
    set in the \gui{Data Source} card. The series is published about one Carrington rotation
    behind, so the most recent weeks are not covered.
  \item[GONG ADAPT] maps from the ADAPT flux-transport model, which estimates the unobserved
    far side instead of leaving it weeks old and usually improves open-field predictions. One
    file holds twelve realizations, which differ mostly on the far side.
  \item[Local FITS file...] a synoptic map you already have. It must be a full-Sun Carrington
    map; plate-carrée maps are reprojected to equal-area automatically.
\end{description}
Downloaded magnetograms are cached, so changing a model parameter and solving again does not
download the map again. Unmeasured pixels, such as the polar gaps in GONG maps, are filled
with zero so that the solution can be computed; their number is reported.

\subsection{Seeds}
\label{sec:pfss-seeds}

\begin{description}
  \item[Uniform grid] an equal-area sample over the visible hemisphere.
  \item[Open-field regions] a coarse survey is traced first, and then lines are started only
    where the field is open. This shows coronal-hole connectivity without closed loops
    dominating the picture.
  \item[Current crop] concentrates the lines inside the crop rectangle.
  \item[Clicked points] traces one line from each point clicked on the disk. This needs the
    PyQtGraph\index{PyQtGraph} renderer.
\end{description}
Tracing costs about 1.6~ms per field line, so a solution takes a few seconds at most, but
beyond about 150 lines the field lines start to hide the image they are compared with. Prefer
\gui{Open-field regions}\index{PFSS!seeds} to a high density when coronal-hole connectivity is what you want.

\section{Computing and displaying the model}

\gui{Compute PFSS}\index{PFSS!compute} downloads the magnetogram if needed, solves the potential field and traces
the field lines in the background; \gui{Clear} removes the overlay and the solution. The
\gui{Show} options toggle the layers without recomputing:
\begin{description}
  \item[Open field] lines reaching the source surface, colored by the sign of the radial field
    at their footpoint: red outward, blue inward.
  \item[Closed loops] lines returning to the photosphere, drawn in gray.
  \item[Open-field boundaries] outlines of the open-field footpoint regions, for comparison
    with the dark coronal holes in 193 or 211~Å images.
  \item[Near side only] (on by default) hides lines rooted on the far hemisphere. Such lines
    are not hidden behind the disk once they rise above the limb, so they would otherwise be
    drawn as arcs outside the disk with no visible footpoint. Clear it for the geometrically
    complete picture.
\end{description}
Stepping through a sequence re-projects the same solution onto each frame instead of solving
again. This is correct, because the synoptic map is valid across the rotation while the
observer's view changes from frame to frame.

\section{Diagnostics}

\gui{Diagnostics…}\index{PFSS!diagnostics} opens a four-panel window that documents the solution
(Figure~\ref{fig:pfss-diagnostics}):
\begin{enumerate}[label=\arabic*.]
  \item the input magnetogram (\(B_r\)) with its provenance;
  \item the radial field at the source surface with its neutral line, where the heliospheric
    current sheet leaves the model --- the model's most directly testable prediction;
  \item the open (red and blue) and closed (dark) footpoint map in Carrington projection;
  \item the solution statistics, including the open flux, the open-area fraction and the
    number of filled pixels.
\end{enumerate}
\gui{Save Figure…} saves all four panels as one image (PNG, PDF or SVG).

\screenshot{solar/pfss_diagnostics}{The PFSS diagnostics window.}{fig:pfss-diagnostics}{The PFSS diagnostics window showing the four panels: input magnetogram, source-surface Br with neutral line, open/closed footpoint map, and solution statistics.}

\begin{note}
  PFSS modeling uses the \texttt{sunkit-magex}\index{sunkit-magex} package \parencite{stansby2020}, which is
  included in the installers. When the application is run from source without it, the card
  remains visible but disabled and shows the command to install it. Sessions store the PFSS
  settings, not the solution (Section~\ref{sec:solar-sessions}).
\end{note}


\setpartpreamble{\vspace{2em}\begin{center}\begin{minipage}{0.72\textwidth}\centering\large
Multi-viewpoint reconstruction of coronal mass ejections and their shocks
with the Graduated Cylindrical Shell and spheroid/ellipsoid
models.\end{minipage}\end{center}}
\part{Three-Dimensional CME Reconstruction}
%% Chapter 20 ----------------------------------------------------------------
\chapter{The GCS CME Fitting Window}
\label{ch:gcs-window}
\index{GCS CME Fitting}\index{Graduated Cylindrical Shell}

A single coronagraph view cannot separate a CME's direction from its size: a wide CME aimed
at the observer and a narrow one crossing the sky can look alike. Two well-separated views
break that ambiguity, and a third over-constrains it. The GCS CME Fitting window reconstructs
a CME in three dimensions by fitting one Graduated Cylindrical Shell (GCS) model
\parencite{thernisien2006,thernisien2011} to up to three coronagraph viewpoints at once, and
optionally a spheroid or ellipsoid to the shock the CME drives
\parencite{kouloumvakos2019,kwon2014}. The parameters follow the conventions of PyThea\index{PyThea}
\parencite{kouloumvakos2022}, so results can be compared directly.

This chapter describes the window and how images are loaded and viewed. Fitting, recording
and exporting are described in Chapter~\ref{ch:gcs-fitting}.

\section{Opening the window}

Open the window with \menu{Analysis > GCS CME Fitting…} in the main window or in Solar Image
Analysis. The window downloads its own images, so nothing needs to be loaded first:
\begin{itemize}
  \item from the main window, the visible time range of the spectrum becomes the event range;
  \item from Solar Image Analysis, the event range is centered on the displayed frame, one hour
    either side.
\end{itemize}
Choosing the menu item again reuses the open window and passes it the new time.

\section{Window layout}

The window is divided into the regions numbered in Figure~\ref{fig:gcs-window} and described
in Table~\ref{tab:gcs-regions}.

\screenshot{gcs/gcs_window_overview}{The GCS CME Fitting window in the equal layout.}{fig:gcs-window}{GCS CME Fitting window with three viewpoints loaded (STEREO-B or another channel, LASCO C2, STEREO-A COR2) in running difference, the GCS wireframe drawn, and the four control cards below. Add numbered callouts: 1 menu bar, 2 toolbar, 3 viewpoint panels A, B and C, 4 Event \& channels card, 5 Model card, 6 Display card, 7 Kinematics card, 8 status bar.}

\begin{table}[!htbp]
  \centering
  \caption{Regions of the GCS CME Fitting window.}
  \label{tab:gcs-regions}
  \small
  \begin{tabularx}{\linewidth}{@{}cP{0.22\linewidth}L@{}}
    \toprule
    \thead{No.} & \thead{Region} & \thead{Description} \\
    \midrule
    \callout{1} & Menu bar & \menu{File}, \menu{Event}, \menu{View}, \menu{Fit} and \menu{Help} (Section~\ref{sec:gcs-menus}). \\
    \callout{2} & Toolbar & Playback, the shared time, image mode, layout and overlay toggles (Section~\ref{sec:gcs-toolbar}). \\
    \callout{3} & Viewpoint panels & Panels A, B and C, each showing one channel's image nearest the shared time, with the models drawn on it. \\
    \callout{4} & Event \& channels card & Event range, frame cap, time-offset limit and the three channels (Section~\ref{sec:gcs-event-card}). \\
    \callout{5} & Model card & The edited model, its parameters, front-point controls, \gui{Refine fit}\index{GCS!refine fit} and \gui{Commit} (Chapter~\ref{ch:gcs-fitting}). \\
    \callout{6} & Display card & Wireframe colors and style, and each image's color map and contrast (Section~\ref{sec:gcs-display-card}). \\
    \callout{7} & Kinematics card & Recorded fits, height--time graph and fit (Section~\ref{sec:gcs-kinematics}). \\
    \callout{8} & Status bar & Messages, including the views taking part in a fit and their separations, and the helioprojective position under the cursor. \\
    \bottomrule
  \end{tabularx}
\end{table}

\section{Toolbar}
\label{sec:gcs-toolbar}

\begin{description}
  \item[Playback] first frame (\kbd{Home}), previous (\kbd{←}), play and pause (\kbd{Space}),
    next (\kbd{→}), and the playback speed (1--30~fps; default 4~fps). With fits recorded, the
    model follows them during playback.
  \item[Time slider and time label] step through the event. Every panel shows its own image
    nearest the selected UTC time, which is displayed beside the slider.
  \item[View] \gui{Raw} (the images as downloaded), \gui{Running diff}\index{GCS!image modes} (each image minus the
    previous one; selected whenever images load) or \gui{Base diff} (each image minus the first
    of the sequence).
  \item[Layout] \gui{Equal}\index{GCS!layout} shows three equal images with the cards beneath (\kbd{0});
    \gui{A}, \gui{B} or \gui{C} enlarges that panel, stacks the other two beside it and moves
    the cards to the right (\kbd{1}--\kbd{3}).
  \item[Solar limb, Axes, Banner] outline the photosphere in every image (on by default; the
    reference against which a CME's height is judged), show arcsecond axes around each image,
    and show the information banner across the top of each image (\kbd{B}).
\end{description}
The keys apply while an image has keyboard focus.

\section{Loading the viewpoints}
\label{sec:gcs-event-card}

The \gui{Event \& channels (UTC)} card (Figure~\ref{fig:gcs-event-card}) defines what is
loaded.

\screenshot[0.5\linewidth]{gcs/gcs_event_card}{The \gui{Event \& channels} card.}{fig:gcs-event-card}{The Event \& channels (UTC) card with the event range, the frame cap, Load all, Max time offset, and the three channel rows A, B and C with their source, frame count, Load button and time range.}

\begin{description}
  \item[Event] the UTC start and end of the event, up to three days. Changing them sets every
    channel's range; each channel's range can then be adjusted on its own.
  \item[\(\leq\) frames] the maximum number of images per channel (2--200; default 60). A longer
    sequence is thinned evenly across the event rather than cut short.
  \item[Load all] (\kbd{Ctrl+L}) loads every channel. Each channel also has its own \gui{Load}
    button and a frame count. \menu{Event > Cancel downloads} stops a slow load.
  \item[Max time offset] the largest difference between an image's time and the shared time
    for the image to take part in a fit (0--60~min; default 5~min). Models are still drawn on
    the other panels for comparison. Use a smaller value for a fast-evolving CME; the default is
    a workflow choice, not an accuracy guarantee.
  \item[A, B, C] the channel of each panel and its time range.
\end{description}

The available channels are listed in Table~\ref{tab:gcs-channels}. By default panels~A, B and
C show STEREO-B COR2, SOHO LASCO C2 and STEREO-A\index{STEREO} COR2, left to right. Channels without data
for the event date are grayed out. The STEREO-B archive ends on 27~September~2014, so for later
events panel~A reports that it has no data; only SOHO and STEREO-A then remain as vantage
points, and another channel in panel~A (LASCO~C3 or COR1-A) adds an image but not a new viewing
direction.

\begin{table}[!htbp]
  \centering
  \caption{Coronagraph channels available for GCS fitting.}
  \label{tab:gcs-channels}
  \small
  \begin{tabularx}{\linewidth}{@{}P{0.3\linewidth}L@{}}
    \toprule
    \thead{Channel} & \thead{Available} \\
    \midrule
    SOHO LASCO C2 & from 21~February 1997 \\
    SOHO LASCO C3 & from 20~February 1997 \\
    STEREO-A COR1 & from 4~December 2006 \\
    STEREO-A COR2 & from 29~December 2006 \\
    STEREO-B COR1 & 12~December 2006 to 27~September 2014 \\
    STEREO-B COR2 & 10~January 2007 to 27~September 2014 \\
    \bottomrule
  \end{tabularx}
\end{table}

The images are Helioviewer\index{Helioviewer} JPEG2000 files, so loading needs network access; downloaded images
are cached and reused. Changing a channel's source or range discards its loaded images and any
pending results. A 15-minute e-CALLISTO file spans only one or two coronagraph images (about
one every 15~minutes for COR2 and every 12~minutes for LASCO~C2), so widen the range to follow a
CME.

\section{Reading the images}

The banner across each image gives the instrument, the actual observation time, its offset
from the shared time and the difference reference. In running difference\index{running difference} the first image of
each channel is differenced against the archive image just before the range, and is shown raw
only when no earlier image exists. Clicked front points belong to the image on which they were
clicked and return when that image is shown again.

\section{The Display card}
\label{sec:gcs-display-card}

\begin{description}
  \item[Wireframe] \gui{GCS…} and \gui{Shock…} choose the colors of the two models (orange
    and sky blue by default); \gui{Width} (0.5--8~px), \gui{Opacity} (10--100\%) and
    \gui{Density} (mesh density, 1--8) set the wireframe style.
  \item[Image A, B, C] select a panel, then set its \gui{Colormap} and contrast sliders. Every
    load resets the color map to gray.
\end{description}

\section{Menus}
\label{sec:gcs-menus}

\begin{xltabular}{\linewidth}{@{}P{0.42\linewidth} L P{0.15\linewidth}@{}}
  \caption{Menus of the GCS CME Fitting window.}\label{tab:gcs-menus}\\
  \toprule \thead{Command} & \thead{Function} & \thead{Shortcut} \\ \midrule\endfirsthead
  \toprule \thead{Command} & \thead{Function} & \thead{Shortcut} \\ \midrule\endhead
  \menu{File > Export analysis JSON…}\index{GCS!JSON export} & Parameters, points and provenance of both models. & \kbd{Ctrl+Shift+S} \\
  \menu{File > Export recorded fits CSV…} & Recorded series of the edited model. & \\
  \menu{File > Export height–time graph…} & Height--time graph with the chosen fit. & \\
  \menu{File > Save viewpoint snapshot…}\index{GCS!viewpoint snapshot} & The three views as shown. & \\
  \menu{File > Export movie (GIF/MP4)…} & Movie of the event with recorded shells\index{GCS!recorded shells}. & \\
  \menu{File > Export fitting report (PDF)…}\index{GCS!fitting report} & Complete fitting report. & \\
  \menu{File > Close window} & Close the window. & \kbd{Ctrl+W} \\
  \menu{Event > Load all channels}\index{GCS!loading} & Load every channel. & \kbd{Ctrl+L} \\
  \menu{Event > Cancel downloads} & Stop running loads. & \\
  \menu{Event > Apply event range to all channels} & Reset every channel to the event range. & \\
  \menu{Event > Go to UTC time…} & Jump to the image nearest a typed time. & \kbd{Ctrl+G} \\
  \menu{Event > First / Previous / Next / Last frame}, \menu{Play / pause} & Navigation and playback. & \kbd{Space} \\
  \menu{View > Raw / Running difference / Base difference} & Image mode. & \\
  \menu{View > Equal viewpoints / Focus A / Focus B / Focus C} & Layout. & \kbd{0}--\kbd{3} \\
  \menu{View > GCS wireframe / Shock wireframe / Solar limb / Arcsecond axes / Image banners} & Overlay toggles. & \kbd{B} \\
  \menu{View > Reset image zoom}, \menu{Fit layout to window} & Restore the zoom and the panel sizes. & \\
  \menu{Fit > Edit GCS flux rope / Edit shock} & Choose the model to edit. & \\
  \menu{Fit > Pick front points}\index{GCS!front points} & Toggle point picking. & \\
  \menu{Fit > Undo last front point} & Remove the last point clicked in any panel. & \kbd{Ctrl+Z} \\
  \menu{Fit > Clear current front points} & Remove the edited model's points. & \\
  \menu{Fit > Refine current model} & Least-squares refinement. & \kbd{Ctrl+R} \\
  \menu{Fit > Record fit at current time} & Commit the fit. & \kbd{Ctrl+Return} \\
  \menu{Fit > Restore / Delete recorded fit at current time} & Revisit or remove a recorded fit. & \\
  \menu{Fit > Reset model parameters} & Reset the edited model. & \\
  \menu{Fit > Fit recorded heights over time} & Kinematic fit. & \\
  \menu{Help > Fitting workflow and scientific limits…} & Summary of the procedure and its limits. & \kbd{F1} \\
  \bottomrule
\end{xltabular}

%% Chapter 21 ----------------------------------------------------------------
\chapter{Fitting the Flux Rope and the Shock}
\label{ch:gcs-fitting}

This chapter describes the models of the GCS CME Fitting window, how they are aligned and
refined against the observed fronts, how fits are recorded over time to build kinematics, and
how results are exported.

\section{The Model card}

The \gui{Model} card (Figure~\ref{fig:gcs-model-card}) shows the parameters of one model at a
time. Choose the model with \gui{Edit} on the card's title row:
\begin{description}
  \item[GCS flux rope] the Graduated Cylindrical Shell, fitted to the CME ejecta.
  \item[Shock] a spheroid or ellipsoid, fitted to the shock the CME drives --- its faint outer
    envelope (Section~\ref{sec:shock-fitting}).
\end{description}
Both models are drawn in every panel (the shell in orange and the shock in sky blue by
default); the edited model is the one that the handles, clicks, \gui{Refine fit}\index{GCS!refine fit},
\gui{Commit} and the Kinematics card act on.

\screenshot[0.45\linewidth]{gcs/gcs_model_card}{The \gui{Model} card with the GCS parameters.}{fig:gcs-model-card}{The Model card in GCS flux rope mode: the six sliders with numeric entry (Lon, Lat, Tilt, Height, α, κ), the derived-geometry readout, Refine fit and Commit GCS, and the Pick front points, wireframe, Undo point and Clear points controls.}

\section{GCS parameters}
\label{sec:gcs-parameters}
\index{GCS!parameters}

Six sliders, each with a numeric entry, control the shell (Table~\ref{tab:gcs-params}).

\begin{table}[!htbp]
  \centering
  \caption{Parameters of the GCS flux-rope model.}
  \label{tab:gcs-params}
  \small
  \begin{tabularx}{\linewidth}{@{}P{0.12\linewidth}P{0.2\linewidth}L@{}}
    \toprule
    \thead{Slider} & \thead{Range} & \thead{Meaning} \\
    \midrule
    Lon & \(-180^\circ\) to \(180^\circ\) & Stonyhurst longitude of the direction of propagation. \\
    Lat & \(-89^\circ\) to \(89^\circ\) & Stonyhurst latitude of the direction of propagation. \\
    Tilt & \(-90^\circ\) to \(90^\circ\) & Rotation of the shell about the propagation direction. \\
    Height & 1.05--30~\(R_\odot\) & Distance of the leading edge (apex) from Sun center --- not the altitude above the surface. \\
    \(\alpha\) & \(1^\circ\)--\(80^\circ\) & Half-angle between the legs of the shell. \\
    \(\kappa\) & 0.05--0.95 & Aspect ratio, controlling the thickness of the shell. \\
    \bottomrule
  \end{tabularx}
\end{table}

Set the direction, tilt, \(\alpha\) and \(\kappa\) before fine-tuning the height. The apex can
also be dragged in any panel: moving it around the Sun turns the direction, and moving it
outwards or inwards changes the height. The readout below the sliders gives the derived leg
height, apex cross-section radius, center distance, and the intrinsic face-on and edge-on
widths.

\begin{background}[Background: GCS geometry]
  The GCS is a hollow croissant: two conical legs joined by a curved front with a circular
  cross-section that grows with distance from the Sun \parencite{thernisien2006,thernisien2011}.
  The face-on width is \(2[\alpha + \arcsin\kappa]\) and the edge-on width \(2\arcsin\kappa\);
  both include the shell thickness and differ from the apparent width in a coronagraph image.
  Because the height is the leading edge, it feeds the height--time fit directly. GCS is a
  flux-rope model, not a shock model.
\end{background}

\section{Clicking the front}
\label{sec:front-points}

With \gui{Pick front points}\index{GCS!front points} ticked, left-click points along the same front in two or more
views. Points are drawn in each panel and are used by \gui{Refine fit}:
\begin{itemize}
  \item a right-click removes that panel's last point;
  \item \gui{Undo point} (\kbd{Ctrl+Z}) removes the point clicked last in any panel;
  \item \gui{Clear points} removes all of the edited model's points on the displayed images;
  \item untick \gui{Pick front points} to navigate without adding points.
\end{itemize}
The \gui{GCS wireframe} and \gui{Shock wireframe} options hide a model temporarily so that the
observed front can be inspected. Each model keeps its own front points.

\section{Refining the fit}
\label{sec:refine}
\index{GCS!refinement}

\gui{Refine fit} (\kbd{Ctrl+R}) makes a local least-squares adjustment of the edited model
through the clicked points. It refines a fit; it cannot find one from scratch, because
longitude, \(\alpha\) and \(\kappa\) are poorly conditioned in GCS, so first align the shell
roughly with the CME by eye. The refinement fits as many parameters as the points support
(Table~\ref{tab:refine-stages}); the status bar names the parameters fitted and how many more
points would free the next one.

\begin{table}[!htbp]
  \centering
  \caption{Staged refinement of the GCS model.}
  \label{tab:refine-stages}
  \small
  \begin{tabularx}{\linewidth}{@{}P{0.38\linewidth}L@{}}
    \toprule
    \thead{Front points} & \thead{Parameters fitted} \\
    \midrule
    2 & Height \\
    4, in two separated views & Height, longitude, latitude \\
    5 & \ldots and tilt \\
    6 & \ldots and \(\alpha\) \\
    7 & All six parameters (adding \(\kappa\)) \\
    \bottomrule
  \end{tabularx}
\end{table}

The status bar also reports the views taking part, their angular separations and the number
of points. Only images within the \gui{Max time offset}\index{GCS!max time offset} of the shared time take part. With a
single view, or with no pair of views between 20° and 160° apart, the direction cannot be
constrained and is held fixed; empty, nearly coincident and nearly opposite viewpoints do not
constrain the direction independently. Formal errors are shown when the local solution is
reliable, are withheld otherwise, and become invalid when the model or the observations
change.

\section{Recording fits}
\label{sec:gcs-recording}

\gui{Commit GCS}\index{GCS!commit} (\kbd{Ctrl+Return}) records the current fit at the shared time, together with
the images it was made from and their observation details. Step to later times and repeat to
build a height--time series. One combination of images cannot be recorded at two different
times, so repeated images never count as independent height measurements.

When stepping, dragging the time slider or playing, the recorded shells\index{GCS!recorded shells} are drawn: each fit on
its own images, interpolated between recorded times (a display, not a fit) and held before the
first and after the last record. A model with nothing recorded keeps its slider values, so
commit before stepping away. \menu{Fit > Restore recorded model at current time} and
\menu{Delete recorded fit at current time} revisit or remove a record; \menu{Fit > Reset model
parameters} resets the edited model.

Moving the event range to a different event sets the recorded fits aside. They return when you
come back to that event, and the JSON export includes them.

\section{Kinematics}
\label{sec:gcs-kinematics}
\index{GCS!kinematics}

The \gui{Kinematics} card (Figure~\ref{fig:gcs-kinematics}) lists the recorded fits of the
edited model beside a height--time graph. For the GCS model the table shows \gui{Time (UT)},
\gui{t (s)}, \gui{Apex (R☉)}, \gui{Lon (°)}, \gui{Lat (°)} and \gui{α (°)}; for the shock,
\(\kappa\) replaces \(\alpha\).

\screenshot{gcs/gcs_kinematics}{A GCS fit with three recorded times; the \gui{Kinematics} card (right) lists the recorded fits beside their height--time graph and linear fit.}{fig:gcs-kinematics}{The Kinematics · GCS card with several recorded fits in the table, the height–time graph with a quadratic fit and error bars, and the Fit, Fit height–time, Clear fits and Export controls.}

\begin{itemize}
  \item Double-click a row to go back to its time and restore its model.
  \item Choose a \gui{Linear}\index{height--time fit!order}, \gui{Quadratic} or \gui{Cubic} fit and click
    \gui{Fit height–time}\index{GCS!kinematics}. A fit of degree \(n\) needs \(n+1\) recorded times, and \(n+2\)
    before errors can be estimated (Section~\ref{sec:tracking-panel}).
  \item \gui{Clear fits} removes the recorded fits of the edited model.
  \item \gui{Export} saves the table as CSV or the graph as a figure.
\end{itemize}
The heights are de-projected under the model assumptions, so the speeds are model-dependent
radial speeds rather than plane-of-sky speeds.

\section{Fitting the shock}
\label{sec:shock-fitting}
\index{shock model}\index{spheroid}\index{ellipsoid}

Choose \gui{Edit: Shock} to fit the shock with a \gui{Spheroid} or an \gui{Ellipsoid}
(Figure~\ref{fig:gcs-shock}). The parameters follow PyThea\index{PyThea}'s conventions exactly
(Table~\ref{tab:shock-model-params}).

\begin{table}[!htbp]
  \centering
  \caption{Parameters of the shock model.}
  \label{tab:shock-model-params}
  \small
  \begin{tabularx}{\linewidth}{@{}P{0.13\linewidth}P{0.2\linewidth}L@{}}
    \toprule
    \thead{Slider} & \thead{Range} & \thead{Meaning} \\
    \midrule
    Lon, Lat & as for GCS & Stonyhurst direction of the shock center and apex. \\
    Height & 1.05--30~\(R_\odot\) & Heliocentric distance of the apex (center distance plus the radial semi-axis \(a\)). \\
    \(\kappa\) & 0.05--2.0 & Self-similar constant \(\kappa = b/(h - 1\,R_\odot)\): the lateral size for a given apex height. \\
    \(\varepsilon\) & \(-0.99\) to \(0.99\) & Signed eccentricity: positive stretches the shock radially (\(a > b\)), negative flattens it (\(a < b\)). \\
    Tilt & \(-90^\circ\) to \(90^\circ\) & Ellipsoid only: rotation about the radial axis; at 0 the \(b\) axis is parallel to the solar equator. \\
    \(\alpha\) (b/c) & 0.5--1.5 & Ellipsoid only: ratio of the two lateral semi-axes. \\
    \bottomrule
  \end{tabularx}
\end{table}

\screenshot{gcs/gcs_shock_fit}{Shock and flux-rope models in three viewpoints.}{fig:gcs-shock}{The three viewpoints with both the GCS shell (orange) and the shock spheroid (sky blue) drawn, Edit: Shock selected, and shock front points clicked on the outer envelope.}

The derived readout gives the center distance and the semi-axes, under PyThea's names. Click the
shock front and use \gui{Refine fit}: the model's exactly computed projected outline is pulled
through the points in every view, with its own sequence of height, direction and shape
parameters. Record shock fits with \gui{Commit Shock}\index{shock model!commit}. The shock keeps its own front points,
recorded fits, kinematics and CSV file. With fewer than three views, \(\varepsilon\) and the tilt
are weakly constrained and an ellipsoid's shape parameters trade off against each other, so
prefer a spheroid.

\section{Exporting results}
\label{sec:gcs-export}
\index{GCS!export}

Images and graphs are drawn again from the data rather than copied from the screen, always on
a white page whatever the application theme; graphs follow the OriginPro style. Figures can be
saved as PNG, PDF, EPS, SVG, TIFF or JPG.
\begin{description}
  \item[Export analysis JSON…] (\kbd{Ctrl+Shift+S}) the current, recorded and set-aside
    parameters of both models, with the clicked points, image times, offsets and observer
    geometry. The file (schema version~2) has a \texttt{shock} section beside the GCS keys. It
    is an analysis export, not a PyThea session file.
  \item[Export recorded fits CSV…] the recorded series of the edited model
    (\file{cme_gcs_fit.csv} or \file{cme_shock_fit.csv}).
  \item[Export height–time graph…] the recorded apex heights with error bars and the chosen
    fit, titled with the model and fit, with the speed and acceleration in the legend.
  \item[Save viewpoint snapshot…] the three views as shown, with arcsecond axes, the
    wireframes, the solar limb, the front points and a note of where each drawn shell comes
    from. With the banners hidden, only each panel's name is printed above it.
  \item[Export movie (GIF/MP4)…] every time step as playback shows it, recorded shells
    included, at the playback speed. MP4 needs the bundled FFmpeg; a GIF is offered when it is
    unavailable.
  \item[Export fitting report (PDF)…] the viewpoints and their separations, the model at the
    current time, every recorded fit drawn on the images it was made from with its parameters
    and observation details, each model's linear, quadratic and cubic height--time fits (each
    with a titled graph and all its parameters, then compared), and the method, its limits and
    references (Figure~\ref{fig:gcs-report}).
\end{description}
Stepping through time for a movie or a report leaves the window as it was: the time shown, the
working models, their formal errors and the front points are restored afterwards.

\screenshot[0.55\linewidth]{gcs/gcs_report_page}{A page of the GCS fitting report.}{fig:gcs-report}{A representative page of the exported PDF fitting report showing a recorded fit drawn on its three images with the parameter table.}

\begin{caution}[Scientific limits]
  GCS models the flux rope, not the shock. Formal fit errors leave out the uncertainty from
  identifying the front, non-simultaneous images, image preparation and the assumed geometry,
  and a small residual does not make a fit unique or accurate \parencite{verbeke2023}. Helioviewer\index{Helioviewer}
  JPEG2000 images are display products: suitable for fitting shapes, not for calibrated
  intensity measurements. \menu{Help > Fitting workflow and scientific limits…} (\kbd{F1})
  summarizes the procedure and these limits.
\end{caution}


% ===========================================================================
\appendix
% ===========================================================================
\addpart{Appendices}
%% Appendix A ----------------------------------------------------------------
\chapter{Keyboard Shortcuts}
\label{app:shortcuts}
\index{keyboard shortcuts}

On macOS, \kbd{Ctrl} in the shortcuts below corresponds to the Command key.

\begin{table}[!htbp]
  \centering
  \caption{Shortcuts of the main window.}
  \label{tab:keys-main}
  \small
  \begin{tabularx}{\linewidth}{@{}P{0.3\linewidth}L@{}}
    \toprule
    \thead{Shortcut} & \thead{Action} \\
    \midrule
    \kbd{Ctrl+O} & Open FITS files \\
    \kbd{Ctrl+Shift+O} & Open a project \\
    \kbd{Ctrl+S} & Save the project \\
    \kbd{Ctrl+Shift+S} & Save the project as a new file \\
    \kbd{Ctrl+E} & Export the figure \\
    \kbd{Ctrl+F} & Export the spectrum as FITS \\
    \kbd{Ctrl+Z} & Undo \\
    \kbd{Ctrl+Shift+Z} & Redo \\
    \kbd{Ctrl+Shift+Left} & Add the previous observation to the timeline \\
    \kbd{Ctrl+Shift+Right} & Add the next observation to the timeline \\
    \kbd{F1} & Open the built-in user guide \\
    Right-click or double-click & Finish drift-rate picking; close a polygon annotation \\
    \bottomrule
  \end{tabularx}
\end{table}

\begin{table}[!htbp]
  \centering
  \caption{Shortcuts of Solar Image Analysis.}
  \label{tab:keys-solar}
  \small
  \begin{tabularx}{\linewidth}{@{}P{0.3\linewidth}L@{}}
    \toprule
    \thead{Shortcut} & \thead{Action} \\
    \midrule
    \kbd{Ctrl+S} & Save the session \\
    \kbd{Ctrl+Shift+S} & Save the session as a new file \\
    \kbd{Ctrl+Return} & Commit the fitted circle (\gui{Circle Fit}) \\
    \kbd{Esc} or right-click & Cancel the measurement pick in progress \\
    \bottomrule
  \end{tabularx}
\end{table}

\begin{table}[!htbp]
  \centering
  \caption{Shortcuts of the GCS CME Fitting window. The single-key shortcuts apply while an image has focus.}
  \label{tab:keys-gcs}
  \small
  \begin{tabularx}{\linewidth}{@{}P{0.3\linewidth}L@{}}
    \toprule
    \thead{Shortcut} & \thead{Action} \\
    \midrule
    \kbd{Space} & Play or pause \\
    \kbd{←} / \kbd{→} & Previous or next frame \\
    \kbd{Home} & First frame \\
    \kbd{B} & Show or hide the image banners \\
    \kbd{0} & Equal layout \\
    \kbd{1}, \kbd{2}, \kbd{3} & Enlarge panel A, B or C \\
    \kbd{Ctrl+L} & Load all channels \\
    \kbd{Ctrl+G} & Go to a UTC time \\
    \kbd{Ctrl+Z} & Undo the last front point \\
    \kbd{Ctrl+R} & Refine the edited model \\
    \kbd{Ctrl+Return} & Record the fit at the current time \\
    \kbd{Ctrl+Shift+S} & Export the analysis JSON \\
    \kbd{Ctrl+W} & Close the window \\
    \kbd{F1} & Fitting workflow and scientific limits \\
    Right-click & Remove the panel's last front point \\
    \bottomrule
  \end{tabularx}
\end{table}

%% Appendix B ----------------------------------------------------------------
\chapter{File Formats and Outputs}
\label{app:file-formats}
\index{file formats}

\section{Files read and written by the application}

\begin{table}[!htbp]
  \centering
  \caption{File types of the \appname{}.}
  \label{tab:file-types}
  \small
  \begin{tabularx}{\linewidth}{@{}P{0.27\linewidth}L@{}}
    \toprule
    \thead{Extension} & \thead{Meaning} \\
    \midrule
    \file{.fit}, \file{.fits}, \file{.fit.gz}, \file{.fits.gz} & Radio spectra in the main window, and solar images in Solar Image Analysis. Upper-case suffixes are accepted. \\
    \file{.efaproj} & Complete analyzer project: data, processing, view and analysis session (Section~\ref{sec:projects}). \\
    \file{.efaview.json} & Portable view configuration: display range, units, thresholds, color map and graph styling (Section~\ref{sec:view-config}). \\
    \file{.ecsolar} & Solar Image Analysis session with embedded FITS frames (Section~\ref{sec:solar-sessions}). \\
    \file{.srs} & Learmonth spectrograph daily archive file (Section~\ref{sec:learmonth}). \\
    \bottomrule
  \end{tabularx}
\end{table}

\section{Exported FITS files}
\label{sec:fits-structure}

A FITS file written with \menu{File > Export As > Export to FIT} has the structure of a
standard CALLISTO file:
\begin{description}
  \item[Primary HDU] the two-dimensional spectrum (frequency~\(\times\)~time) with the chosen
    \texttt{BITPIX}, and the original header updated with the keywords in
    Table~\ref{tab:fits-keywords}.
  \item[\texttt{AXIS} extension] a binary table with one row and the columns
    \texttt{FREQUENCY} (MHz) and \texttt{TIME} (s).
\end{description}

\begin{table}[!htbp]
  \centering
  \caption{Header keywords added or updated on FITS export.}
  \label{tab:fits-keywords}
  \small
  \begin{tabularx}{\linewidth}{@{}P{0.2\linewidth}L@{}}
    \toprule
    \thead{Keyword} & \thead{Content} \\
    \midrule
    \texttt{HISTORY} & ``Exported by e-CALLISTO FITS Analyzer'', the plot type, the units shown, and the source file (or the first and last source of a combined dataset). \\
    \texttt{COMBINED} & \texttt{T} for a combined dataset, \texttt{F} otherwise. \\
    \texttt{COMBMETH} & Combination method: \texttt{time}, \texttt{frequency} or \texttt{time\_frequency}. \\
    \texttt{NFILES} & Number of source files of a combined dataset. \\
    \texttt{BUNIT} & Intensity unit of the data. \\
    \texttt{DATAMIN}, \texttt{DATAMAX} & Minimum and maximum data values. \\
    \texttt{BSCALE}, \texttt{BZERO} & Scaling for integer data. \\
    \bottomrule
  \end{tabularx}
\end{table}

\section{Summary of outputs}

\begin{xltabular}{\linewidth}{@{}P{0.33\linewidth} P{0.24\linewidth} L@{}}
  \caption{Outputs produced by the application.}\label{tab:outputs}\\
  \toprule \thead{Output} & \thead{Format} & \thead{Produced by} \\ \midrule\endfirsthead
  \toprule \thead{Output} & \thead{Format} & \thead{Produced by} \\ \midrule\endhead
  Spectrum figure & PNG, PDF, EPS, SVG, TIFF, JPG & \menu{File > Export As > Export Figure} \\
  Processed spectrum & FITS & \menu{File > Export As > Export to FIT} \\
  Provenance report & Markdown and JSON & \menu{File > Export As > Export Provenance Report...} \\
  Analysis log & CSV and TXT & \menu{File > Export As > Export Analysis Log...} \\
  Project report & PDF & \menu{File > Generate Project Report...} \\
  Batch images & PNG (300~dpi) & Batch processor \\
  Maximum-intensity points & CSV (time channel, time, frequency) & Maximum Intensities, \menu{File > Save As} \\
  Analyzer results & Excel \file{.xlsx} (Table~\ref{tab:excel-columns}) & Burst Analyzer, \gui{Save Data} \\
  Analysis graphs & PNG, PDF, EPS, SVG, TIFF, JPG & Maximum Intensities, Analyzer, band-splitting \\
  Spectral overview & PNG, PDF, EPS, SVG, TIFF & Downloader, \gui{Spectral Overview} \\
  Station comparison & PNG, PDF, EPS, SVG, TIFF & Multi-Station Comparison \\
  Learmonth conversion & FIT & Learmonth browser \\
  GOES, SEP, Dst, Kp data & PNG and CSV & Solar-event viewers \\
  SunPy plots and statistics & PNG, PDF, SVG, TIFF, JPG; CSV & SunPy Explorer \\
  Solar images & PNG; cropped FITS; MP4, GIF & Solar Image Analysis \\
  Active regions & CSV & Solar Image Analysis \\
  CME tracking & CSV; PNG, PDF, EPS, SVG, TIFF, JPG & Solar Image Analysis tracking panel \\
  PFSS diagnostics & PNG, PDF, SVG & PFSS diagnostics window \\
  GCS analysis & JSON; CSV; figures; GIF, MP4; PDF report & GCS CME Fitting \\
  Diagnostics bundle & ZIP & \menu{About > Report a Bug...} \\
  \bottomrule
\end{xltabular}

%% Appendix C ----------------------------------------------------------------
\chapter{Data Sources}
\label{app:data-sources}
\index{data sources}

Table~\ref{tab:data-sources} lists the archives from which the application retrieves data.
Access requires a network connection; most archives are public and need no registration, the
exception being JSOC exports, which need a registered e-mail address. When you publish results
based on these data, acknowledge the providers as their data policies request.

{\small
\begin{xltabular}{\linewidth}{@{}P{0.25\linewidth} P{0.33\linewidth} L@{}}
  \caption{Archives used by the \appname{}.}\label{tab:data-sources}\\
  \toprule \thead{Data} & \thead{Provider and address} & \thead{Used by} \\ \midrule\endfirsthead
  \toprule \thead{Data} & \thead{Provider and address} & \thead{Used by} \\ \midrule\endhead
  e-CALLISTO FITS spectra & e-CALLISTO archive, FHNW, \url{https://soleil.i4ds.ch/solarradio/data/2002-20yy_Callisto/} & Downloader, Timeline, Multi-Station Comparison \\
  e-CALLISTO burst lists & Monstein catalog, \url{https://soleil.i4ds.ch/solarradio/data/BurstLists/2010-yyyy_Monstein/} & Downloader, \gui{Burst List} \\
  Learmonth spectrograph & Australian Bureau of Meteorology, Space Weather Services, \url{https://downloads.sws.bom.gov.au/wdc/wdc_spec/data/learmonth/raw/} & Learmonth browser \\
  STEREO/SWAVES level-2 & NASA Space Physics Data Facility, \url{https://spdf.gsfc.nasa.gov/pub/data/stereo/combined/swaves/level2_cdf/} & SWAVES loader \\
  GOES XRS and SGPS & NOAA National Centers for Environmental Information, \url{https://data.ngdc.noaa.gov/platforms/solar-space-observing-satellites/} & GOES viewers, GOES overlay \\
  Dst index & World Data Center for Geomagnetism, Kyoto, \url{https://wdc.kugi.kyoto-u.ac.jp} & Dst viewer \\
  Kp index & GFZ German Research Centre for Geosciences, \url{https://kp.gfz.de} & Kp viewer \\
  SOHO/LASCO CME catalog & CDAW Data Center, NASA GSFC, \url{https://cdaw.gsfc.nasa.gov} & CME catalog \\
  SDO/AIA and SDO/HMI & Joint Science Operations Center, \url{https://jsoc.stanford.edu}; Virtual Solar Observatory & Solar Image Analysis, SunPy Explorer, PFSS (HMI synoptic) \\
  SOHO/EIT, SOHO/LASCO, STEREO/SECCHI, GOES/SUVI, PROBA2/SWAP & Virtual Solar Observatory and mission archives, through SunPy & Solar Image Analysis, SunPy Explorer, overlay layers \\
  Helioviewer JPEG2000 images & Helioviewer Project, \url{https://api.helioviewer.org} & Live Preview, GCS CME Fitting \\
  GONG synoptic and ADAPT maps & National Solar Observatory, through SunPy & PFSS \\
  NOAA active-region labels & Heliophysics Event Knowledgebase (HEK), through SunPy & Active-region labels \\
  Software updates & GitHub releases, \url{https://github.com/SaanDev/e-Callisto_FITS_Analyzer} & Update check \\
  \bottomrule
\end{xltabular}}

\begin{note}
  Archive availability changes over time. The calibrated SOHO/LASCO and SOHO/EIT archives lag
  real time by many months, the HMI synoptic series is published about one Carrington rotation
  behind, and STEREO-B data end in 2014. Downloaded files are cached locally and can be
  reused without network access.
\end{note}

%% Appendix D ----------------------------------------------------------------
\chapter{Models and Formulas}
\label{app:formulas}
\index{formulas}

This appendix collects the relations used to compute the quantities reported by the
application. Heights \(r\) are heliocentric distances in solar radii (\(R_\odot = 696\,340\)~km),
electron densities \(n_e\) are in cm\(^{-3}\) and frequencies \(f\) in MHz.

\section{Intensity scales}

\begin{description}
  \item[CALLISTO digits to dB] the CALLISTO receiver's logarithmic detector maps 2500~mV onto
    256 digits with a slope of 25.4~mV/dB, so that
    \(\Delta I_\mathrm{dB} = \Delta I_\mathrm{digits}\times 2500/(256\times25.4)\approx
    0.384\)~dB per digit. The Plotutil median-dB recipe used by \gui{Median (dB)} and the
    batch processor uses \(2500/(255\times25.4)\).
  \item[ARTEMIS-IV] \(4096\) ADC counts over a 70~dB range: 58.51 counts per dB.
  \item[GOES flare class] from the peak 0.1--0.8~nm flux \(F\): A for \(F<10^{-7}\), B for
    \(10^{-7}\le F<10^{-6}\), C for \(10^{-6}\le F<10^{-5}\), M for \(10^{-5}\le F<10^{-4}\) and
    X for \(F\ge10^{-4}\)~W\,m\(^{-2}\); the number is \(F\) divided by the lower limit of the class (\(10^{-8}\) for A up to \(10^{-4}\) for X), for example M2.5 for \(F = 2.5\times10^{-5}\)~W\,m\(^{-2}\).
\end{description}

\section{Radio burst analysis}

\paragraph{Plasma frequency.} The fundamental emission frequency is the local plasma frequency
\begin{equation}
  f_p = 8.98\times10^{-3}\sqrt{n_e}\quad\text{MHz},
\end{equation}
so \(n_e = (f/8.98\times10^{-3})^2\). For harmonic emission \(f = 2f_p\); the application divides
the observed frequencies and drift rates by two before computing heights and speeds.

\paragraph{Coronal density models.} Each model is multiplied by the fold number \(N\)
(\(n_e \to N\,n_e\)).

\begin{table}[!htbp]
  \centering
  \caption{Coronal electron-density models.}
  \label{tab:density-formulas}
  \small
  \begin{tabularx}{\linewidth}{@{}P{0.22\linewidth}L@{}}
    \toprule
    \thead{Model} & \thead{\(n_e(r)\) in cm\(^{-3}\)} \\
    \midrule
    Newkirk & \(4.2\times10^{4}\times10^{4.32/r}\) \\
    Saito & \(1.36\times10^{6}\,r^{-2.14} + 1.68\times10^{8}\,r^{-6.13}\) \\
    Leblanc & \(3.3\times10^{5}\,r^{-2} + 4.1\times10^{6}\,r^{-4} + 8.0\times10^{7}\,r^{-6}\) \\
    Baumbach--Allen & \(10^{8}\left(2.99\,r^{-16} + 1.55\,r^{-6} + 0.036\,r^{-1.5}\right)\) \\
    Mann & \(5.14\times10^{9}\exp\!\left[13.83\left(1/r - 1\right)\right]\) \\
    \bottomrule
  \end{tabularx}
\end{table}

For the Newkirk model the emission height has the closed form
\begin{equation}
  r = \frac{4.32\,\ln 10}{\ln\!\left[n_e/(N\times4.2\times10^{4})\right]}.
\end{equation}
For the other models the height is found numerically between \(1\,R_\odot\) and 215~\(R_\odot\)
(1~AU); frequencies outside this range have no height.

\paragraph{Power-law backbone and drift rate.}
\begin{equation}
  f(t) = a\,(t-t_0)^{-b},\qquad
  \frac{\mathrm{d}f}{\mathrm{d}t} = -a\,b\,(t-t_0)^{-b-1}.
\end{equation}
The goodness of fit is reported as \(R^2 = 1 - \sum(f_i-\hat f_i)^2/\sum(f_i-\bar f)^2\) and
\(\mathrm{RMSE} = \sqrt{\overline{(f_i-\hat f_i)^2}}\).

\paragraph{Shock speed.}
\begin{equation}
  v = R_\odot\left|\frac{\mathrm{d}f}{\mathrm{d}t}\right|\left|\frac{\mathrm{d}r}{\mathrm{d}f}\right|,
  \qquad \frac{\mathrm{d}r}{\mathrm{d}f} = \frac{2\,n_e}{f\,\left|\mathrm{d}n_e/\mathrm{d}r\right|}.
\end{equation}
The \emph{initial} values are evaluated at the starting frequency, the 90th percentile of the
fitted frequencies; the \emph{average} values are means along the backbone with their standard
errors.

\paragraph{Band splitting} \parencite{vrsnak2002}.
\begin{gather}
  X = \left(\frac{f_u}{f_l}\right)^2,\qquad
  M_A = \sqrt{\frac{X\,(X+5)}{2\,(4-X)}},\qquad
  V_A = \frac{V_s}{M_A},\\
  B = V_A\sqrt{\mu_0\,m_p\,n_e}\quad(\text{SI units; reported in gauss}).
\end{gather}

\section{Image measurements}

\paragraph{Circle fit.} For a sphere of fitted radius \(r\) resting on the solar surface, the
height of its top is \(h = 1\,R_\odot + 2r\).

\paragraph{Height--time fits.}
\begin{align}
  \text{linear:}\quad & h = h_0 + v\,t,\\
  \text{quadratic:}\quad & h = h_0 + v_0\,t + \tfrac12 a\,t^2,\\
  \text{cubic:}\quad & h = h_0 + v_0\,t + \tfrac12 a_0\,t^2 + \tfrac16 j\,t^3.
\end{align}
Uncertainties are propagated from the covariance of the fitted coefficients.

\section{Three-dimensional CME models}

\paragraph{GCS} \parencite{thernisien2006,thernisien2011}. Parameters: Stonyhurst longitude and
latitude of the propagation direction, tilt, leading-edge height \(h\), half-angle \(\alpha\) and
aspect ratio \(\kappa\). Intrinsic widths:
\begin{equation}
  w_\text{face-on} = 2\left[\alpha + \arcsin\kappa\right],\qquad
  w_\text{edge-on} = 2\arcsin\kappa.
\end{equation}

\paragraph{Shock spheroid and ellipsoid} \parencite{kouloumvakos2022}. With the radial semi-axis
\(a\) and lateral semi-axes \(b\) and \(c\):
\begin{equation}
  h = r_\text{center} + a,\qquad \kappa = \frac{b}{h - 1\,R_\odot},\qquad
  \varepsilon = \begin{cases} +\sqrt{1-(b/a)^2} & a > b,\\ -\sqrt{1-(a/b)^2} & a < b,\end{cases}
  \qquad \alpha = \frac{b}{c}.
\end{equation}
A spheroid has \(c = b\) (\(\alpha = 1\), no tilt).

%% Appendix E ----------------------------------------------------------------
\chapter{Troubleshooting}
\label{app:troubleshooting}
\index{troubleshooting}

\newcommand{\problem}[1]{\par\medskip\noindent{\sffamily\bfseries\color{accentdark}#1}\par\nobreak\smallskip}

\section{Main window}

\problem{Buttons and menu items are grayed out.}
Most analysis controls, and the \gui{Graph Properties} panel, are enabled only after a file is
loaded. \gui{Isolate Burst} additionally requires background subtraction.

\problem{Rectangle zoom does nothing.}
Click \gui{Lock zooming and panning} first, then \gui{Rectangular Zooming}, and drag a
rectangle. Each activation allows one rectangle.

\problem{Drift-rate picking does not stop.}
Right-click or double-click to finish the series of points.

\problem{Selected files cannot be combined.}
The message names the reason --- for example non-consecutive timestamps, mixed stations or an
incomplete time~\(\times\)~focus-code grid. Check the rules in Table~\ref{tab:combine-rules}.

\problem{The UT time axis is unavailable.}
UT display, UT display ranges and UT alignment in the Multi-Station Comparison need the
\texttt{TIME-OBS} keyword in the FITS header. Use seconds instead.

\problem{The thresholds were reset after background subtraction.}
Limits chosen on raw data would saturate the background-subtracted scale and are therefore
reset. Set them again, or apply a preset.

\problem{The application closed unexpectedly.}
Restart it and accept the offer to recover the autosave snapshot, or use \menu{File > Recover
Last Session...}.

\problem{A save fails on Windows.}
The proposed folder may not be writable (for example inside \file{C:\\Program Files}); choose
another folder when asked.

\section{Downloads and archives}

\problem{The station list is empty or a station is missing.}
Station lists contain only stations with data on the selected UTC days. Check the date and the
network connection.

\problem{A listed burst has no FITS files.}
Catalog detections do not guarantee that matching files are archived. Increase
\gui{Extra time}, or search another station.

\problem{Recent SOHO/LASCO or SOHO/EIT data cannot be found.}
The calibrated archives lag real time by many months. Use \gui{Find Latest} to locate the newest
archived data, or \gui{Live Preview (Helioviewer)} for recent quicklook images.

\problem{JSOC downloads, binned frames or cutouts are refused.}
Enter a JSOC notify e-mail registered at
\url{https://jsoc.stanford.edu/ajax/register_email.html} in the \gui{Data Source} card. Binned
and cutout frame sizes require the JSOC source.

\problem{The disk is filling up.}
Empty the caches with the \gui{Clear Cache} buttons of the e-CALLISTO Downloader and the SunPy
Explorer, or with \menu{Data > Reset All (clear cache \& defaults)} in Solar Image Analysis.

\section{Solar Image Analysis and GCS fitting}

\problem{The PFSS card is disabled.}
When running from source, install the \texttt{sunkit-magex} package; the card shows the command.

\problem{Clicked PFSS seeds or the crop rectangle do not respond.}
These interactions need the \gui{PyQtGraph} renderer (\gui{Display \& Crop} card).

\problem{A GCS panel reports that it has no data.}
The STEREO-B archive ends on 27~September~2014. Switch that panel to another channel.

\problem{GCS refinement does not change the direction.}
The direction is constrained only by at least four points in two views separated by
20°--160°. Add points in a second, well-separated view.

\problem{MP4 export is not offered.}
MP4 export needs FFmpeg, which is bundled with the installers. If it is unavailable, save a GIF
instead.

\section{Installation}

\problem{macOS reports that the application cannot be verified.}
Open \menu{System Settings > Privacy \& Security} and click \gui{Open Anyway}, or Control-click
the application and choose \gui{Open}.

\problem{The application does not start on Linux.}
Install the package with \texttt{apt}, which also installs the required Qt/OpenGL libraries
(\texttt{libgl1}, \texttt{libegl1}, \texttt{libxkbcommon-x11-0}, \texttt{libxcb-cursor0}).

\problem{Qt reports ``DLL load failed'' when running from source on Windows.}
Rebuild the virtual environment with \file{packaging/windows/repair_windows_venv.ps1}.

\medskip
If a problem persists, report it with \menu{About > Report a Bug...}
(Section~\ref{sec:bug-report}).


% ===========================================================================
\backmatter
% ===========================================================================
%% Glossary ------------------------------------------------------------------
\chapter{Glossary}
\label{ch:glossary}
\index{glossary}

\begin{description}[style=nextline, leftmargin=1.6em, itemsep=4pt]
  \item[ADAPT] Air Force Data Assimilative Photospheric Flux Transport model; synoptic
    magnetograms that estimate the unobserved far side of the Sun.
  \item[AIA] Atmospheric Imaging Assembly on SDO; full-disk EUV and UV imager.
  \item[Alfvén Mach number (\(M_A\))] Ratio of the shock speed to the Alfvén speed.
  \item[Alfvén speed (\(V_A\))] Speed of Alfvén waves, \(B/\sqrt{\mu_0\rho}\); depends on the
    magnetic field and plasma density.
  \item[ARTEMIS-IV] Solar radio spectrograph at Thermopylae, Greece.
  \item[Background subtraction] Removal of each frequency channel's steady receiver and sky
    level from a dynamic spectrum.
  \item[Band splitting] Division of a Type~II emission lane into two parallel bands emitted
    upstream and downstream of the shock.
  \item[Base difference] Each image minus the first image of the sequence.
  \item[\texttt{BITPIX}] FITS keyword giving the data type of the stored values.
  \item[Carrington coordinates] Heliographic coordinates that rotate with the Sun at the
    Carrington rate; used for synoptic maps.
  \item[CDAW] Coordinated Data Analysis Workshop data center at NASA GSFC, host of the
    SOHO/LASCO CME catalog.
  \item[CME] Coronal mass ejection: an eruption of magnetized plasma from the corona.
  \item[Compression ratio (\(X\))] Ratio of the downstream to the upstream plasma density at a
    shock.
  \item[Coronagraph] Telescope that blocks the solar disk with an occulter to observe the faint
    corona.
  \item[Density model] Function \(n_e(r)\) giving the coronal electron density at a
    heliocentric distance.
  \item[Derotation] Compensation for solar differential rotation over an image sequence.
  \item[Digit] Raw CALLISTO intensity value (analog-to-digital converter unit, ADU), about
    0.384~dB per digit.
  \item[Drift rate] Rate of change of the emission frequency of a burst, \(\mathrm{d}f/\mathrm{d}t\).
  \item[Dst index] Disturbance storm-time index; measure of the ring current and geomagnetic
    storm strength.
  \item[Dynamic spectrum] Radio intensity as a function of frequency and time.
  \item[e-CALLISTO] World-wide network of CALLISTO solar radio spectrometers.
  \item[EIT] Extreme-ultraviolet Imaging Telescope on SOHO.
  \item[EUV] Extreme ultraviolet.
  \item[EUVI] Extreme Ultraviolet Imager of STEREO/SECCHI.
  \item[FITS] Flexible Image Transport System, the standard astronomical data format.
  \item[Focus code] Two-digit identifier of a CALLISTO receiver or frequency band at a station.
  \item[Fold number] Factor multiplying a coronal density model to represent denser structures.
  \item[Fundamental, harmonic] Plasma emission at the local plasma frequency, or at twice it.
  \item[GCS] Graduated Cylindrical Shell, a geometric flux-rope model of CMEs.
  \item[GONG] Global Oscillation Network Group, a ground-based network providing synoptic
    magnetograms.
  \item[GOES] Geostationary Operational Environmental Satellites of NOAA, carrying the XRS,
    SGPS and SUVI instruments.
  \item[HEK] Heliophysics Event Knowledgebase.
  \item[Helioprojective coordinates] Observer-centered angular coordinates on the sky, in
    arcseconds from Sun center.
  \item[Helioviewer] Service providing JPEG2000 solar images for browsing.
  \item[HI] Heliospheric Imagers (HI1, HI2) of STEREO/SECCHI.
  \item[HMI] Helioseismic and Magnetic Imager on SDO.
  \item[J-map] Time--elongation map built from a slit through a sequence of images.
  \item[JSOC] Joint Science Operations Center, the SDO data archive.
  \item[Kp index] Planetary three-hourly index of geomagnetic activity, from 0 to 9.
  \item[LASCO] Large Angle and Spectrometric Coronagraph (C2, C3) on SOHO.
  \item[Level 1, 1.5, 1b, 2] Processing levels of solar image data, from archive products to
    registered or composite images.
  \item[NRGF] Normalizing-radial-graded filter, which flattens the radial brightness gradient of
    coronagraph images.
  \item[Occulter] Disk in a coronagraph that blocks the bright inner field of view.
  \item[Open and closed field] Magnetic field lines that reach the source surface, or return to
    the photosphere.
  \item[OriginPro style] The publication graph style used for exported figures: closed frame,
    inward ticks and a white background.
  \item[PFSS] Potential-field source-surface model of the coronal magnetic field.
  \item[Plasma frequency] Natural oscillation frequency of electrons in a plasma,
    \(f_p \approx 8.98\sqrt{n_e}\)~kHz.
  \item[Position angle (PA)] Angle on the sky measured from solar north toward east.
  \item[Power law] Function of the form \(a\,t^{-b}\).
  \item[\(R_\odot\)] Solar radius, 696\,340~km.
  \item[RFI] Radio-frequency interference from terrestrial transmitters.
  \item[Ridge] The line of peak intensity of a burst in a dynamic spectrum.
  \item[Running difference] Each image minus the previous image.
  \item[SDO] Solar Dynamics Observatory.
  \item[SECCHI] Sun Earth Connection Coronal and Heliospheric Investigation on STEREO (EUVI,
    COR1, COR2, HI1, HI2).
  \item[SEP] Solar energetic particles.
  \item[SGPS] Solar and Galactic Proton Sensor on GOES-R series satellites.
  \item[SOHO] Solar and Heliospheric Observatory.
  \item[Source surface] Sphere in the PFSS model beyond which the field is radial.
  \item[STEREO] Solar Terrestrial Relations Observatory (Ahead and Behind spacecraft).
  \item[Stonyhurst coordinates] Heliographic coordinates with longitude zero facing the Earth.
  \item[SUVI] Solar Ultraviolet Imager on GOES.
  \item[SWAVES] Radio and plasma wave instrument on STEREO.
  \item[Synoptic map] Full-Sun map assembled over a solar rotation.
  \item[Type II burst] Slowly drifting radio burst produced by a shock wave.
  \item[Type III burst] Fast-drifting radio burst produced by electron beams.
  \item[UT] Universal Time.
  \item[VSO] Virtual Solar Observatory.
  \item[XRS] X-Ray Sensor on GOES, with short (0.05--0.4~nm) and long (0.1--0.8~nm) channels.
\end{description}


\printbibliography[heading=bibintoc, title=Bibliography]

\printindex

\end{document}
